\input{preamble}

% Début du chapitre.
%
\begin{document}

\title{Constructions de schémas}


\maketitle

\phantomsection
\label{section-phantom}

\tableofcontents

\section{Introduction}
\label{section-introduction}

\noindent
Dans ce chapitre, nous présentons des procédés permettant de construire des
schémas à partir d'autres schémas. Une référence fondamentale est \cite{EGA}.


\section{Recollement relatif}
\label{section-relative-glueing}

\noindent
Le lemme suivant est utile lorsque l'on cherche à construire un schéma $X$
sur $S$ et que l'on sait déjà construire sa restriction aux ouverts affines
de $S$. L'énoncé est en fait tout à fait général et vaut pour les espaces
(localement) annelés, bien que la démonstration soit rédigée dans le langage
des schémas.

\begin{lemma}
\label{lemma-relative-glueing}
Soient $S$ un schéma et $\mathcal{B}$ une base de sa topologie.
Supposons données les familles suivantes :
\begin{enumerate}
\item pour tout $U \in \mathcal{B}$, un schéma $f_U : X_U \to U$ sur $U$ ;
\item pour $U, V \in \mathcal{B}$ avec $V \subset U$, un morphisme
$\rho^U_V : X_V \to X_U$ sur $U$.
\end{enumerate}
Supposons que
\begin{enumerate}
\item[(a)] chaque $\rho^U_V$ induise un isomorphisme
$X_V \to f_U^{-1}(V)$ de schémas sur $V$ ;
\item[(b)] pour $W, V, U \in \mathcal{B}$ tels que
$W \subset V \subset U$, on ait $\rho^U_W = \rho^U_V \circ \rho ^V_W$.
\end{enumerate}
Il existe alors un morphisme de schémas $f : X \to S$ et des isomorphismes
$i_U : f^{-1}(U) \to X_U$ au-dessus de $U \in \mathcal{B}$ tels que,
pour $V, U \in \mathcal{B}$ avec $V \subset U$, la composée
$$
\xymatrix{
X_V \ar[r]^{i_V^{-1}} &
f^{-1}(V) \ar[rr]^{\text{inclusion}} & &
f^{-1}(U) \ar[r]^{i_U} &
X_U
}
$$
soit $\rho^U_V$. De plus, $X$ est unique à unique isomorphisme près sur $S$.
\end{lemma}

\begin{proof}
Nous allons appliquer Schémas, Lemme \ref{schemes-lemma-glue-functors}.
Définissons d'abord un foncteur contravariant $F$ de la catégorie des
schémas vers celle des ensembles. Pour un schéma $T$, posons
$$
F(T) =
\left\{
\begin{matrix}
(g, \{h_U\}_{U \in \mathcal{B}}),
\ g : T \to S, \ h_U : g^{-1}(U) \to X_U, \\
f_U \circ h_U = g|_{g^{-1}(U)},
\ h_U|_{g^{-1}(V)} = \rho^U_V \circ h_V
\ \forall\ V, U \in \mathcal{B}, V \subset U
\end{matrix}
\right\}.
$$
L'application de restriction $F(T) \to F(T')$ associée à un morphisme
$T' \to T$ est donnée par composition. Pour $W \in \mathcal{B}$,
considérons le sous-foncteur $F_W \subset F$ constitué des systèmes
$(g, \{h_U\})$ tels que $g(T) \subset W$.

\medskip\noindent
Montrons d'abord que $F$ est un faisceau pour la topologie de Zariski.
Soient $T$ un schéma, $T = \bigcup V_i$ un recouvrement ouvert et
$\xi_i \in F(V_i)$ des éléments tels que
$\xi_i|_{V_i \cap V_j} = \xi_j|_{V_i \cap V_j}$.
Écrivons $\xi_i = (g_i, \{h_{i, U}\})$.
Les morphismes $g_i$ se recollent en un unique morphisme global
$g : T \to S$. De plus, $g^{-1}(U) = \bigcup g_i^{-1}(U)$ ; les morphismes
$h_{i, U} : g_i^{-1}(U) \to X_U$ se recollent donc en un unique morphisme
$h_U : g^{-1}(U) \to X_U$. On vérifie aisément que le système
$(g, \{h_U\})$ appartient à $F(T)$. Le foncteur $F$ satisfait donc la
condition de faisceau pour la topologie de Zariski.

\medskip\noindent
Vérifions ensuite que chaque $F_W$, $W \in \mathcal{B}$, est représentable.
Nous affirmons que la transformation de foncteurs
$$
F_W \longrightarrow \Mor(-, X_W), \ (g, \{h_U\}) \longmapsto h_W
$$
est un isomorphisme. Soient $T$ un schéma et $\alpha : T \to X_W$ un
morphisme. Posons $g = f_W \circ \alpha$.
Pour $U \in \mathcal{B}$ avec $U \subset W$, définissons
$h_U : g^{-1}(U) \to X_U$ comme la composée
$(\rho^W_U)^{-1} \circ \alpha|_{g^{-1}(U)}$.
C'est possible puisque $\alpha(g^{-1}(U))$ est contenu dans $f_W^{-1}(U)$,
grâce à la condition (a). On a alors $f_U \circ h_U = g|_{g^{-1}(U)}$.
Si, de plus, $V \in \mathcal{B}$ et $V \subset U \subset W$, la condition
(b) donne $\rho^U_V \circ h_V = h_U|_{g^{-1}(V)}$.
Il reste à définir $h_U$ pour un élément quelconque $U \in \mathcal{B}$.
Comme $\mathcal{B}$ est une base de la topologie de $S$, on peut choisir
un recouvrement ouvert $U \cap W = \bigcup U_i$ avec $U_i \in \mathcal{B}$.
Puisque $g$ est à valeurs dans $W$, on a
$g^{-1}(U) = g^{-1}(U \cap W) = \bigcup g^{-1}(U_i)$.
Considérons les morphismes
$h_i = \rho^U_{U_i} \circ h_{U_i} : g^{-1}(U_i) \to X_U$.
La condition (b) permet aisément de vérifier
$h_i|_{g^{-1}(U_i) \cap g^{-1}(U_j)} = h_j|_{g^{-1}(U_i) \cap g^{-1}(U_j)}$.
Ces morphismes se recollent donc en le morphisme cherché
$h_U : g^{-1}(U) \to X_U$. Nous omettons la vérification élémentaire que
$(g, \{h_U\})$ appartient à $F_W(T)$ et s'envoie sur $\alpha$ par
l'application affichée ci-dessus.

\medskip\noindent
Vérifions ensuite que chaque inclusion $F_W \subset F$ est représentable
par des immersions ouvertes. Cela résulte directement des définitions.

\medskip\noindent
Enfin, vérifions que la famille $(F_W)_{W \in \mathcal{B}}$ recouvre $F$.
Cela résulte de la construction et du fait que $\mathcal{B}$ est une base
de la topologie de $S$.

\medskip\noindent
Soit $X$ un schéma représentant $F$, et soit
$(f, \{i_U\}) \in F(X)$ une « famille universelle ».
Chaque $F_W$ étant représentable par $X_W$ via la transformation de
foncteurs ci-dessus, $i_W : f^{-1}(W) \to X_W$ est un isomorphisme.
Le lemme est démontré.
\end{proof}

\begin{lemma}
\label{lemma-relative-glueing-sheaves}
Soient $S$ un schéma et $\mathcal{B}$ une base de sa topologie.
Supposons données les familles suivantes :
\begin{enumerate}
\item pour tout $U \in \mathcal{B}$, un schéma $f_U : X_U \to U$ sur $U$ ;
\item pour tout $U \in \mathcal{B}$, un faisceau quasi-cohérent
$\mathcal{F}_U$ sur $X_U$ ;
\item pour $U, V \in \mathcal{B}$ avec $V \subset U$, un morphisme
$\rho^U_V : X_V \to X_U$ ;
\item pour $U, V \in \mathcal{B}$ avec $V \subset U$, un morphisme
$\theta^U_V : (\rho^U_V)^*\mathcal{F}_U \to \mathcal{F}_V$.
\end{enumerate}
Supposons que
\begin{enumerate}
\item[(a)] chaque $\rho^U_V$ induise un isomorphisme
$X_V \to f_U^{-1}(V)$ de schémas sur $V$ ;
\item[(b)] chaque $\theta^U_V$ soit un isomorphisme ;
\item[(c)] pour $W, V, U \in \mathcal{B}$ avec $W \subset V \subset U$,
on ait $\rho^U_W = \rho^U_V \circ \rho ^V_W$ ;
\item[(d)] pour $W, V, U \in \mathcal{B}$ avec $W \subset V \subset U$,
on ait $\theta^U_W = \theta^V_W \circ (\rho^V_W)^*\theta^U_V$.
\end{enumerate}
Il existe alors un morphisme de schémas $f : X \to S$, un faisceau
quasi-cohérent $\mathcal{F}$ sur $X$ et des isomorphismes
$i_U : f^{-1}(U) \to X_U$ et
$\theta_U : i_U^*\mathcal{F}_U \to \mathcal{F}|_{f^{-1}(U)}$
au-dessus de $U \in \mathcal{B}$ tels que, pour $V, U \in \mathcal{B}$
avec $V \subset U$, la composée
$$
\xymatrix{
X_V \ar[r]^{i_V^{-1}} &
f^{-1}(V) \ar[rr]^{\text{inclusion}} & &
f^{-1}(U) \ar[r]^{i_U} &
X_U
}
$$
soit $\rho^U_V$ et que la composée
\begin{equation}
\label{equation-glue}
(\rho^U_V)^*\mathcal{F}_U
=
(i_V^{-1})^*((i_U^*\mathcal{F}_U)|_{f^{-1}(V)})
\xrightarrow{\theta_U|_{f^{-1}(V)}}
(i_V^{-1})^*(\mathcal{F}|_{f^{-1}(V)})
\xrightarrow{\theta_V^{-1}}
\mathcal{F}_V
\end{equation}
soit $\theta^U_V$. De plus, $(X, \mathcal{F})$ est unique à unique
isomorphisme près sur $S$.
\end{lemma}

\begin{proof}
Le Lemme \ref{lemma-relative-glueing} fournit le schéma $X$ sur $S$ et les
isomorphismes $i_U$. Pour $U \in \mathcal{B}$, posons
$\mathcal{F}'_U = i_U^*\mathcal{F}_U$ ; c'est un
$\mathcal{O}_{f^{-1}(U)}$-module quasi-cohérent. Les morphismes
$$
\mathcal{F}'_U|_{f^{-1}(V)} =
i_U^*\mathcal{F}_U|_{f^{-1}(V)} =
i_V^*(\rho^U_V)^*\mathcal{F}_U \xrightarrow{i_V^*\theta^U_V}
i_V^*\mathcal{F}_V = \mathcal{F}'_V
$$
définissent des isomorphismes
$(\theta')^U_V : \mathcal{F}'_U|_{f^{-1}(V)} \to \mathcal{F}'_V$
pour $V \subset U$ dans $\mathcal{B}$. La condition (d) exprime précisément
leur compatibilité pour un triplet $W \subset V \subset U$ d'éléments de
$\mathcal{B}$. On peut ainsi définir des isomorphismes
$$
\varphi_{12} :
\mathcal{F}'_{U_1}|_{f^{-1}(U_1 \cap U_2)}
\longrightarrow
\mathcal{F}'_{U_2}|_{f^{-1}(U_1 \cap U_2)}
$$
pour $U_1, U_2 \in \mathcal{B}$, en recouvrant l'intersection
$U_1 \cap U_2 = \bigcup V_j$ par des éléments $V_j$ de $\mathcal{B}$ et en
posant
$$
\varphi_{12}|_{f^{-1}(V_j)} =
\left((\theta')^{U_2}_{V_j}\right)^{-1}
\circ
(\theta')^{U_1}_{V_j}.
$$
Nous omettons la vérification que ces morphismes se recollent bien en un
morphisme $\varphi_{12}$ et satisfont la condition de cocycle d'une donnée
de recollement de faisceaux (au sens de Faisceaux,
Section \ref{sheaves-section-glueing-sheaves}). Par Faisceaux,
Lemme \ref{sheaves-lemma-glue-sheaves}, on obtient le faisceau
$\mathcal{F}$ sur $X$. Nous omettons la vérification de
(\ref{equation-glue}).
\end{proof}

\begin{remark}
\label{remark-relative-glueing-functorial}
Les constructions des Lemmes \ref{lemma-relative-glueing} et
\ref{lemma-relative-glueing-sheaves} sont fonctorielles.
Supposons en effet données deux familles
$(f_U : X_U \to U, \rho^U_V)$ et
$(g_U : Y_U \to U, \sigma^U_V)$ comme dans le
Lemme \ref{lemma-relative-glueing}. Supposons que, pour chaque
$U \in \mathcal{B}$, on dispose d'un morphisme $h_U : X_U \to Y_U$ sur
$U$, compatible avec les restrictions $\rho^U_V$ et $\sigma^U_V$.
La fonctorialité signifie que l'on obtient un morphisme de schémas
$h : X \to Y$ sur $S$ dont les restrictions sont les $h_U$, où
$f : X \to S$ est construit à partir de
$(f_U : X_U \to U, \rho^U_V)$ et $g : Y \to S$ à partir de
$(g_U : Y_U \to U, \sigma^U_V)$.

\medskip\noindent
De même, supposons données deux familles
$(f_U : X_U \to U, \mathcal{F}_U, \rho^U_V, \theta^U_V)$ et
$(g_U : Y_U \to U, \mathcal{G}_U, \sigma^U_V, \eta^U_V)$ comme dans le
Lemme \ref{lemma-relative-glueing-sheaves}.
Supposons que, pour chaque $U \in \mathcal{B}$, on dispose d'un morphisme
$h_U : X_U \to Y_U$ sur $U$, compatible avec $\rho^U_V$ et $\sigma^U_V$,
et d'un morphisme $\tau_U : h_U^*\mathcal{G}_U \to \mathcal{F}_U$
compatible avec $\theta^U_V$ et $\eta^U_V$.
La fonctorialité signifie que ces données définissent un morphisme de
schémas $h : X \to Y$ sur $S$ dont les restrictions sont les $h_U$, ainsi
qu'un morphisme $h^*\mathcal{G} \to \mathcal{F}$ dont les restrictions
sont les $\tau_U$. Ici, $(f : X \to S, \mathcal{F})$ est construit à
partir de $(f_U : X_U \to U, \mathcal{F}_U, \rho^U_V, \theta^U_V)$ et
$(g : Y \to S, \mathcal{G})$ à partir de
$(g_U : Y_U \to U, \mathcal{G}_U, \sigma^U_V, \eta^U_V)$.

\medskip\noindent
Nous omettons les vérifications, ainsi qu'une formulation convenable de
l'« équivalence de catégories » entre données de recollement relatif et
objets relatifs.
\end{remark}


\section{Spectre relatif par recollement}
\label{section-spec-via-glueing}

\begin{situation}
\label{situation-relative-spec}
Ici, $S$ est un schéma et $\mathcal{A}$ une $\mathcal{O}_S$-algèbre
quasi-cohérente : c'est un faisceau de $\mathcal{O}_S$-algèbres qui est
quasi-cohérent comme $\mathcal{O}_S$-module.
\end{situation}

\noindent
Dans cette section, nous indiquons comment construire un morphisme de schémas
$$
\underline{\Spec}_S(\mathcal{A}) \longrightarrow S
$$
en recollant les spectres $\Spec(\Gamma(U, \mathcal{A}))$, où $U$ parcourt
les ouverts affines de $S$. Montrons d'abord que les spectres des algèbres
de sections de $\mathcal{A}$ sur ces ouverts affines forment une famille
de schémas satisfaisant les conditions du Lemme \ref{lemma-relative-glueing}.

\begin{lemma}
\label{lemma-spec-inclusion}
Dans la Situation \ref{situation-relative-spec}, soient
$U \subset U' \subset S$ des ouverts affines. Posons
$A = \mathcal{A}(U)$ et $A' = \mathcal{A}(U')$.
L'homomorphisme d'anneaux $A' \to A$ induit un morphisme
$\Spec(A) \to \Spec(A')$, et le diagramme
$$
\xymatrix{
\Spec(A) \ar[r] \ar[d] &
\Spec(A') \ar[d] \\
U \ar[r] &
U'
}
$$
est cartésien.
\end{lemma}

\begin{proof}
Posons $R = \mathcal{O}_S(U)$ et $R' = \mathcal{O}_S(U')$.
L'homomorphisme $R \otimes_{R'} A' \to A$ est un isomorphisme puisque
$\mathcal{A}$ est quasi-cohérente (voir, par exemple, Schémas,
Lemme \ref{schemes-lemma-widetilde-pullback}). Le résultat découle de la
description du produit fibré de schémas affines dans Schémas,
Lemme \ref{schemes-lemma-fibre-product-affine-schemes}.
\end{proof}

\noindent
En particulier, le morphisme $\Spec(A) \to \Spec(A')$ du lemme est une
immersion ouverte.

\begin{lemma}
\label{lemma-transitive-spec}
Dans la Situation \ref{situation-relative-spec}, soient
$U \subset U' \subset U'' \subset S$ des ouverts affines. Posons
$A = \mathcal{A}(U)$, $A' = \mathcal{A}(U')$, $A'' = \mathcal{A}(U'')$.
La composée des morphismes $\Spec(A) \to \Spec(A')$ et
$\Spec(A') \to \Spec(A'')$ du Lemme \ref{lemma-spec-inclusion} est le
morphisme $\Spec(A) \to \Spec(A'')$ du Lemme \ref{lemma-spec-inclusion}.
\end{lemma}

\begin{proof}
Cela résulte de ce que $A'' \to A$ est la composée de $A'' \to A'$ et de
$A' \to A$, puisque $\mathcal{A}$ est un faisceau.
\end{proof}

\begin{lemma}
\label{lemma-glue-relative-spec}
Dans la Situation \ref{situation-relative-spec}, il existe un morphisme
de schémas
$$
\pi : \underline{\Spec}_S(\mathcal{A}) \longrightarrow S
$$
ayant les propriétés suivantes :
\begin{enumerate}
\item pour tout ouvert affine $U \subset S$, il existe un isomorphisme
$i_U : \pi^{-1}(U) \to \Spec(\mathcal{A}(U))$ sur $U$ ;
\item pour des ouverts affines $U \subset U' \subset S$, la composée
$$
\xymatrix{
\Spec(\mathcal{A}(U)) \ar[r]^{i_U^{-1}} &
\pi^{-1}(U) \ar[rr]^{\text{inclusion}} & &
\pi^{-1}(U') \ar[r]^{i_{U'}} &
\Spec(\mathcal{A}(U'))
}
$$
est l'immersion ouverte du Lemme \ref{lemma-spec-inclusion}.
\end{enumerate}
De plus, $\underline{\Spec}_S(\mathcal{A})$ est unique à unique
isomorphisme près sur $S$.
\end{lemma}

\begin{proof}
Cela résulte immédiatement des Lemmes \ref{lemma-relative-glueing},
\ref{lemma-spec-inclusion} et \ref{lemma-transitive-spec}.
L'unicité est énoncée dans la dernière phrase du
Lemme \ref{lemma-relative-glueing}.
\end{proof}











\section{Le spectre relatif comme foncteur}
\label{section-spec}

\noindent
Plaçons-nous dans la Situation \ref{situation-relative-spec}, c'est-à-dire que
$S$ est un schéma et que $\mathcal{A}$ est un faisceau quasi-cohérent de
$\mathcal{O}_S$-algèbres.

\medskip\noindent
Pour tout $f : T \to S$, l'image inverse
$f^*\mathcal{A}$ est un faisceau quasi-cohérent de $\mathcal{O}_T$-algèbres.
Nous allons considérer les paires $(f : T \to S, \varphi)$ où
$f$ est un morphisme de schémas et $\varphi : f^*\mathcal{A} \to \mathcal{O}_T$
un morphisme de $\mathcal{O}_T$-algèbres. Remarquons que cela revient à se
donner un homomorphisme de $f^{-1}\mathcal{O}_S$-algèbres
$\varphi : f^{-1}\mathcal{A} \to \mathcal{O}_T$ ; voir
Faisceaux, Lemme \ref{sheaves-lemma-adjointness-tensor-restrict}.
Cela revient également à se donner un homomorphisme de $\mathcal{O}_S$-algèbres
$\varphi : \mathcal{A} \to f_*\mathcal{O}_T$ ; voir
Faisceaux, Lemme \ref{sheaves-lemma-adjoint-push-pull-modules}.
Nous utiliserons sans autre mention ces trois façons d'envisager $\varphi$.

\medskip\noindent
Étant donnés une telle paire
$(f : T \to S, \varphi)$ et un morphisme $a : T' \to T$, on obtient
une seconde paire $(f' = f \circ a, \varphi' = a^*\varphi)$, appelée
l'image inverse de $(f, \varphi)$. On peut notamment décrire
$\varphi' = a^*\varphi$ comme la composée
$\mathcal{A} \to f_*\mathcal{O}_T \to f'_*\mathcal{O}_{T'}$,
où la seconde application est $f_*a^\sharp$ avec
$a^\sharp : \mathcal{O}_T \to a_*\mathcal{O}_{T'}$.
On a ainsi défini un foncteur
\begin{eqnarray}
\label{equation-spec}
F : \Sch^{opp} & \longrightarrow & \textit{Ens} \\
T & \longmapsto & F(T) = \{\text{couples }(f, \varphi) \text{ comme ci-dessus}\}
\nonumber
\end{eqnarray}

\begin{lemma}
\label{lemma-spec-base-change}
Dans la Situation \ref{situation-relative-spec}, soit $F$ le foncteur
associé ci-dessus à $(S, \mathcal{A})$.
Soit $g : S' \to S$ un morphisme de schémas.
Posons $\mathcal{A}' = g^*\mathcal{A}$ et soit $F'$ le
foncteur associé ci-dessus à $(S', \mathcal{A}')$.
Il existe alors un isomorphisme canonique
$$
F' \cong h_{S'} \times_{h_S} F
$$
de foncteurs.
\end{lemma}

\begin{proof}
Se donner une paire $(f' : T \to S', \varphi' : (f')^*\mathcal{A}' \to \mathcal{O}_T)$
revient à se donner une paire $(f, \varphi : f^*\mathcal{A} \to \mathcal{O}_T)$
ainsi qu'une factorisation $f = g \circ f'$. En effet, avec ces notations,
$(f')^* \mathcal{A}' = (f')^*g^*\mathcal{A} = f^*\mathcal{A}$.
D'où le lemme.
\end{proof}

\begin{lemma}
\label{lemma-spec-affine}
Dans la Situation \ref{situation-relative-spec}, soit $F$ le foncteur
associé ci-dessus à $(S, \mathcal{A})$.
Si $S$ est affine, alors $F$ est représentable par le
schéma affine $\Spec(\Gamma(S, \mathcal{A}))$.
\end{lemma}

\begin{proof}
Écrivons $S = \Spec(R)$ et $A = \Gamma(S, \mathcal{A})$.
Alors $A$ est une $R$-algèbre et $\mathcal{A} = \widetilde A$.
L'homomorphisme d'anneaux $R \to A$ donne une application canonique
$$
f_{univ} : \Spec(A)
\longrightarrow
S = \Spec(R).
$$
On a
$f_{univ}^*\mathcal{A} =  \widetilde{A \otimes_R A}$
par Schémas, Lemme \ref{schemes-lemma-widetilde-pullback}.
Il existe donc une application canonique
$$
\varphi_{univ} :
f_{univ}^*\mathcal{A} = \widetilde{A \otimes_R A}
\longrightarrow
\widetilde A = \mathcal{O}_{\Spec(A)}
$$
provenant de l'application de $A$-modules $A \otimes_R A \to A$,
$a \otimes a' \mapsto aa'$. Nous affirmons que la paire
$(f_{univ}, \varphi_{univ})$ représente $F$ dans ce cas.
Autrement dit, pour tout schéma $T$, l'application
$$
\Mor(T, \Spec(A)) \longrightarrow \{\text{couples } (f, \varphi)\},\quad
a \longmapsto (f_{univ} \circ a, a^*\varphi_{univ})
$$
est bijective.

\medskip\noindent
Construisons l'application inverse.
À toute paire $(f : T \to S, \varphi)$ est associé l'homomorphisme d'anneaux
$$
\xymatrix{
A = \Gamma(S, \mathcal{A}) \ar[r]^{f^*} &
\Gamma(T, f^*\mathcal{A}) \ar[r]^{\varphi} &
\Gamma(T, \mathcal{O}_T)
}
$$
Celui-ci induit un morphisme de schémas $T \to \Spec(A)$
par Schémas, Lemme \ref{schemes-lemma-morphism-into-affine}.

\medskip\noindent
Nous omettons de vérifier que cette application est inverse de celle
affichée ci-dessus.
\end{proof}

\begin{lemma}
\label{lemma-spec}
Dans la Situation \ref{situation-relative-spec},
le foncteur $F$ est représentable par un schéma.
\end{lemma}

\begin{proof}
Nous allons appliquer Schémas, Lemme \ref{schemes-lemma-glue-functors}.

\medskip\noindent
Vérifions d'abord que $F$ satisfait la condition de faisceau pour la
topologie de Zariski. Supposons que $T$ soit un schéma,
que $T = \bigcup_{i \in I} U_i$ soit un recouvrement ouvert
et que $(f_i, \varphi_i) \in F(U_i)$ vérifient
$(f_i, \varphi_i)|_{U_i \cap U_j} = (f_j, \varphi_j)|_{U_i \cap U_j}$.
Il en résulte que les morphismes $f_i : U_i \to S$
se recollent en un morphisme de schémas $f : T \to S$ tel que
$f|_{U_i} = f_i$ ; voir Schémas, Section \ref{schemes-section-glueing-schemes}.
Ainsi, $f_i^*\mathcal{A} = f^*\mathcal{A}|_{U_i}$ et, par hypothèse, les
morphismes $\varphi_i$ coïncident sur $U_i \cap U_j$. Ils se recollent donc,
par Faisceaux, Section \ref{sheaves-section-glueing-sheaves}, en un
morphisme de $\mathcal{O}_T$-algèbres $f^*\mathcal{A} \to \mathcal{O}_T$.
Cela prouve que $F$ satisfait la condition de faisceau pour la topologie
de Zariski.

\medskip\noindent
Soit $S = \bigcup_{i \in I} U_i$ un recouvrement par des ouverts affines.
Soit $F_i \subset F$ le sous-foncteur formé des
paires $(f : T \to S, \varphi)$ telles que
$f(T) \subset U_i$.

\medskip\noindent
Il faut montrer que chaque $F_i$ est représentable.
C'est le cas, car $F_i$ s'identifie au foncteur associé à $U_i$ muni de
la $\mathcal{O}_{U_i}$-algèbre quasi-cohérente $\mathcal{A}|_{U_i}$,
par le Lemme \ref{lemma-spec-base-change}.
Le résultat découle donc du Lemme \ref{lemma-spec-affine}.

\medskip\noindent
Montrons ensuite que $F_i \subset F$ est représentable par des immersions ouvertes.
Soit $(f : T \to S, \varphi) \in F(T)$. Posons $V_i = f^{-1}(U_i)$.
La définition de $F_i$ montre que, pour tout $a : T' \to T$,
on a $a^*(f, \varphi) \in F_i(T')$ si et seulement si $a(T') \subset V_i$.
C'est précisément ce qu'il fallait démontrer.

\medskip\noindent
Enfin, il faut montrer que la famille $(F_i)_{i \in I}$
recouvre $F$. Soit $(f : T \to S, \varphi) \in F(T)$.
Posons $V_i = f^{-1}(U_i)$. Puisque $S = \bigcup_{i \in I} U_i$
est un recouvrement ouvert de $S$, on voit que $T = \bigcup_{i \in I} V_i$
est un recouvrement ouvert de $T$. De plus, $(f, \varphi)|_{V_i} \in F_i(V_i)$.
Le lemme est ainsi démontré.
\end{proof}

\begin{lemma}
\label{lemma-glueing-gives-functor-spec}
Dans la Situation \ref{situation-relative-spec}, le schéma
$\pi : \underline{\Spec}_S(\mathcal{A}) \to S$
construit dans le Lemme \ref{lemma-glue-relative-spec} et le schéma
représentant le foncteur $F$ sont canoniquement isomorphes comme schémas
au-dessus de $S$.
\end{lemma}

\begin{proof}
Soit $X \to S$ le schéma représentant le foncteur $F$.
Considérons le faisceau de $\mathcal{O}_S$-algèbres
$\mathcal{R} = \pi_*\mathcal{O}_{\underline{\Spec}_S(\mathcal{A})}$.
Par construction de $\underline{\Spec}_S(\mathcal{A})$,
on dispose d'isomorphismes $\mathcal{A}(U) \to \mathcal{R}(U)$
pour tout ouvert affine $U \subset S$ ; cela résulte de la partie (1) du
Lemme \ref{lemma-glue-relative-spec}.
Pour des ouverts $U \subset U' \subset S$, ces isomorphismes sont
compatibles aux applications de restriction ; cela résulte de la partie (2) du
Lemme \ref{lemma-glue-relative-spec}.
Par Faisceaux, Lemme \ref{sheaves-lemma-restrict-basis-equivalence-modules},
ils proviennent donc d'un isomorphisme de $\mathcal{O}_S$-algèbres
$\varphi : \mathcal{A} \to \mathcal{R}$. On obtient ainsi un élément
$(\pi, \varphi) \in F(\underline{\Spec}_S(\mathcal{A}))$.
Puisque $X$ représente le foncteur $F$, il lui correspond un
morphisme de schémas $can : \underline{\Spec}_S(\mathcal{A}) \to X$
au-dessus de $S$.

\medskip\noindent
Soit $U \subset S$ un ouvert affine. Soit $F_U \subset F$ le
sous-foncteur correspondant aux paires $(f, \varphi)$ définies sur des
schémas $T$ tels que $f(T) \subset U$. Le changement de base
$X_U$ représente manifestement $F_U$. D'autre part, d'après le
Lemme \ref{lemma-spec-affine}, $F_U$ est représenté par
$\Spec(\mathcal{A}(U)) = \pi^{-1}(U)$. Autrement dit, $X_U \cong \pi^{-1}(U)$.
Nous omettons de vérifier que cette identification est induite par le
changement de base à $U$ du morphisme $can$.
\end{proof}

\begin{definition}
\label{definition-relative-spec}
Soit $S$ un schéma. Soit $\mathcal{A}$ un faisceau quasi-cohérent de
$\mathcal{O}_S$-algèbres. Le {\it spectre relatif de $\mathcal{A}$ au-dessus de
$S$}, ou simplement le {\it spectre de $\mathcal{A}$ au-dessus de $S$}, est le schéma
construit dans le Lemme \ref{lemma-glue-relative-spec} qui représente le
foncteur $F$ (\ref{equation-spec}) ; voir le
Lemme \ref{lemma-glueing-gives-functor-spec}.
On le note $\pi : \underline{\Spec}_S(\mathcal{A}) \to S$.
La « famille universelle » est un morphisme de $\mathcal{O}_S$-algèbres
$$
\mathcal{A}
\longrightarrow
\pi_*\mathcal{O}_{\underline{\Spec}_S(\mathcal{A})}
$$
\end{definition}

\noindent
Le lemme suivant affirme notamment que la formation du spectre relatif
commute aux changements de base.

\begin{lemma}
\label{lemma-spec-properties}
Soit $S$ un schéma. Soit $\mathcal{A}$ un faisceau quasi-cohérent
de $\mathcal{O}_S$-algèbres. Soit
$\pi : \underline{\Spec}_S(\mathcal{A}) \to S$
le spectre relatif de $\mathcal{A}$ au-dessus de $S$.
\begin{enumerate}
\item Pour tout ouvert affine $U \subset S$, l'image inverse
$\pi^{-1}(U)$ est affine.
\item Pour tout morphisme $g : S' \to S$, on a
$S' \times_S \underline{\Spec}_S(\mathcal{A}) =
\underline{\Spec}_{S'}(g^*\mathcal{A})$.
\item
L'application universelle
$$
\mathcal{A}
\longrightarrow
\pi_*\mathcal{O}_{\underline{\Spec}_S(\mathcal{A})}
$$
est un isomorphisme de $\mathcal{O}_S$-algèbres.
\end{enumerate}
\end{lemma}

\begin{proof}
La partie (1) résulte de la description du spectre relatif par recollement ;
voir le Lemme \ref{lemma-glue-relative-spec}.
La partie (2) découle immédiatement du Lemme \ref{lemma-spec-base-change}.
La partie (3) est locale sur $S$ et, lorsque $S$ est affine, elle résulte
par exemple du Lemme \ref{lemma-spec-affine}.
\end{proof}

\begin{lemma}
\label{lemma-canonical-morphism}
Soit $f : X \to S$ un morphisme quasi-compact et quasi-séparé de schémas.
Par Schémas, Lemme \ref{schemes-lemma-push-forward-quasi-coherent},
le faisceau $f_*\mathcal{O}_X$ est un faisceau quasi-cohérent de
$\mathcal{O}_S$-algèbres. Il existe un morphisme canonique
$$
can : X \longrightarrow \underline{\Spec}_S(f_*\mathcal{O}_X)
$$
de schémas au-dessus de $S$.
Pour tout ouvert affine $U \subset S$, la restriction $can|_{f^{-1}(U)}$
s'identifie au morphisme canonique
$$
f^{-1}(U) \longrightarrow \Spec(\Gamma(f^{-1}(U), \mathcal{O}_X))
$$
provenant de Schémas, Lemme \ref{schemes-lemma-morphism-into-affine}.
\end{lemma}

\begin{proof}
Par la définition de $\underline{\Spec}$ comme schéma représentant le
foncteur $F$, le morphisme provient de l'application canonique
$\varphi : f^*f_*\mathcal{O}_X \to \mathcal{O}_X$ (qui, par adjonction entre
image inverse et image directe, correspond à
$\text{id} : f_*\mathcal{O}_X \to f_*\mathcal{O}_X$).
L'assertion concernant la restriction à $f^{-1}(U)$
résulte de la description du spectre relatif au-dessus d'un ouvert affine ;
voir le Lemme \ref{lemma-spec-affine}.
\end{proof}










\section{Espace affine de dimension n}
\label{section-affine-n-space}

\noindent
Comme application du spectre relatif, définissons l'espace affine de
dimension $n$ au-dessus d'un schéma de base $S$. Pour tout entier $n \geq 0$,
on peut considérer le faisceau quasi-cohérent de $\mathcal{O}_S$-algèbres
$\mathcal{O}_S[T_1, \ldots, T_n]$. Il est quasi-cohérent car, comme faisceau
de $\mathcal{O}_S$-modules, il est simplement la somme directe de copies de
$\mathcal{O}_S$ indexées par les multi-indices.

\begin{definition}
\label{definition-affine-n-space}
Soient $S$ un schéma et $n \geq 0$.
Le schéma
$$
\mathbf{A}^n_S =
\underline{\Spec}_S(\mathcal{O}_S[T_1, \ldots, T_n])
$$
au-dessus de $S$ est appelé {\it espace affine de dimension $n$ au-dessus de $S$}.
Si $S = \Spec(R)$ est affine, on l'appelle aussi
{\it espace affine de dimension $n$ au-dessus de $R$} et on le note $\mathbf{A}^n_R$.
\end{definition}

\noindent
Remarquons que $\mathbf{A}^n_R = \Spec(R[T_1, \ldots, T_n])$.
Pour tout morphisme de schémas $g : S' \to S$, on a
$g^*\mathcal{O}_S[T_1, \ldots, T_n] = \mathcal{O}_{S'}[T_1, \ldots, T_n]$ ;
ainsi, $\mathbf{A}^n_{S'} = S' \times_S \mathbf{A}^n_S$ est le changement
de base. Une définition équivalente de l'espace affine de dimension $n$
est donc donnée par la formule
$$
\mathbf{A}^n_S = S \times_{\Spec(\mathbf{Z})} \mathbf{A}^n_{\mathbf{Z}}.
$$
En outre, un morphisme d'un $S$-schéma $f : X \to S$
vers $\mathbf{A}^n_S$ est déterminé par un homomorphisme de
$\mathcal{O}_S$-algèbres
$\mathcal{O}_S[T_1, \ldots, T_n] \to f_*\mathcal{O}_X$.
Cela revient manifestement à se donner les images des $T_i$.
Autrement dit, se donner un morphisme de $X$ vers $\mathbf{A}^n_S$ au-dessus de $S$
revient à se donner $n$ éléments
$h_1, \ldots, h_n \in \Gamma(X, \mathcal{O}_X)$.
















\section{Fibrés vectoriels}
\label{section-vector-bundle}

\noindent
Soit $S$ un schéma.
Soit $\mathcal{E}$ un faisceau quasi-cohérent de $\mathcal{O}_S$-modules.
Par Modules, Lemme \ref{modules-lemma-whole-tensor-algebra-permanence},
l'algèbre symétrique $\text{Sym}(\mathcal{E})$ de
$\mathcal{E}$ sur $\mathcal{O}_S$
est un faisceau quasi-cohérent de $\mathcal{O}_S$-algèbres.
On peut donc lui appliquer la construction du spectre relatif.

\begin{definition}
\label{definition-vector-bundle}
Soit $S$ un schéma. Soit $\mathcal{E}$ un
$\mathcal{O}_S$-module quasi-cohérent\footnote{Le lecteur pourrait s'attendre
à ce que l'on suppose ici $\mathcal{E}$ localement libre de type fini. Nous ne
le faisons pas, afin de rester cohérents avec \cite[II, Definition 1.7.8]{EGA}.}.
Le {\it fibré vectoriel associé à $\mathcal{E}$} est
$$
\mathbf{V}(\mathcal{E}) = \underline{\Spec}_S(\text{Sym}(\mathcal{E})).
$$
\end{definition}

\noindent
Le fibré vectoriel associé à $\mathcal{E}$ est muni d'une structure
supplémentaire. En effet, on a une graduation
$$
\pi_*\mathcal{O}_{\mathbf{V}(\mathcal{E})} =
\bigoplus\nolimits_{n \geq 0} \text{Sym}^n(\mathcal{E}).
$$
qui fait de $\pi_*\mathcal{O}_{\mathbf{V}(\mathcal{E})}$
une $\mathcal{O}_S$-algèbre graduée. Réciproquement, on retrouve
$\mathcal{E}$ comme sa partie de degré $1$.
Nous définissons donc un fibré vectoriel abstrait comme suit.

\begin{definition}
\label{definition-abstract-vector-bundle}
Soit $S$ un schéma. Un {\it fibré vectoriel $\pi : V \to S$ au-dessus de $S$}
est un morphisme affine de schémas tel que $\pi_*\mathcal{O}_V$ soit muni
d'une structure de $\mathcal{O}_S$-algèbre graduée
$\pi_*\mathcal{O}_V = \bigoplus\nolimits_{n \geq 0} \mathcal{E}_n$
telle que $\mathcal{E}_0 = \mathcal{O}_S$ et que les applications
$$
\text{Sym}^n(\mathcal{E}_1) \longrightarrow \mathcal{E}_n
$$
soient des isomorphismes pour tout $n \geq 0$. Un {\it morphisme de fibrés
vectoriels au-dessus de $S$} est un morphisme $f : V \to V'$ tel que
l'application induite
$$
f^* : \pi'_*\mathcal{O}_{V'} \longrightarrow \pi_*\mathcal{O}_V
$$
soit compatible avec les graduations données.
\end{definition}

\noindent
L'espace affine de dimension $n$,
$\mathbf{A}^n_S$, est un exemple de fibré vectoriel au-dessus de $S$ ;
voir la Définition \ref{definition-affine-n-space}.
En effet,
$\mathcal{O}_S[T_1, \ldots, T_n] = \text{Sym}(\mathcal{O}_S^{\oplus n})$.

\begin{lemma}
\label{lemma-category-vector-bundles}
La catégorie des fibrés vectoriels au-dessus d'un schéma $S$ est
anti-équivalente à la catégorie des $\mathcal{O}_S$-modules quasi-cohérents.
\end{lemma}

\begin{proof}
Omis. Indication : dans un sens, on utilise le foncteur
$\underline{\Spec}_S(\text{Sym}^*_{\mathcal{O}_S}(-))$,
et dans l'autre le foncteur
$(\pi : V \to S) \leadsto (\pi_*\mathcal{O}_V)_1$, où l'indice
signifie que l'on prend la partie de degré $1$.
\end{proof}




\section{Cônes}
\label{section-cone}

\noindent
En géométrie algébrique, les cônes correspondent aux algèbres graduées.
Selon nos conventions, un anneau ou une algèbre gradué $A$ est muni d'une graduation
$A = \bigoplus_{d \geq 0} A_d$ par les entiers positifs ou nuls ; voir
Algèbre, Section \ref{algebra-section-graded}.

\begin{definition}
\label{definition-cone}
Soit $S$ un schéma. Soit $\mathcal{A}$ une
$\mathcal{O}_S$-algèbre graduée quasi-cohérente. Supposons que
$\mathcal{O}_S \to \mathcal{A}_0$ soit un isomorphisme\footnote{On impose
souvent que $\mathcal{A}$ soit engendrée par $\mathcal{A}_1$ sur
$\mathcal{O}_S$. Nous ne le supposons pas, afin de rester cohérents avec
\cite[II, (8.3.1)]{EGA}.}.
Le {\it cône associé à $\mathcal{A}$}, ou le
{\it cône affine associé à $\mathcal{A}$}, est
$$
C(\mathcal{A}) = \underline{\Spec}_S(\mathcal{A}).
$$
\end{definition}

\noindent
Le cône associé à un faisceau gradué de $\mathcal{O}_S$-algèbres
est muni d'une structure supplémentaire : on obtient une graduation
$$
\pi_*\mathcal{O}_{C(\mathcal{A})} =
\bigoplus\nolimits_{n \geq 0} \mathcal{A}_n
$$
On peut donc définir un cône abstrait comme suit.

\begin{definition}
\label{definition-abstract-cone}
Soit $S$ un schéma. Un {\it cône $\pi : C \to S$ au-dessus de $S$} est un
morphisme affine de schémas tel que $\pi_*\mathcal{O}_C$ soit muni d'une
structure de $\mathcal{O}_S$-algèbre graduée
$\pi_*\mathcal{O}_C = \bigoplus\nolimits_{n \geq 0} \mathcal{A}_n$
telle que $\mathcal{A}_0 = \mathcal{O}_S$. Un {\it morphisme de cônes}
de $\pi : C \to S$ vers $\pi' : C' \to S$
est un morphisme $f : C \to C'$ tel que l'application induite
$$
f^* : \pi'_*\mathcal{O}_{C'} \longrightarrow \pi_*\mathcal{O}_C
$$
soit compatible avec les graduations données.
\end{definition}

\noindent
Tout fibré vectoriel est un exemple de cône. En fait, la catégorie des
fibrés vectoriels au-dessus de $S$ est une sous-catégorie pleine de la
catégorie des cônes au-dessus de $S$.

\section{Proj d'un anneau gradué}
\label{section-proj}

\noindent
Dans cette section, nous construisons le Proj d'un anneau gradué
en suivant \cite[II, Section 2]{EGA}.

\medskip\noindent
Soit $S$ un anneau gradué. Considérons l'espace topologique $\text{Proj}(S)$
associé à $S$ ; voir Algèbre, Section \ref{algebra-section-proj}.
Nous allons munir cet espace d'un faisceau d'anneaux $\mathcal{O}_{\text{Proj}(S)}$
de sorte que la paire $(\text{Proj}(S), \mathcal{O}_{\text{Proj}(S)})$
ainsi obtenue soit un schéma.

\medskip\noindent
Rappelons que $\text{Proj}(S)$ possède une base d'ouverts $D_{+}(f)$,
$f \in S_d$, $d \geq 1$, que nous appelons {\it ouverts standards} ; voir
Algèbre, Section \ref{algebra-section-proj}. Cette terminologie sous-entendra
toujours que $f$ est homogène de degré strictement positif, même si nous
omettons de le préciser. En outre, l'intersection de deux ouverts standards
est encore un ouvert standard :
$D_{+}(f) \cap D_{+}(g) = D_{+}(fg)$ pour $f, g \in S$ homogènes de degré
strictement positif.

\begin{lemma}
\label{lemma-standard-open}
Soit $S$ un anneau gradué. Soit $f \in S$ homogène de degré strictement positif.
\begin{enumerate}
\item Si $g\in S$ est homogène de degré strictement positif
et $D_{+}(g) \subset D_{+}(f)$, alors
\begin{enumerate}
\item $f$ est inversible dans $S_g$, et
$f^{\deg(g)}/g^{\deg(f)}$ est inversible dans $S_{(g)}$,
\item $g^e = af$ pour un certain $e \geq 1$ et un certain $a \in S$ homogène,
\item il existe un homomorphisme canonique de $S$-algèbres $S_f \to S_g$,
\item il existe un homomorphisme canonique de $S_0$-algèbres
$S_{(f)} \to S_{(g)}$, compatible avec l'application $S_f \to S_g$,
\item l'application $S_{(f)} \to S_{(g)}$ induit un isomorphisme
$$
(S_{(f)})_{g^{\deg(f)}/f^{\deg(g)}} \cong S_{(g)},
$$
\item ces applications induisent un diagramme commutatif d'espaces
topologiques
$$
\xymatrix{
D_{+}(g) \ar[d] &
\{\text{idéaux premiers }\mathbf{Z}\text{-gradués de }S_g\} \ar[l] \ar[r] \ar[d] &
\Spec(S_{(g)}) \ar[d] \\
D_{+}(f) &
\{\text{idéaux premiers }\mathbf{Z}\text{-gradués de }S_f\} \ar[l] \ar[r] &
\Spec(S_{(f)})
}
$$
où les applications horizontales sont des homéomorphismes et les applications
verticales des immersions ouvertes,
\item pour tout $S$-module gradué $M$, il existe des applications canoniques
compatibles de $S_f$-modules et de $S_{(f)}$-modules
$M_f \to M_g$ et $M_{(f)} \to M_{(g)}$, et
\item l'application $M_{(f)} \to M_{(g)}$ induit un isomorphisme
$$
(M_{(f)})_{g^{\deg(f)}/f^{\deg(g)}} \cong M_{(g)}.
$$
\end{enumerate}
\item Tout recouvrement ouvert de $D_{+}(f)$ admet un raffinement qui est un
recouvrement ouvert fini de la forme $D_{+}(f) = \bigcup_{i = 1}^n D_{+}(g_i)$.
\item Soient $g_1, \ldots, g_n \in S$ homogènes de degré strictement positif.
Alors $D_{+}(f) \subset \bigcup D_{+}(g_i)$ si et seulement si
$g_1^{\deg(f)}/f^{\deg(g_1)}, \ldots, g_n^{\deg(f)}/f^{\deg(g_n)}$
engendrent l'idéal unité de $S_{(f)}$.
\end{enumerate}
\end{lemma}

\begin{proof}
Rappelons que $D_{+}(g) = \Spec(S_{(g)})$, l'identification étant donnée par
les homomorphismes d'anneaux $S \to S_g \leftarrow S_{(g)}$ ; voir
Algèbre, Lemme \ref{algebra-lemma-topology-proj}.
Ainsi, $f^{\deg(g)}/g^{\deg(f)}$ est un élément de $S_{(g)}$ qui n'appartient
à aucun idéal premier ; il est donc inversible, par
Algèbre, Lemme \ref{algebra-lemma-Zariski-topology}.
On en déduit (a).
Écrivons l'inverse de $f$ dans $S_g$ sous la forme $a/g^d$.
On peut remplacer $a$ par sa partie homogène de degré
$d\deg(g) - \deg(f)$.
Cela signifie que $g^d - af$ est annulé par une puissance de $g$, d'où
$g^e = af$ pour un certain $a \in S$ homogène de degré
$e\deg(g) - \deg(f)$. Ceci prouve (b).
Pour (c), l'application $S_f \to S_g$ existe, d'après (a), par la propriété
universelle de la localisation ; on peut aussi la définir en envoyant $b/f^n$
sur $a^nb/g^{ne}$. Elle induit manifestement une application entre les
sous-anneaux d'éléments de degré zéro $S_{(f)} \to S_{(g)}$.
De même, on peut définir $M_f \to M_g$ et $M_{(f)} \to M_{(g)}$ en envoyant
$x/f^n$ sur $a^nx/g^{ne}$. Les assertions qui présentent $S_{(g)}$,
respectivement $M_{(g)}$, comme une localisation principale de $S_{(f)}$,
respectivement de $M_{(f)}$, résultent aussitôt des formules précédentes.
Les applications du diagramme commutatif d'espaces topologiques correspondent
aux homomorphismes d'anneaux ci-dessus. Les flèches horizontales sont des
homéomorphismes par Algèbre, Lemme \ref{algebra-lemma-topology-proj}.
Les flèches verticales sont des immersions ouvertes, puisque celle de gauche
est l'inclusion d'un sous-ensemble ouvert.

\medskip\noindent
L'ouvert $D_{+}(f)$ est quasi-compact, car il est homéomorphe à
$\Spec(S_{(f)})$ ; voir Algèbre, Lemme \ref{algebra-lemma-quasi-compact}.
La deuxième assertion résulte donc directement du fait que les ouverts
standards forment une base de la topologie.

\medskip\noindent
La troisième assertion découle directement de
Algèbre, Lemme \ref{algebra-lemma-Zariski-topology}.
\end{proof}

\noindent
Dans Faisceaux, Section \ref{sheaves-section-bases}, nous avons défini
la notion de faisceau sur une base et montré qu'elle est essentiellement
équivalente à celle de faisceau sur l'espace ; voir Faisceaux, Lemmes
\ref{sheaves-lemma-extend-off-basis} et
\ref{sheaves-lemma-extend-off-basis-structures}. En outre, nous avons montré
dans Faisceaux, Lemme
\ref{sheaves-lemma-cofinal-systems-coverings-standard-case}, qu'il suffit de
vérifier la condition de faisceau sur un système cofinal de recouvrements
ouverts de chaque ouvert standard. D'après le lemme précédent, il suffit donc
de la vérifier sur les recouvrements finis par des ouverts standards.

\begin{definition}
\label{definition-standard-covering}
Soit $S$ un anneau gradué. Supposons que
$D_{+}(f) \subset \text{Proj}(S)$ soit un ouvert standard. Un
{\it recouvrement ouvert standard} de $D_{+}(f)$ est un recouvrement
$D_{+}(f) = \bigcup_{i = 1}^n D_{+}(g_i)$, où
$g_1, \ldots, g_n \in S$ sont homogènes de degré strictement positif.
\end{definition}

\noindent
Soit $S$ un anneau gradué et soit $M$ un $S$-module gradué. Nous allons
définir un préfaisceau $\widetilde M$ sur la base des ouverts standards.
Supposons que $U \subset \text{Proj}(S)$ soit un ouvert standard.
Si $f, g \in S$ sont homogènes de degré strictement positif et vérifient
$D_{+}(f) = D_{+}(g)$, le Lemme \ref{lemma-standard-open} fournit des
applications canoniques $M_{(f)} \to M_{(g)}$ et
$M_{(g)} \to M_{(f)}$, inverses l'une de l'autre. On peut donc choisir
n'importe quel $f$ tel que $U = D_{+}(f)$ et poser
$$
\widetilde M(U) = M_{(f)}.
$$
Remarquons que, si $D_{+}(g) \subset D_{+}(f)$, le
Lemme \ref{lemma-standard-open} donne une application canonique
$$
\widetilde M(D_{+}(f)) = M_{(f)} \longrightarrow
M_{(g)} = \widetilde M(D_{+}(g)).
$$
On obtient ainsi un préfaisceau de groupes abéliens sur la base des ouverts
standards. Si $M = S$, alors $\widetilde S$ est un préfaisceau d'anneaux sur
cette base. Pour $M$ quelconque, $\widetilde M$ est un préfaisceau de
$\widetilde S$-modules sur la même base.

\medskip\noindent
Calculons la fibre de $\widetilde M$ en un point
$x \in \text{Proj}(S)$. Supposons que $x$ corresponde à l'idéal premier
homogène $\mathfrak p \subset S$. Par définition de la fibre, on a
$$
\widetilde M_x
=
\colim_{f\in S_d, d > 0, f\not\in \mathfrak p} M_{(f)}
$$
L'ensemble $\{f \in S_d, d > 0, f \not \in \mathfrak p\}$ est ici préordonné
par la règle
$f \geq f' \Leftrightarrow D_{+}(f) \subset D_{+}(f')$.
Si $f_1, f_2 \in S \setminus \mathfrak p$ sont homogènes de degré strictement
positif, alors $f_1f_2 \geq f_1$ pour cet ordre. Dans Algèbre,
Section \ref{algebra-section-proj}, nous avons défini
$M_{(\mathfrak p)}$ comme le module dont les éléments sont les fractions
$x/f$, où $x$ et $f$ sont homogènes,
$\deg(x) = \deg(f)$ et $f \not \in \mathfrak p$.
Comme $\mathfrak p \in \text{Proj}(S)$, il existe un élément
$f_0 \in S$, homogène de degré strictement positif, qui n'appartient pas à
$\mathfrak p$. Puisque $x/f = f_0x/ff_0$, on peut toujours supposer que le
dénominateur d'un élément de $M_{(\mathfrak p)}$ est de degré strictement
positif. Ces remarques donnent aussitôt
$$
\widetilde M_x = M_{(\mathfrak p)}.
$$

\medskip\noindent
Vérifions maintenant la condition de faisceau pour les recouvrements ouverts
standards. Si
$D_{+}(f) = \bigcup_{i = 1}^n D_{+}(g_i)$, cette condition équivaut à
l'exactitude de la suite
$$
0 \to M_{(f)} \to \bigoplus M_{(g_i)} \to \bigoplus M_{(g_ig_j)}.
$$
Comme $D_{+}(g_i) = D_{+}(fg_i)$, on peut récrire cette suite sous la forme
$$
0 \to M_{(f)} \to \bigoplus M_{(fg_i)} \to \bigoplus M_{(fg_ig_j)}.
$$
D'après le Lemme \ref{lemma-standard-open}, les éléments
$g_1^{\deg(f)}/f^{\deg(g_1)}, \ldots, g_n^{\deg(f)}/f^{\deg(g_n)}$
engendrent l'idéal unité de $S_{(f)}$, et les modules
$M_{(fg_i)}$ et $M_{(fg_ig_j)}$ sont les localisations principales du
$S_{(f)}$-module $M_{(f)}$ par rapport à ces éléments et à leurs produits.
On peut donc appliquer Algèbre, Lemme \ref{algebra-lemma-cover-module}, au
$S_{(f)}$-module $M_{(f)}$ et aux éléments
$g_1^{\deg(f)}/f^{\deg(g_1)}, \ldots, g_n^{\deg(f)}/f^{\deg(g_n)}$.
La suite est ainsi exacte. D'après les remarques précédentes,
$\widetilde M$ est donc un faisceau sur la base des ouverts standards.

\medskip\noindent
Il résulte alors de Faisceaux, Section \ref{sheaves-section-bases}, qu'il
existe un unique faisceau d'anneaux $\mathcal{O}_{\text{Proj}(S)}$ qui coïncide
avec $\widetilde S$ sur les ouverts standards. Le calcul des fibres effectué
ci-dessus et Algèbre, Lemme \ref{algebra-lemma-proj-prime}, montrent que toutes
les fibres de ce faisceau d'anneaux sont des anneaux locaux.

\medskip\noindent
De même, pour tout $S$-module gradué $M$, il existe un unique faisceau
$\mathcal{F}$ de $\mathcal{O}_{\text{Proj}(S)}$-modules qui coïncide avec
$\widetilde M$ sur les ouverts standards ; voir Faisceaux,
Lemme \ref{sheaves-lemma-extend-off-basis-module}.

\begin{definition}
\label{definition-structure-sheaf}
Soit $S$ un anneau gradué.
\begin{enumerate}
\item Le {\it faisceau structural $\mathcal{O}_{\text{Proj}(S)}$ du spectre
homogène de $S$} est l'unique faisceau d'anneaux
$\mathcal{O}_{\text{Proj}(S)}$ qui coïncide avec $\widetilde S$ sur la base
des ouverts standards.
\item L'espace localement annelé
$(\text{Proj}(S), \mathcal{O}_{\text{Proj}(S)})$ est appelé le
{\it spectre homogène} de $S$ et noté $\text{Proj}(S)$.
\item Le faisceau de $\mathcal{O}_{\text{Proj}(S)}$-modules qui prolonge
$\widetilde M$ à tous les ouverts de $\text{Proj}(S)$ est appelé le faisceau
de $\mathcal{O}_{\text{Proj}(S)}$-modules associé à $M$. Lui aussi est noté
$\widetilde M$.
\end{enumerate}
\end{definition}

\noindent
Résumons les résultats obtenus jusqu'ici.

\begin{lemma}
\label{lemma-proj-sheaves}
Soit $S$ un anneau gradué, soit $M$ un $S$-module gradué et soit
$\widetilde M$ le faisceau de $\mathcal{O}_{\text{Proj}(S)}$-modules associé
à $M$.
\begin{enumerate}
\item Pour tout $f \in S$ homogène de degré strictement positif, on a
$$
\Gamma(D_{+}(f), \mathcal{O}_{\text{Proj}(S)}) = S_{(f)}.
$$
\item Pour tout $f\in S$ homogène de degré strictement positif, on a
$\Gamma(D_{+}(f), \widetilde M) = M_{(f)}$ en tant que $S_{(f)}$-module.
\item Si $D_{+}(g) \subset D_{+}(f)$, les applications de restriction de
$\mathcal{O}_{\text{Proj}(S)}$ et de $\widetilde M$ sont les applications
$S_{(f)} \to S_{(g)}$ et $M_{(f)} \to M_{(g)}$ du
Lemme \ref{lemma-standard-open}.
\item Soit $\mathfrak p$ un idéal premier homogène de $S$ qui ne contient pas
$S_{+}$, et soit $x \in \text{Proj}(S)$ le point correspondant. On a
$\mathcal{O}_{\text{Proj}(S), x} = S_{(\mathfrak p)}$.
\item Soit $\mathfrak p$ un idéal premier homogène de $S$ qui ne contient pas
$S_{+}$, et soit $x \in \text{Proj}(S)$ le point correspondant. On a
$(\widetilde M)_x = M_{(\mathfrak p)}$ en tant que
$S_{(\mathfrak p)}$-module.
\item
\label{item-map}
Il existe un homomorphisme canonique d'anneaux
$
S_0 \longrightarrow \Gamma(\text{Proj}(S), \widetilde S)
$
et un morphisme canonique de $S_0$-modules
$
M_0 \longrightarrow \Gamma(\text{Proj}(S), \widetilde M)
$
compatibles avec les descriptions ci-dessus des sections sur les ouverts
standards et des fibres.
\end{enumerate}
Toutes ces identifications sont en outre fonctorielles en $M$. En particulier,
le foncteur $M \mapsto \widetilde M$ est un
foncteur exact de la catégorie des $S$-modules gradués vers la catégorie des
$\mathcal{O}_{\text{Proj}(S)}$-modules.
\end{lemma}

\begin{proof}
Les assertions (1) à (5) résultent de la discussion précédente.
L'assertion (6) découle de l'existence des applications canoniques
$M_0 \to M_{(f)}$, $x \mapsto x/1$, compatibles avec les applications de
restriction décrites en (3). L'exactitude du foncteur
$M \mapsto \widetilde M$ résulte de celle du foncteur
$M \mapsto M_{(\mathfrak p)}$ (voir Algèbre,
Lemme \ref{algebra-lemma-proj-prime}) et du fait que l'exactitude d'une suite
courte de faisceaux se vérifie sur les fibres ; voir Modules,
Lemme \ref{modules-lemma-abelian}.
\end{proof}

\begin{remark}
\label{remark-global-sections-not-isomorphism}
L'application de $M_0$ dans les sections globales de $\widetilde M$ est en
général loin d'être un isomorphisme. Un exemple élémentaire consiste à prendre
$S = k[x, y, z]$ avec $1 = \deg(x) = \deg(y) = \deg(z)$ (ou un nombre
quelconque de variables) et
$M = S/(x^{100}, y^{100}, z^{100})$. On voit facilement que
$\widetilde M = 0$, tandis que $M_0 = k$.
\end{remark}

\begin{lemma}
\label{lemma-standard-open-proj}
Soit $S$ un anneau gradué et soit $f \in S$ homogène de degré strictement
positif. Supposons que $D(g) \subset \Spec(S_{(f)})$ soit un ouvert standard.
Il existe alors un élément $h \in S$ homogène de degré strictement positif tel
que $D(g)$ corresponde à $D_{+}(h) \subset D_{+}(f)$ par l'homéomorphisme
d'Algèbre, Lemme \ref{algebra-lemma-topology-proj}. On peut même choisir $h$
de sorte que $g = h/f^n$ pour un certain $n$.
\end{lemma}

\begin{proof}
Écrivons $g = h/f^n$, où $h$ est homogène de degré strictement positif et
$n \geq 1$. Si $D_{+}(h)$ n'est pas contenu dans $D_{+}(f)$, remplaçons $h$
par $hf$ et $n$ par $n + 1$. L'élément $h$ a alors la forme voulue et
$D_{+}(h) \subset D_{+}(f)$ correspond à
$D(g) \subset \Spec(S_{(f)})$.
\end{proof}

\begin{lemma}
\label{lemma-proj-scheme}
Soit $S$ un anneau gradué. L'espace localement annelé $\text{Proj}(S)$ est un
schéma. Les ouverts standards $D_{+}(f)$ sont des ouverts affines. Pour tout
$S$-module gradué $M$, le faisceau $\widetilde M$ est un faisceau
quasi-cohérent de $\mathcal{O}_{\text{Proj}(S)}$-modules.
\end{lemma}

\begin{proof}
Considérons un ouvert standard $D_{+}(f) \subset \text{Proj}(S)$.
D'après les Lemmes \ref{lemma-standard-open} et \ref{lemma-proj-sheaves},
on a $\Gamma(D_{+}(f), \mathcal{O}_{\text{Proj}(S)}) = S_{(f)}$, et il existe
un homéomorphisme $\varphi : D_{+}(f) \to \Spec(S_{(f)})$.
Pour tout ouvert standard $D(g) \subset \Spec(S_{(f)})$, choisissons
$h \in S_{+}$ comme dans le Lemme \ref{lemma-standard-open-proj}.
Alors $\varphi^{-1}(D(g)) = D_{+}(h)$ et, d'après les Lemmes
\ref{lemma-proj-sheaves} et \ref{lemma-standard-open},
$$
\Gamma(D_{+}(h), \mathcal{O}_{\text{Proj}(S)})
=
S_{(h)}
=
(S_{(f)})_{h^{\deg(f)}/f^{\deg(h)}}
=
(S_{(f)})_g
=
\Gamma(D(g), \mathcal{O}_{\Spec(S_{(f)})}).
$$
La restriction de $\mathcal{O}_{\text{Proj}(S)}$ à $D_{+}(f)$ correspond donc,
par l'homéomorphisme $\varphi$, exactement au faisceau
$\mathcal{O}_{\Spec(S_{(f)})}$ défini dans Schémas,
Section \ref{schemes-section-affine-schemes}. Il s'ensuit que $D_{+}(f)$ est
un schéma affine isomorphe à $\Spec(S_{(f)})$ par $\varphi$, et donc que
$\text{Proj}(S)$ est un schéma.

\medskip\noindent
On montre exactement de la même manière que $\widetilde M$ est un faisceau
quasi-cohérent de $\mathcal{O}_{\text{Proj}(S)}$-modules. En effet, l'argument
précédent donne
$$
\widetilde M|_{D_{+}(f)} \cong \varphi^*\left(\widetilde{M_{(f)}}\right),
$$
ce qui montre que $\widetilde M$ est quasi-cohérent.
\end{proof}

\begin{lemma}
\label{lemma-proj-separated}
Soit $S$ un anneau gradué. Le schéma $\text{Proj}(S)$ est séparé.
\end{lemma}

\begin{proof}
Il faut montrer que le morphisme canonique
$\text{Proj}(S) \to \Spec(\mathbf{Z})$ est séparé. Utilisons Schémas,
Lemme \ref{schemes-lemma-characterize-separated}. Il suffit de montrer que,
pour toute paire d'ouverts standards $D_{+}(f)$ et $D_{+}(g)$,
l'intersection $D_{+}(f) \cap D_{+}(g) = D_{+}(fg)$ est affine (ce qui est
clair) et que l'homomorphisme d'anneaux
$$
S_{(f)} \otimes_{\mathbf{Z}} S_{(g)} \longrightarrow S_{(fg)}
$$
est surjectif. Tout élément $s$ de $S_{(fg)}$ est de la forme
$s = h/(f^ng^m)$, où $h \in S$ est homogène de degré
$n\deg(f) + m\deg(g)$. En multipliant $h$ par un monôme convenable
$f^ig^j$, on peut supposer que $n = n' \deg(g)$ et
$m = m' \deg(f)$. On peut alors récrire $s$ sous la forme
$s = h/f^{(n' + m')\deg(g)} \cdot
f^{m'\deg(g)}/g^{m'\deg(f)}$. Ainsi, $s$ appartient bien à l'image de la
flèche affichée.
\end{proof}

\begin{lemma}
\label{lemma-proj-quasi-compact}
Soit $S$ un anneau gradué. Le schéma $\text{Proj}(S)$ est quasi-compact si et
seulement s'il existe un nombre fini d'éléments homogènes
$f_1, \ldots, f_n \in S_{+}$ tels que
$S_{+} \subset \sqrt{(f_1, \ldots, f_n)}$. Dans ce cas,
$\text{Proj}(S) = D_+(f_1) \cup \ldots \cup D_+(f_n)$.
\end{lemma}

\begin{proof}
Si une telle famille d'éléments existe, les ouverts affines standards
$D_{+}(f_i)$ recouvrent $\text{Proj}(S)$ d'après Algèbre,
Lemme \ref{algebra-lemma-topology-proj}. Réciproquement, si
$\text{Proj}(S)$ est quasi-compact, on peut le recouvrir par un nombre fini
d'ouverts standards $D_{+}(f_i)$, $i = 1, \ldots, n$, et le lemme que nous
venons de citer donne
$S_{+} \subset \sqrt{(f_1, \ldots, f_n)}$.
\end{proof}

\begin{lemma}
\label{lemma-structure-morphism-proj}
Soit $S$ un anneau gradué. Le schéma $\text{Proj}(S)$ possède un morphisme
canonique vers le schéma affine $\Spec(S_0)$, qui coïncide sur les espaces
topologiques avec l'application d'Algèbre,
Définition \ref{algebra-definition-proj}.
\end{lemma}

\begin{proof}
Nous avons vu plus haut que notre construction de $\widetilde S$, respectivement
de $\widetilde M$, donne un faisceau de $S_0$-algèbres, respectivement de
$S_0$-modules. On obtient donc un morphisme par Schémas,
Lemme \ref{schemes-lemma-morphism-into-affine}. Restreint à $D_{+}(f)$, ce
morphisme provient de l'homomorphisme canonique d'anneaux
$S_0 \to S_{(f)}$. Les applications $S \to S_f$ et $S_{(f)} \to S_f$ sont
des homomorphismes de $S_0$-algèbres ; voir le
Lemme \ref{lemma-standard-open}. Par conséquent, si l'idéal premier homogène
$\mathfrak p \subset S$ correspond à l'idéal premier $\mathbf{Z}$-gradué
$\mathfrak p' \subset S_f$ et à l'idéal premier ordinaire
$\mathfrak p'' \subset S_{(f)}$, ces trois idéaux ont la même image inverse
dans $S_0$.
\end{proof}

\begin{lemma}
\label{lemma-proj-valuative-criterion}
Soit $S$ un anneau gradué. Si $S$ est une algèbre de type fini sur $S_0$,
le morphisme $\text{Proj}(S) \to \Spec(S_0)$ satisfait les parties existence
et unicité du critère valuatif ; voir Schémas,
Définition \ref{schemes-definition-valuative-criterion}.
\end{lemma}

\begin{proof}
L'unicité résulte de ce que $\text{Proj}(S)$ est séparé
(Lemme \ref{lemma-proj-separated} et Schémas,
Lemme \ref{schemes-lemma-separated-implies-valuative}).
Choisissons des éléments homogènes $x_i \in S_{+}$, $i = 1, \ldots, n$,
qui engendrent $S$ sur $S_0$. Posons $d_i = \deg(x_i)$ et
$d = \text{lcm}\{d_i\}$. Supposons donné un diagramme
$$
\xymatrix{
\Spec(K) \ar[r] \ar[d] & \text{Proj}(S) \ar[d] \\
\Spec(A) \ar[r] & \Spec(S_0)
}
$$
comme dans Schémas, Définition \ref{schemes-definition-valuative-criterion}.
Notons $v : K^* \to \Gamma$ la valuation de $A$ ; voir Algèbre,
Définition \ref{algebra-definition-value-group}. On peut choisir un élément
homogène $f \in S_{+}$ tel que $\Spec(K)$ s'envoie dans $D_{+}(f)$. On obtient
alors un diagramme commutatif d'homomorphismes d'anneaux
$$
\xymatrix{
K & S_{(f)} \ar[l]^{\varphi} \\
A \ar[u] & S_0 \ar[l] \ar[u]
}
$$
Après renumérotation, on peut supposer que
$\varphi(x_i^{\deg(f)}/f^{d_i})$ est non nul pour $i = 1, \ldots, r$ et nul
pour $i = r + 1, \ldots, n$. Comme les ouverts $D_{+}(x_i)$ recouvrent
$\text{Proj}(S)$, on a $r \geq 1$. Choisissons
$i_0 \in \{1, \ldots, r\}$ de sorte que
$\gamma_i = (d/d_i)v(\varphi(x_i^{\deg(f)}/f^{d_i}))$ soit minimal dans
$\Gamma$. Pour alléger les notations, posons $x_0 = x_{i_0}$ et
$d_0 = d_{i_0}$. L'homomorphisme $\varphi$ se factorise par un homomorphisme
$\varphi' : S_{(fx_0)} \to K$, qui fournit un homomorphisme
$S_{(x_0)} \to S_{(fx_0)} \to K$.
La $S_0$-algèbre $S_{(x_0)}$ est engendrée par les éléments
$x_1^{e_1} \ldots x_n^{e_n}/x_0^{e_0}$ tels que
$\sum e_i d_i = e_0 d_0$. Si $e_i > 0$ pour un certain $i > r$, alors
$\varphi'(x_1^{e_1} \ldots x_n^{e_n}/x_0^{e_0}) = 0$.
Si $e_i = 0$ pour $i > r$, on a
\begin{align*}
d \deg(f) v(\varphi'(x_1^{e_1} \ldots x_r^{e_r}/x_0^{e_0}))
& =
d v(\varphi'(x_1^{e_1 \deg(f)} \ldots x_r^{e_r \deg(f)}/x_0^{e_0 \deg(f)})) \\
& =
d \sum e_i v(\varphi'(x_i^{\deg(f)}/f^{d_i}))
- e_0 v(\varphi'(x_0^{\deg(f)}/f^{d_0})) \\
& =
\sum e_i d_i \gamma_i - e_0 d_0 \gamma_0 \\
& \geq
\sum e_i d_i \gamma_0 - e_0 d_0 \gamma_0 = 0
\end{align*}
car $\gamma_0$ est minimal parmi les $\gamma_i$. Il en résulte que
$\varphi'$ envoie $S_{(x_0)}$ dans $A$. Le morphisme de schémas correspondant
$\Spec(A) \to \Spec(S_{(x_0)}) = D_{+}(x_0)
\subset \text{Proj}(S)$ est le morphisme qui complète le premier diagramme
commutatif de cette démonstration.
\end{proof}

\noindent
Nous avons aussi vu dans la démonstration du
Lemme \ref{lemma-proj-valuative-criterion} que, sous les hypothèses de ce
lemme, le morphisme $\text{Proj}(S) \to \Spec(S_0)$ est quasi-compact.
Il résulte donc de Schémas,
Proposition \ref{schemes-proposition-characterize-universally-closed}, que
$\text{Proj}(S) \to \Spec(S_0)$ est universellement fermé dans cette
situation. Les exemples suivants montrent que ces résultats ne subsistent
pas sans hypothèse sur l'anneau gradué $S$.

\begin{example}
\label{example-not-existence-valuative-big-proj}
Soit $\mathbf{C}[X_1, X_2, X_3, \ldots]$ la $\mathbf{C}$-algèbre graduée dans
laquelle chaque $X_i$ est de degré $1$. Considérons l'homomorphisme d'anneaux
$$
\mathbf{C}[X_1, X_2, X_3, \ldots]
\longrightarrow
\mathbf{C}[t^\alpha ; \alpha \in \mathbf{Q}_{\geq 0}]
$$
qui envoie $X_i$ sur $t^{1/i}$. Le membre de droite devient un anneau de
valuation $A$ après localisation en l'idéal
$\mathfrak m = (t^\alpha ; \alpha > 0)$. Soit $K$ le corps des fractions de
$A$. L'homomorphisme ci-dessus donne un morphisme
$\Spec(K) \to \text{Proj}(\mathbf{C}[X_1, X_2, X_3, \ldots])$ qui ne se
prolonge pas en un morphisme défini sur tout $\Spec(A)$. En effet, l'image de
$\Spec(A)$ serait contenue dans l'un des $D_{+}(X_i)$ ; mais
$X_{i + 1}/X_i$ devrait alors s'envoyer sur un élément de $A$, ce qui est
impossible puisqu'il s'envoie sur $t^{1/(i + 1) - 1/i}$.
\end{example}

\begin{example}
\label{example-not-existence-valuative-small-proj}
Soit $R = \mathbf{C}[t]$ et
$$
S = R[X_1, X_2, X_3, \ldots]/(X_i^2 - tX_{i + 1}).
$$
La graduation est définie par $R = S_0$ et
$\deg(X_i) = 2^{i - 1}$. Remarquons que, si
$\mathfrak p \in \text{Proj}(S)$, alors $t \not \in \mathfrak p$ ; sinon,
$\mathfrak p$ contiendrait tous les $X_i$, ce qui est impossible pour un point
du spectre homogène. On a donc
$D_{+}(X_i) = D_{+}(X_{i + 1})$ pour tout $i$. Ainsi,
$\text{Proj}(S)$ est quasi-compact ; il est même affine, puisqu'il est égal à
$D_{+}(X_1)$. On voit facilement que l'image de
$\text{Proj}(S) \to \Spec(R)$ est $D(t)$. Ce morphisme n'est donc pas fermé.
Le critère valuatif ne peut par conséquent pas s'appliquer, puisqu'il
entraînerait que le morphisme est fermé (voir Schémas,
Proposition \ref{schemes-proposition-characterize-universally-closed}).
\end{example}

\begin{example}
\label{example-trivial-proj}
Soit $A$ un anneau. Munissons $S = A[T]$ de sa structure de $A$-algèbre
graduée pour laquelle $T$ est de degré $1$. Alors le morphisme canonique
$\text{Proj}(S) \to \Spec(A)$ (voir le
Lemme \ref{lemma-structure-morphism-proj}) est un isomorphisme.
\end{example}

\begin{example}
\label{example-open-subset-proj}
Soit $X = \Spec(A)$ un schéma affine et soit $U \subset X$ un sous-schéma
ouvert. Munissons $A[T]$ de la graduation définie par $\deg T = 1$.
Soit $S$ le sous-anneau de $A[T]$ engendré par $A$ et par tous les $fT^i$,
où $i \ge 0$ et où $f \in A$ vérifie $D(f) \subset U$. Nous affirmons que
$S$ est un anneau gradué tel que $S_0 = A$ et
$\text{Proj}(S) \cong U$, et que cet isomorphisme identifie le morphisme
canonique $\text{Proj}(S) \to \Spec(A)$ du
Lemme \ref{lemma-structure-morphism-proj} à l'inclusion $U \subset X$.

\medskip\noindent
Supposons que $\mathfrak p \in \text{Proj}(S)$ contienne tout élément
$fT \in S_1$. Alors il contient tout générateur $fT^i$ avec $i \ge 1$, car
$(fT^i)^2 = (fT)(fT^{2i-1}) \in \mathfrak p$ et $\mathfrak p$ est radical.
On aurait donc $\mathfrak p \supset S_+$, ce qui est impossible. Par
conséquent, les sous-schémas affines ouverts standards
$\{D_+(fT)\}_{fT \in S_1}$ recouvrent $\text{Proj}(S)$.

\medskip\noindent
Remarquons que, si $fT \in S_1$, l'inclusion $S \subset A[T]$ induit un
isomorphisme gradué de $S[(fT)^{-1}]$ sur
$A[T, T^{-1}, f^{-1}]$. L'ouvert standard
$D_+(fT) \cong \Spec(S_{(fT)})$ est donc isomorphe à
$\Spec(A[T, T^{-1}, f^{-1}]_0) = \Spec(A[f^{-1}])$.
Il est clair que cet isomorphisme est la restriction du morphisme canonique
$\text{Proj}(S) \to \Spec(A)$. Si, de plus, $gT \in S_1$, alors
$S[(fT)^{-1}, (gT)^{-1}] \cong A[T, T^{-1}, f^{-1}, g^{-1}]$
comme anneaux gradués, de sorte que
$D_+(fT) \cap D_+(gT) \cong \Spec(A[f^{-1}, g^{-1}])$.
Ainsi, $\text{Proj}(S)$ est la réunion des sous-schémas ouverts $D_+(fT)$,
qui sont isomorphes par le morphisme canonique aux sous-schémas ouverts
$D(f) \subset X$, et leurs intersections dans $\text{Proj}(S)$ correspondent
à leurs intersections dans $X$. Le morphisme canonique est donc un
isomorphisme de $\text{Proj}(S)$ sur la réunion de tous les
$D(f) \subset U$, c'est-à-dire sur $U$.
\end{example}


\section{Faisceaux quasi-cohérents sur Proj}
\label{section-quasi-coherent-proj}

\noindent
Soit $S$ un anneau gradué et soit $M$ un $S$-module gradué. Nous avons vu au
Lemme \ref{lemma-proj-sheaves} comment construire sur $\text{Proj}(S)$ un
faisceau quasi-cohérent de modules $\widetilde{M}$ ainsi qu'une application
\begin{equation}
\label{equation-map-global-sections}
M_0 \longrightarrow \Gamma(\text{Proj}(S), \widetilde{M})
\end{equation}
de la partie de degré $0$ de $M$ vers les sections globales de
$\widetilde{M}$. La partie de degré $0$ du $n$-ième décalé $M(n)$ du module
gradué $M$ (voir Algèbre, Section \ref{algebra-section-graded}) est égale à
$M_n$. On obtient donc des applications
\begin{equation}
\label{equation-map-global-sections-degree-n}
M_n \longrightarrow \Gamma(\text{Proj}(S), \widetilde{M(n)}).
\end{equation}
Nous voudrions pouvoir effectuer cette opération pour tout faisceau
quasi-cohérent $\mathcal{F}$ sur $\text{Proj}(S)$. Pour ce faire, nous
prendrons son produit tensoriel avec le $n$-ième tordu du faisceau structural ;
voir la Définition \ref{definition-twist}. Le lemme suivant relie les deux
notions.

\begin{lemma}
\label{lemma-widetilde-tensor}
Soit $S$ un anneau gradué. Soit
$(X, \mathcal{O}_X) = (\text{Proj}(S), \mathcal{O}_{\text{Proj}(S)})$ le
schéma du Lemme \ref{lemma-proj-scheme}. Soit $f \in S_{+}$ homogène.
Soit $x \in X$ un point correspondant à l'idéal premier homogène
$\mathfrak p \subset S$. Soient $M$ et $N$ des $S$-modules gradués.
Il existe un morphisme canonique de $\mathcal{O}_{\text{Proj}(S)}$-modules
$$
\widetilde M \otimes_{\mathcal{O}_X} \widetilde N
\longrightarrow
\widetilde{M \otimes_S N}
$$
qui induit sur les sections au-dessus de $D_{+}(f)$ l'application canonique
$
M_{(f)} \otimes_{S_{(f)}} N_{(f)}
\to
(M \otimes_S N)_{(f)}
$
et sur les fibres en $x$ l'application canonique
$
M_{(\mathfrak p)} \otimes_{S_{(\mathfrak p)}} N_{(\mathfrak p)}
\to
(M \otimes_S N)_{(\mathfrak p)}.
$
En outre, le diagramme suivant est commutatif :
$$
\xymatrix{
M_0 \otimes_{S_0} N_0 \ar[r] \ar[d] &
(M \otimes_S N)_0 \ar[d] \\
\Gamma(X, \widetilde M \otimes_{\mathcal{O}_X} \widetilde N) \ar[r] &
\Gamma(X, \widetilde{M \otimes_S N})
}
$$
où les applications verticales sont données par
(\ref{equation-map-global-sections}).
\end{lemma}

\begin{proof}
Construire le morphisme affiché revient à construire une application
$\mathcal{O}_X$-bilinéaire
$$
\widetilde M \times \widetilde N
\longrightarrow
\widetilde{M \otimes_S N} ;
$$
voir Modules, Section \ref{modules-section-tensor-product}. Il suffit de la
définir sur les sections au-dessus des ouverts $D_{+}(f)$, de manière
compatible avec les applications de restriction. Sur $D_{+}(f)$, on utilise
l'application $S_{(f)}$-bilinéaire
$M_{(f)} \times N_{(f)} \to (M \otimes_S N)_{(f)}$,
$(x/f^n, y/f^m) \mapsto (x \otimes y)/f^{n + m}$. Nous omettons les détails.
\end{proof}

\begin{remark}
\label{remark-not-isomorphism}
En général, le morphisme construit au Lemme \ref{lemma-widetilde-tensor}
n'est pas un isomorphisme. Voici un exemple. Soit $k$ un corps. Prenons
$S = k[x, y, z]$, où $k$ est placé en degré $0$ et
$\deg(x) = 1$, $\deg(y) = 2$, $\deg(z) = 3$.
Posons $M = S(1)$ et $N = S(2)$ ; pour les notations, voir Algèbre,
Section \ref{algebra-section-graded}. Alors $M \otimes_S N = S(3)$.
Remarquons que
\begin{eqnarray*}
S_z
& = &
k[x, y, z, 1/z] \\
S_{(z)}
& = &
k[x^3/z, xy/z, y^3/z^2]
\cong
k[u, v, w]/(uw - v^3) \\
M_{(z)} & = & S_{(z)} \cdot x + S_{(z)} \cdot y^2/z \subset S_z \\
N_{(z)} & = & S_{(z)} \cdot y + S_{(z)} \cdot x^2 \subset S_z \\
S(3)_{(z)} & = & S_{(z)} \cdot z \subset S_z
\end{eqnarray*}
Considérons l'idéal maximal $\mathfrak m = (u, v, w) \subset S_{(z)}$.
On voit sans difficulté que $M_{(z)}/\mathfrak mM_{(z)}$ et
$N_{(z)}/\mathfrak mN_{(z)}$ sont tous deux de dimension $2$ sur
$\kappa(\mathfrak m)$. En revanche,
$S(3)_{(z)}/\mathfrak mS(3)_{(z)}$ est de dimension $1$.
L'application $M_{(z)} \otimes N_{(z)} \to S(3)_{(z)}$ n'est donc pas un
isomorphisme.
\end{remark}


\section{Faisceaux inversibles sur Proj}
\label{section-invertible-on-proj}

\noindent
Rappelons la construction, donnée dans Algèbre,
Section \ref{algebra-section-graded}, du module tordu $M(n)$ associé à un
module gradué sur un anneau gradué.

\begin{definition}
\label{definition-twist}
Soit $S$ un anneau gradué et soit $X = \text{Proj}(S)$.
\begin{enumerate}
\item Nous posons $\mathcal{O}_X(n) = \widetilde{S(n)}$. Ce faisceau est
appelé le $n$-ième {\it tordu du faisceau structural de $\text{Proj}(S)$}.
\item Pour tout faisceau de $\mathcal{O}_X$-modules $\mathcal{F}$, nous posons
$\mathcal{F}(n) = \mathcal{F} \otimes_{\mathcal{O}_X} \mathcal{O}_X(n)$.
\end{enumerate}
\end{definition}

\noindent
Nous allons utiliser le Lemme \ref{lemma-widetilde-tensor} pour construire
plusieurs morphismes canoniques. Comme
$S(n) \otimes_S S(m) = S(n + m)$, il existe des morphismes canoniques
\begin{equation}
\label{equation-multiply}
\mathcal{O}_X(n) \otimes_{\mathcal{O}_X} \mathcal{O}_X(m)
\longrightarrow
\mathcal{O}_X(n + m).
\end{equation}
Ces morphismes ne sont pas des isomorphismes en général ; voir l'exemple de la
Remarque \ref{remark-not-isomorphism}. Le même exemple montre que
$\mathcal{O}_X(n)$ n'est en général {\it pas} un faisceau inversible sur $X$.
En prenant le produit tensoriel avec un $\mathcal{O}_X$-module quelconque
$\mathcal{F}$, on obtient des morphismes
\begin{equation}
\label{equation-multiply-on-sheaf}
\mathcal{O}_X(n) \otimes_{\mathcal{O}_X} \mathcal{F}(m)
\longrightarrow
\mathcal{F}(n + m).
\end{equation}
Les applications (\ref{equation-map-global-sections-degree-n}) sur les
sections globales donnent un homomorphisme d'anneaux gradués
\begin{equation}
\label{equation-global-sections}
S \longrightarrow \bigoplus\nolimits_{n \geq 0} \Gamma(X, \mathcal{O}_X(n)).
\end{equation}
Pour un $\mathcal{O}_X$-module quelconque $\mathcal{F}$, les morphismes
(\ref{equation-multiply-on-sheaf}) donnent une structure de module gradué
\begin{equation}
\label{equation-global-sections-module}
\bigoplus\nolimits_{n \geq 0} \Gamma(X, \mathcal{O}_X(n))
\times
\bigoplus\nolimits_{m \in \mathbf{Z}} \Gamma(X, \mathcal{F}(m))
\longrightarrow
\bigoplus\nolimits_{m \in \mathbf{Z}} \Gamma(X, \mathcal{F}(m))
\end{equation}
et, par (\ref{equation-global-sections}), une structure de $S$-module.
Plus généralement, si $M$ est un $S$-module gradué, on a
$M(n) = M \otimes_S S(n)$. On obtient donc des morphismes
\begin{equation}
\label{equation-multiply-more-generally}
\widetilde M(n)
=
\widetilde M
\otimes_{\mathcal{O}_X}
\mathcal{O}_X(n)
\longrightarrow
\widetilde{M(n)}.
\end{equation}
Sur les sections globales, (\ref{equation-map-global-sections-degree-n})
définit un morphisme de $S$-modules gradués
\begin{equation}
\label{equation-global-sections-more-generally}
M \longrightarrow
\bigoplus\nolimits_{n \in \mathbf{Z}} \Gamma(X, \widetilde{M(n)}).
\end{equation}
Voici un fait important qui résulte presque immédiatement des définitions.

\begin{lemma}
\label{lemma-when-invertible}
Soit $S$ un anneau gradué et posons $X = \text{Proj}(S)$. Soit
$f \in S$ homogène de degré $d > 0$. Les faisceaux
$\mathcal{O}_X(nd)|_{D_{+}(f)}$ sont inversibles, et même triviaux, pour tout
$n \in \mathbf{Z}$ (voir Modules,
Définition \ref{modules-definition-invertible}). Les morphismes
(\ref{equation-multiply}) restreints à $D_{+}(f)$
$$
\mathcal{O}_X(nd)|_{D_{+}(f)} \otimes_{\mathcal{O}_{D_{+}(f)}}
\mathcal{O}_X(m)|_{D_{+}(f)}
\longrightarrow
\mathcal{O}_X(nd + m)|_{D_{+}(f)},
$$
les morphismes (\ref{equation-multiply-on-sheaf}) restreints à $D_+(f)$
$$
\mathcal{O}_X(nd)|_{D_{+}(f)} \otimes_{\mathcal{O}_{D_{+}(f)}}
\mathcal{F}(m)|_{D_{+}(f)}
\longrightarrow
\mathcal{F}(nd + m)|_{D_{+}(f)},
$$
et les morphismes (\ref{equation-multiply-more-generally}) restreints à
$D_{+}(f)$
$$
\widetilde M(nd)|_{D_{+}(f)}
=
\widetilde M|_{D_{+}(f)}
\otimes_{\mathcal{O}_{D_{+}(f)}}
\mathcal{O}_X(nd)|_{D_{+}(f)}
\longrightarrow
\widetilde{M(nd)}|_{D_{+}(f)}
$$
sont des isomorphismes pour tous $n, m \in \mathbf{Z}$.
\end{lemma}

\begin{proof}
Les homomorphismes de $S$-modules gradués $S \to S(nd)$ et
$M \to M(nd)$ définis par $x \mapsto f^n x$ deviennent des isomorphismes
après inversion de $f$. Le premier donne
$S_{(f)} \cong S(nd)_{(f)}$, d'où un isomorphisme
$\mathcal{O}_{D_{+}(f)} \cong \mathcal{O}_X(nd)|_{D_{+}(f)}$.
Le second montre que l'application
$S(nd)_{(f)} \otimes_{S_{(f)}} M_{(f)} \to M(nd)_{(f)}$ est un isomorphisme.
Le cas du morphisme (\ref{equation-multiply-on-sheaf}) résulte du cas du
morphisme (\ref{equation-multiply}).
\end{proof}

\begin{lemma}
\label{lemma-apply-modules}
Soit $S$ un anneau gradué, soit $M$ un $S$-module gradué et posons
$X = \text{Proj}(S)$. Supposons que $X$ soit recouvert par les ouverts
standards $D_+(f)$ avec $f \in S_1$, ce qui est par exemple le cas si $S$ est
engendré par $S_1$ sur $S_0$. Alors les faisceaux $\mathcal{O}_X(n)$ sont
inversibles et les morphismes (\ref{equation-multiply}),
(\ref{equation-multiply-on-sheaf}) et
(\ref{equation-multiply-more-generally}) sont des isomorphismes.
En particulier, ces morphismes induisent des isomorphismes
$$
\mathcal{O}_X(1)^{\otimes n} \cong
\mathcal{O}_X(n)
\quad
\text{et}
\quad
\widetilde{M} \otimes_{\mathcal{O}_X} \mathcal{O}_X(n) =
\widetilde{M}(n) \cong \widetilde{M(n)}
$$
Ainsi, (\ref{equation-map-global-sections-degree-n}) devient une application
\begin{equation}
\label{equation-map-global-sections-degree-n-simplified}
M_n \longrightarrow \Gamma(X, \widetilde{M}(n))
\end{equation}
et (\ref{equation-global-sections-more-generally}) devient une application
\begin{equation}
\label{equation-global-sections-more-generally-simplified}
M \longrightarrow
\bigoplus\nolimits_{n \in \mathbf{Z}} \Gamma(X, \widetilde{M}(n)).
\end{equation}
\end{lemma}

\begin{proof}
Sous les hypothèses du lemme, les ouverts $D_{+}(f)$ avec $f \in S_1$
recouvrent $X$ ; le résultat découle donc du
Lemme \ref{lemma-when-invertible}.
\end{proof}

\begin{lemma}
\label{lemma-where-invertible}
Soit $S$ un anneau gradué, posons $X = \text{Proj}(S)$ et fixons un entier
$d \geq 1$. Les deux ouverts suivants de $X$ sont égaux :
\begin{enumerate}
\item le plus grand ouvert $W = W_d \subset X$ tel que chaque
$\mathcal{O}_X(dn)|_W$ soit inversible et que tous les morphismes de
multiplication
$\mathcal{O}_X(nd)|_W \otimes_{\mathcal{O}_W} \mathcal{O}_X(md)|_W
\to \mathcal{O}_X(nd + md)|_W$
(voir \ref{equation-multiply}) soient des isomorphismes ;
\item la réunion des ouverts $D_{+}(fg)$, où $f, g \in S$ sont homogènes et
$\deg(f) = \deg(g) + d$.
\end{enumerate}
En outre, tous les morphismes
$\widetilde M(nd)|_W = \widetilde M|_W \otimes_{\mathcal{O}_W}
\mathcal{O}_X(nd)|_W \to \widetilde{M(nd)}|_W$
(voir \ref{equation-multiply-more-generally}) sont des isomorphismes.
\end{lemma}

\begin{proof}
Si $x \in D_{+}(fg)$ avec $\deg(f) = \deg(g) + d$, alors, sur
$D_{+}(fg)$, les faisceaux $\mathcal{O}_X(dn)$ sont engendrés par l'élément
$(f/g)^n = f^{2n}/(fg)^n$. En raisonnant comme dans la démonstration du
Lemme \ref{lemma-when-invertible}, on en déduit que $x$ appartient à l'ouvert
$W$ défini en (1).

\medskip\noindent
Réciproquement, supposons que $\mathcal{O}_X(dn)$ soit libre de rang $1$ pour
tout $n$, sur un voisinage ouvert $V$ de $x \in X$, et que tous les morphismes
de multiplication
$\mathcal{O}_X(nd)|_V \otimes_{\mathcal{O}_V} \mathcal{O}_X(md)|_V
\to \mathcal{O}_X(nd + md)|_V$ soient des isomorphismes.
On peut choisir $h \in S_{+}$ homogène tel que
$x \in D_{+}(h) \subset V$. Par définition des tordus du faisceau structural,
il existe un élément $s$ de $(S_h)_d$ tel que $s^n$ soit une base de
$(S_h)_{nd}$ comme $S_{(h)}$-module pour tout $n \in \mathbf{Z}$.
Écrivons $s = f/h^m$ avec $m \geq 1$ et
$f \in S_{d + m \deg(h)}$. Posons $g = h^m$, de sorte que $s = f/g$.
Par construction, $x \in D_{+}(g)$. Remarquons que
$g^d \in (S_h)_{d\deg(g)}$. Par hypothèse, cet élément est un multiple de
$s^{\deg(g)} = f^{\deg(g)}/g^{\deg(g)}$ ; écrivons
$g^d = a/g^e \cdot f^{\deg(g)}/g^{\deg(g)}$. On en déduit
$g^{d + e + \deg(g)} = a f^{\deg(g)}$, puis $x \in D_{+}(f)$.
Ainsi, $x$ appartient à l'ensemble défini en (2).

\medskip\noindent
L'existence, sur l'ouvert affine $D_{+}(fg)$, de la section génératrice
$s = f/g$ dont les puissances engendrent librement les faisceaux de modules
$\mathcal{O}_X(nd)$ montre aisément que les morphismes
$\widetilde M(nd)|_W = \widetilde M|_W \otimes_{\mathcal{O}_W}
\mathcal{O}_X(nd)|_W \to \widetilde{M(nd)}|_W$
(voir \ref{equation-multiply-more-generally}) sont des isomorphismes.
Comparer avec la démonstration du Lemme \ref{lemma-when-invertible}.
\end{proof}

\noindent
Rappelons que, d'après Modules, Lemme \ref{modules-lemma-s-open}, si
$\mathcal{L}$ est un faisceau inversible sur un espace localement annelé $X$
et si $s$ est une section globale de $\mathcal{L}$, alors l'ensemble
$X_s = \{x \in X \mid s \not \in \mathfrak m_x\mathcal{L}_x\}$ est ouvert.

\begin{lemma}
\label{lemma-principal-open}
Soit $S$ un anneau gradué, posons $X = \text{Proj}(S)$ et fixons un entier
$d \geq 1$. Soit $W = W_d \subset X$ le sous-schéma ouvert défini dans le
Lemme \ref{lemma-where-invertible}. Soient $n \geq 1$ et $f \in S_{nd}$.
Notons $s \in \Gamma(W, \mathcal{O}_W(nd))$ l'image de $f$ par
(\ref{equation-global-sections}), restreinte à $W$. Alors
$$
W_s = D_{+}(f) \cap W.
$$
\end{lemma}

\begin{proof}
Soit $D_{+}(ab) \subset W$ un ouvert affine standard, où $a, b \in S$ sont
homogènes et $\deg(a) = \deg(b) + d$. Remarquons que
$D_{+}(ab) \cap D_{+}(f) = D_{+}(abf)$. D'autre part, la restriction de $s$
à $D_{+}(ab)$ correspond à l'élément
$f/1 = b^nf/a^n (a/b)^n \in (S_{ab})_{nd}$.
Nous avons vu dans la démonstration du Lemme \ref{lemma-where-invertible} que
$(a/b)^n$ engendre $\mathcal{O}_W(nd)$ au-dessus de $D_{+}(ab)$.
Il s'ensuit que $W_s \cap D_{+}(ab)$ est l'ouvert principal associé à
$b^nf/a^n \in \mathcal{O}_X(D_{+}(ab))$. Le résultat est alors clair.
\end{proof}

\noindent
Le lemme suivant énonce les propriétés que nous utiliserons plus loin pour
caractériser les schémas munis d'un faisceau inversible ample.

\begin{lemma}
\label{lemma-ample-on-proj}
Soit $S$ un anneau gradué, soit $X = \text{Proj}(S)$ et soit
$Y \subset X$ un sous-schéma ouvert quasi-compact. Notons
$\mathcal{O}_Y(n)$ la restriction de $\mathcal{O}_X(n)$ à $Y$.
Il existe un entier $d \geq 1$ tel que
\begin{enumerate}
\item le sous-schéma $Y$ est contenu dans l'ouvert $W_d$ défini au
Lemme \ref{lemma-where-invertible} ;
\item le faisceau $\mathcal{O}_Y(dn)$ est inversible pour tout
$n \in \mathbf{Z}$ ;
\item tous les morphismes
$\mathcal{O}_Y(nd) \otimes_{\mathcal{O}_Y} \mathcal{O}_Y(m)
\longrightarrow
\mathcal{O}_Y(nd + m)$
de l'Équation (\ref{equation-multiply}) sont des isomorphismes ;
\item tous les morphismes
$\widetilde M(nd)|_Y = \widetilde M|_Y \otimes_{\mathcal{O}_Y}
\mathcal{O}_X(nd)|_Y \to \widetilde{M(nd)}|_Y$
(voir \ref{equation-multiply-more-generally}) sont des isomorphismes ;
\item si $f \in S_{nd}$ et si $s \in \Gamma(Y, \mathcal{O}_Y(nd))$ désigne
l'image de $f$ par (\ref{equation-global-sections}), restreinte à $Y$, alors
$D_{+}(f) \cap Y = Y_s$ ;
\item les ouverts $Y_s$, où $s \in \Gamma(Y, \mathcal{O}_Y(nd))$ et
$n \geq 1$, forment une base de la topologie de $Y$ ;
\item les ouverts affines $Y_s \subset Y$, où
$s \in \Gamma(Y, \mathcal{O}_Y(nd))$ et $n \geq 1$, forment une base de la
topologie de $Y$.
\end{enumerate}
\end{lemma}

\begin{proof}
Comme $Y$ est quasi-compact, il existe un nombre fini d'éléments homogènes
$f_i \in S_{+}$, $i = 1, \ldots, n$, tels que les ouverts standards
$D_{+}(f_i)$ recouvrent $Y$. Posons $d_i = \deg(f_i)$ et
$d = d_1 \ldots d_n$. Comme
$D_{+}(f_i) = D_{+}(f_i^{d/d_i})$, la caractérisation (2) du
Lemme \ref{lemma-where-invertible}, ou bien (1) et le
Lemme \ref{lemma-when-invertible}, donne immédiatement $Y \subset W_d$.
Le Lemme \ref{lemma-where-invertible} montre alors que (1) entraîne (2), (3)
et (4). Remarquons que (3) est un cas particulier de (4).
L'assertion (5) résulte du Lemme \ref{lemma-principal-open}.
Enfin, (6) et (7) résultent du fait que les ouverts $D_{+}(f)$ forment une
base de la topologie de $X$ et sont affines.
\end{proof}

\begin{lemma}
\label{lemma-comparison-proj-quasi-coherent}
Soit $S$ un anneau gradué, posons $X = \text{Proj}(S)$ et soit
$\mathcal{F}$ un $\mathcal{O}_X$-module quasi-cohérent. Munissons
$M = \bigoplus_{n \in \mathbf{Z}} \Gamma(X, \mathcal{F}(n))$ de sa structure
de $S$-module gradué à l'aide de
(\ref{equation-global-sections-module}) et (\ref{equation-global-sections}).
Il existe alors un morphisme canonique de $\mathcal{O}_X$-modules
$$
\widetilde{M} \longrightarrow \mathcal{F}
$$
fonctoriel en $\mathcal{F}$, tel que l'application induite
$M_0 \to \Gamma(X, \mathcal{F})$ soit l'identité.
\end{lemma}

\begin{proof}
Soit $f \in S$ homogène de degré $d > 0$. Rappelons que, d'après le
Lemme \ref{lemma-proj-sheaves}, la restriction
$\widetilde{M}|_{D_{+}(f)}$ correspond au $S_{(f)}$-module $M_{(f)}$.
On peut donc définir une application canonique
$$
M_{(f)} \longrightarrow \Gamma(D_+(f), \mathcal{F}),\quad
m/f^n \longmapsto m|_{D_+(f)} \otimes f|_{D_+(f)}^{-n},
$$
qui a un sens parce que $f|_{D_+(f)}$ est une section trivialisante du
faisceau inversible $\mathcal{O}_X(d)|_{D_+(f)}$ ; voir le
Lemme \ref{lemma-when-invertible} et sa démonstration. Comme $\widetilde{M}$
est quasi-cohérent, on obtient ainsi un morphisme canonique
$$
\widetilde{M}|_{D_+(f)} \longrightarrow \mathcal{F}|_{D_+(f)}
$$
par Schémas, Lemme \ref{schemes-lemma-compare-constructions}.
Ces morphismes se recollent sur les intersections et définissent donc un
morphisme global. Nous omettons la vérification du recollement ainsi que la
démonstration de la dernière assertion.
\end{proof}


\section{Fonctorialité de Proj}
\label{section-functoriality-proj}

\noindent
Un homomorphisme d'anneaux gradués $\psi : A \to B$ ne donne pas toujours lieu à
un morphisme entre les spectres homogènes projectifs associés. La raison en est que
l'image réciproque $\psi^{-1}(\mathfrak q)$
d'un idéal premier homogène $\mathfrak q \subset B$ peut
contenir l'idéal $A_{+}$ des éléments de degré strictement positif même si $\mathfrak q$ ne
contient pas $B_{+}$.
L'énoncé correct est le suivant.

\begin{lemma}
\label{lemma-morphism-proj}
Soient $A$, $B$ deux anneaux gradués.
Posons $X = \text{Proj}(A)$ et $Y = \text{Proj}(B)$.
Soit $\psi : A \to B$ un homomorphisme d'anneaux gradués.
Posons
$$
U(\psi)
=
\bigcup\nolimits_{f \in A_{+}\ \text{homogène}} D_{+}(\psi(f))
\subset Y.
$$
Il existe alors un morphisme canonique de schémas
$$
r_\psi :
U(\psi)
\longrightarrow
X
$$
et un homomorphisme de $\mathcal{O}_{U(\psi)}$-algèbres $\mathbf{Z}$-graduées
$$
\theta = \theta_\psi :
r_\psi^*\left(
\bigoplus\nolimits_{d \in \mathbf{Z}} \mathcal{O}_X(d)
\right)
\longrightarrow
\bigoplus\nolimits_{d \in \mathbf{Z}} \mathcal{O}_{U(\psi)}(d).
$$
Le triplet $(U(\psi), r_\psi, \theta)$ est
caractérisé par les propriétés suivantes :
\begin{enumerate}
\item Pour tout $d \geq 0$, le diagramme
$$
\xymatrix{
A_d \ar[d] \ar[rr]_{\psi} & &
B_d \ar[d] \\
\Gamma(X, \mathcal{O}_X(d)) \ar[r]^-\theta &
\Gamma(U(\psi), \mathcal{O}_Y(d)) &
\Gamma(Y, \mathcal{O}_Y(d)) \ar[l]
}
$$
est commutatif.
\item Pour tout $f \in A_{+}$ homogène,
on a $r_\psi^{-1}(D_{+}(f)) = D_{+}(\psi(f))$ et
la restriction de $r_\psi$ à $D_{+}(\psi(f))$
correspond à l'homomorphisme d'anneaux
$A_{(f)} \to B_{(\psi(f))}$ induit par $\psi$.
\end{enumerate}
\end{lemma}

\begin{proof}
Il est clair que la condition (2) détermine de manière unique le morphisme de schémas
et l'ouvert $U(\psi)$. Fixons $f \in A_d$ avec $d \geq 1$.
Notons que
$\mathcal{O}_X(n)|_{D_{+}(f)}$ correspond au
$A_{(f)}$-module $(A_f)_n$ et que
$\mathcal{O}_Y(n)|_{D_{+}(\psi(f))}$ correspond au
$B_{(\psi(f))}$-module $(B_{\psi(f)})_n$. Autrement dit, la restriction de $\theta$
à $D_{+}(\psi(f))$ correspond à un homomorphisme d'algèbres
$\mathbf{Z}$-graduées sur $B_{(\psi(f))}$
$$
A_f \otimes_{A_{(f)}} B_{(\psi(f))}
\longrightarrow
B_{\psi(f)}
$$
La condition (1) détermine les images de tous les éléments de $A$.
Puisque $f$ est un élément inversible envoyé sur $\psi(f)$,
on voit que $1/f^m$ est envoyé sur $1/\psi(f)^m$. Il en résulte aussitôt
que $\theta$ est uniquement déterminé, à savoir qu'il est
donné par la formule
$$
a/f^m \otimes b/\psi(f)^e \longmapsto \psi(a)b/\psi(f)^{m + e}.
$$
Pour démontrer l'existence, remarquons que la preuve d'unicité ci-dessus fournit
une prescription bien définie pour le morphisme $r$ et l'homomorphisme $\theta$
sur tout ouvert standard de la forme
$D_{+}(\psi(f)) \subset U(\psi)$ à valeurs dans $D_{+}(f)$.
Notons-les $r_f$ et $\theta_f$.
Il suffit donc de vérifier que, si $D_{+}(f) \subset D_{+}(g)$
pour des éléments homogènes $f, g \in A_{+}$, alors la restriction de
$r_g$ à $D_{+}(\psi(f))$ coïncide avec $r_f$. Cela résulte clairement des
formules données ci-dessus pour $r$ et $\theta$.
\end{proof}

\begin{lemma}
\label{lemma-morphism-proj-transitive}
Soient $A$, $B$ et $C$ des anneaux gradués.
Posons $X = \text{Proj}(A)$, $Y = \text{Proj}(B)$ et $Z = \text{Proj}(C)$.
Soient $\varphi : A \to B$, $\psi : B \to C$ des homomorphismes d'anneaux gradués.
Alors on a
$$
U(\psi \circ \varphi) = r_\psi^{-1}(U(\varphi))
\quad
\text{et}
\quad
r_{\psi \circ \varphi}
=
r_\varphi \circ r_\psi|_{U(\psi \circ \varphi)}.
$$
De plus, on a
$$
\theta_\psi \circ r_\psi^*\theta_\varphi
=
\theta_{\psi \circ \varphi}
$$
avec les notations évidentes.
\end{lemma}

\begin{proof}
La démonstration est omise.
\end{proof}

\begin{lemma}
\label{lemma-surjective-graded-rings-map-proj}
Avec les hypothèses et les notations du lemme \ref{lemma-morphism-proj} ci-dessus.
Supposons que $A_d \to B_d$ soit surjectif pour tout $d \gg 0$. Alors
\begin{enumerate}
\item $U(\psi) = Y$,
\item $r_\psi : Y \to X$ est une immersion fermée, et
\item les homomorphismes $\theta : r_\psi^*\mathcal{O}_X(n) \to \mathcal{O}_Y(n)$
sont surjectifs, mais ne sont pas des isomorphismes en général (même si $A \to B$ est
surjectif).
\end{enumerate}
\end{lemma}

\begin{proof}
La partie (1) résulte de la définition de $U(\psi)$ et du fait que
$D_{+}(f) = D_{+}(f^n)$ pour tout $n > 0$. Pour $f \in A_{+}$ homogène,
on voit que $A_{(f)} \to B_{(\psi(f))}$ est surjectif car
tout élément de $B_{(\psi(f))}$ peut être représenté par une fraction
$b/\psi(f)^n$ avec $n$ arbitrairement grand (ce qui force le degré de
$b \in B$ à être grand). Cela démontre (2).
Le même argument montre que l'homomorphisme
$$
A_f \to B_{\psi(f)}
$$
est surjectif, ce qui démontre la surjectivité de $\theta$.
Pour un exemple où cet homomorphisme n'est pas un isomorphisme,
considérons l'anneau gradué $A = k[x, y]$, où $k$ est un corps,
avec $\deg(x) = 1$, $\deg(y) = 2$. Posons $I = (x)$, de sorte que
$B = k[y]$. Notons que $\mathcal{O}_Y(1) = 0$ dans ce cas.
Mais on voit facilement que $r_\psi^*\mathcal{O}_X(1)$
n'est pas nul. (Il existe des exemples moins artificiels.)
\end{proof}

\begin{lemma}
\label{lemma-eventual-iso-graded-rings-map-proj}
Avec les hypothèses et les notations du lemme \ref{lemma-morphism-proj} ci-dessus.
Supposons que $A_d \to B_d$ soit un isomorphisme pour tout $d \gg 0$. Alors
\begin{enumerate}
\item $U(\psi) = Y$,
\item $r_\psi : Y \to X$ est un isomorphisme, et
\item les homomorphismes $\theta : r_\psi^*\mathcal{O}_X(n) \to \mathcal{O}_Y(n)$
sont des isomorphismes.
\end{enumerate}
\end{lemma}

\begin{proof}
On a (1) par le lemme \ref{lemma-surjective-graded-rings-map-proj}.
Soit $f \in A_{+}$ homogène. L'hypothèse sur $\psi$ implique que
$A_f \to B_f$ est un isomorphisme (détails omis). Il est donc clair que
les restrictions de $r_\psi$ et de $\theta$ au-dessus de $D_{+}(f)$ sont des isomorphismes.
Le lemme en résulte.
\end{proof}

\begin{lemma}
\label{lemma-surjective-graded-rings-generated-degree-1-map-proj}
Avec les hypothèses et les notations du lemme \ref{lemma-morphism-proj} ci-dessus.
Supposons que $A_d \to B_d$ soit surjectif pour $d \gg 0$ et que $A$ soit engendré
par $A_1$ comme $A_0$-algèbre. Alors
\begin{enumerate}
\item $U(\psi) = Y$,
\item $r_\psi : Y \to X$ est une immersion fermée, et
\item les homomorphismes $\theta : r_\psi^*\mathcal{O}_X(n) \to \mathcal{O}_Y(n)$
sont des isomorphismes.
\end{enumerate}
\end{lemma}

\begin{proof}
Par les lemmes \ref{lemma-eventual-iso-graded-rings-map-proj} et
\ref{lemma-morphism-proj-transitive},
on peut remplacer $B$ par l'image de $A \to B$
sans modifier $Y$ ni les faisceaux $\mathcal{O}_Y(n)$.
On peut donc supposer que $A \to B$ est surjectif. Par le
lemme \ref{lemma-surjective-graded-rings-map-proj}, on obtient (1) et (2),
ainsi que la surjectivité dans (3).
Par le lemme \ref{lemma-apply-modules}, on voit que
$\mathcal{O}_X(n)$ et $\mathcal{O}_Y(n)$
sont tous deux inversibles. Ainsi, $\theta$ est un isomorphisme.
\end{proof}

\begin{lemma}
\label{lemma-base-change-map-proj}
\begin{reference}
\cite[II, Proposition 2.8.10]{EGA}
\end{reference}
Avec les hypothèses et les notations du lemme \ref{lemma-morphism-proj} ci-dessus.
Supposons qu'il existe un homomorphisme d'anneaux $R \to A_0$ et un homomorphisme d'anneaux
$R \to R'$ tels que $B = R' \otimes_R A$. Alors
\begin{enumerate}
\item $U(\psi) = Y$,
\item le diagramme
$$
\xymatrix{
Y = \text{Proj}(B) \ar[r]_{r_\psi} \ar[d] &
\text{Proj}(A) = X \ar[d] \\
\Spec(R') \ar[r] &
\Spec(R)
}
$$
est un carré cartésien, et
\item les homomorphismes $\theta : r_\psi^*\mathcal{O}_X(n) \to \mathcal{O}_Y(n)$
sont des isomorphismes.
\end{enumerate}
\end{lemma}

\begin{proof}
Cela résulte immédiatement de l'examen de ce qui se passe sur les ouverts
standards $D_{+}(f)$ pour $f \in A_{+}$.
\end{proof}

\begin{lemma}
\label{lemma-localization-map-proj}
Avec les hypothèses et les notations du lemme \ref{lemma-morphism-proj} ci-dessus.
Supposons qu'il existe $g \in A_0$ tel que $\psi$ induise un
isomorphisme $A_g \to B$. Alors
$U(\psi) = Y$, $r_\psi : Y \to X$ est une immersion ouverte
qui identifie $Y$ à l'image réciproque
de $D(g) \subset \Spec(A_0)$. De plus, l'homomorphisme $\theta$
est un isomorphisme.
\end{lemma}

\begin{proof}
C'est un cas particulier du lemme \ref{lemma-base-change-map-proj} ci-dessus.
\end{proof}

\begin{lemma}
\label{lemma-d-uple}
Soit $S$ un anneau gradué. Soit $d \geq 1$. Posons $S' = S^{(d)}$ avec les notations
d'Algèbre, section \ref{algebra-section-graded}. Posons
$X = \text{Proj}(S)$ et $X' = \text{Proj}(S')$. Il existe un
isomorphisme canonique de schémas $i : X \to X'$ tel que
\begin{enumerate}
\item pour tout $S$-module gradué $M$, en posant $M' = M^{(d)}$,
on ait un isomorphisme canonique $\widetilde{M} \to i^*\widetilde{M'}$,
\item on ait des isomorphismes canoniques
$\mathcal{O}_{X}(nd) \to i^*\mathcal{O}_{X'}(n)$
\end{enumerate}
et ces isomorphismes sont compatibles avec les homomorphismes de multiplication
du lemme \ref{lemma-widetilde-tensor}, et donc avec les homomorphismes
(\ref{equation-multiply}),
(\ref{equation-multiply-on-sheaf}),
(\ref{equation-global-sections}),
(\ref{equation-global-sections-module}),
(\ref{equation-multiply-more-generally}) et
(\ref{equation-global-sections-more-generally}) (voir la démonstration pour les
énoncés précis).
\end{lemma}

\begin{proof}
L'homomorphisme injectif d'anneaux $S' \to S$ (qui n'est pas un homomorphisme d'anneaux gradués
en raison de nos conventions) induit un morphisme $j : \Spec(S) \to \Spec(S')$.
Étant donné un idéal premier homogène $\mathfrak p \subset S$, on voit que
$\mathfrak p' = j(\mathfrak p) = S' \cap \mathfrak p$
est un idéal premier homogène de $S'$.
De plus, si $f \in S_+$ est homogène et $f \not \in \mathfrak p$, alors
$f^d \in S'_+$ et $f^d \not \in \mathfrak p'$. Réciproquement, si
$\mathfrak p' \subset S'$ est un idéal premier homogène ne contenant pas un
élément homogène $f \in S'_+$, alors
$\mathfrak p = \{g \in S \mid g^d \in \mathfrak p'\}$ est un
idéal premier homogène de $S$ ne contenant pas $f$, dont l'image par $j$
est $\mathfrak p'$. Pour voir que $\mathfrak p$ est un idéal, notons
que, si $g, h \in \mathfrak p$, alors
$(g + h)^{2d} \in \mathfrak p'$ par la formule du binôme.
Puisque $\mathfrak p'$ est premier, on a donc $g + h \in \mathfrak p$
par définition de $\mathfrak p$.
On voit ainsi que $j$ induit un homéomorphisme $i : X \to X'$.
De plus, étant donné $f \in S_+$ homogène, on a
$S_{(f)} \cong S'_{(f^d)}$. Comme ces isomorphismes sont compatibles
avec les homomorphismes de restriction du
lemme \ref{lemma-standard-open}, on voit qu'il existe un
isomorphisme $i^\sharp : i^{-1}\mathcal{O}_{X'} \to \mathcal{O}_X$ entre
les faisceaux structuraux ; ainsi, $i$ est un isomorphisme
de schémas.

\medskip\noindent
Soit $M$ un $S$-module gradué. Étant donné $f \in S_+$ homogène, on a
$M_{(f)} \cong M'_{(f^d)}$ ; exactement de la même manière que ci-dessus,
on obtient donc l'isomorphisme de (1). Les isomorphismes de (2) sont un cas particulier
de (1) pour $M = S(nd)$, ce qui donne $M' = S'(n)$. Soient $M$ et $N$
des $S$-modules gradués. Alors on a
$$
M' \otimes_{S'} N' =
(M \otimes_S N)^{(d)} =
(M \otimes_S N)'
$$
comme on le vérifie directement à partir des définitions. Cela étant,
la compatibilité avec les homomorphismes de multiplication du
lemme \ref{lemma-widetilde-tensor} équivaut à la commutativité du diagramme
$$
\xymatrix{
\widetilde M \otimes_{\mathcal{O}_X} \widetilde N
\ar[d]_{(1) \otimes (1)} \ar[r] &
\widetilde{M \otimes_S N} \ar[d]^{(1)} \\
i^*(\widetilde{M'} \otimes_{\mathcal{O}_{X'}} \widetilde{N'}) \ar[r] &
i^*(\widetilde{M' \otimes_{S'} N'})
}
$$
Cela se voit en examinant la construction des homomorphismes
sur l'ouvert $D_+(f) = D_+(f^d)$, où la flèche horizontale supérieure
est donnée par l'application
$M_{(f)} \times N_{(f)} \to (M \otimes_S N)_{(f)}$
et la flèche horizontale inférieure par l'application
$M'_{(f^d)} \times N'_{(f^d)} \to (M' \otimes_{S'} N')_{(f^d)}$.
Puisque ces applications coïncident via les identifications
$M_{(f)} = M'_{(f^d)}$, etc., on obtient la compatibilité voulue.
Nous omettons la démonstration des autres compatibilités.
\end{proof}


\section{Morphismes vers Proj}
\label{section-morphisms-proj}

\noindent
Soit $S$ un anneau gradué. Soit $X = \text{Proj}(S)$ le spectre homogène de
$S$ et soit $d \geq 1$ un entier. Considérons le sous-schéma ouvert
\begin{equation}
\label{equation-Ud}
U_d = \bigcup\nolimits_{f  \in S_d} D_{+}(f)
\quad\subset\quad
X = \text{Proj}(S).
\end{equation}
Remarquons que $d | d' \Rightarrow U_d \subset U_{d'}$ et que
$X = \bigcup_d U_d$. Ni $X$ ni $U_d$ ne sont nécessairement quasi-compacts ;
voir Algèbre, Lemme \ref{algebra-lemma-topology-proj}. Écrivons
$\mathcal{O}_{U_d}(n) = \mathcal{O}_X(n)|_{U_d}$. D'après le
Lemme \ref{lemma-when-invertible}, $\mathcal{O}_{U_d}(nd)$ est un
$\mathcal{O}_{U_d}$-module inversible pour tout $n \in \mathbf{Z}$, et tous
les morphismes de multiplication
$\mathcal{O}_{U_d}(nd) \otimes_{\mathcal{O}_{U_d}} \mathcal{O}_{U_d}(m)
\to \mathcal{O}_{U_d}(nd + m)$ de (\ref{equation-multiply}) sont des
isomorphismes. En particulier,
$\mathcal{O}_{U_d}(nd) \cong \mathcal{O}_{U_d}(d)^{\otimes n}$.
L'homomorphisme d'anneaux gradués (\ref{equation-global-sections}), composé
avec la restriction à $U_d$, donne un homomorphisme d'anneaux gradués
\begin{equation}
\label{equation-psi-d}
\psi^d : S^{(d)} \longrightarrow \Gamma_*(U_d, \mathcal{O}_{U_d}(d)).
\end{equation}
Pour la notation $S^{(d)}$, voir Algèbre,
Section \ref{algebra-section-graded}. Pour la notation $\Gamma_*$, voir
Modules, Définition \ref{modules-definition-gamma-star}. En outre, puisque les
ouverts $D_{+}(f)$, $f \in S_d$, recouvrent $U_d$, le faisceau
$\mathcal{O}_{U_d}(d)$ est engendré par les sections appartenant à l'image de
$\psi^d_1 : S^{(d)}_1 = S_d \to \Gamma(U_d, \mathcal{O}_{U_d}(d))$ ; voir
Modules, Définition \ref{modules-definition-globally-generated}.

\medskip\noindent
Soit $Y$ un schéma et soit $\varphi : Y \to X$ un morphisme de schémas.
Supposons que l'image $\varphi(Y)$ soit contenue dans le sous-schéma ouvert
$U_d$ de $X$. La discussion qui suit Modules,
Définition \ref{modules-definition-gamma-star}, fournit un homomorphisme
d'anneaux gradués
$$
\Gamma_*(U_d, \mathcal{O}_{U_d}(d))
\longrightarrow
\Gamma_*(Y, \varphi^*\mathcal{O}_X(d)).
$$
Sa composée avec $\psi^d$ donne un homomorphisme d'anneaux gradués
\begin{equation}
\label{equation-psi-phi-d}
\psi_\varphi^d :
S^{(d)}
\longrightarrow
\Gamma_*(Y, \varphi^*\mathcal{O}_X(d))
\end{equation}
tel que le faisceau inversible $\varphi^*\mathcal{O}_X(d)$ soit engendré par
les sections appartenant à l'image de
$(S^{(d)})_1 = S_d \to \Gamma(Y, \varphi^*\mathcal{O}_X(d))$.

\begin{lemma}
\label{lemma-converse-construction}
Soient $S$ un anneau gradué, $X = \text{Proj}(S)$, $d \geq 1$ et
$U_d \subset X$ comme ci-dessus. Soit $Y$ un schéma, soit $\mathcal{L}$ un
faisceau inversible sur $Y$ et soit
$\psi : S^{(d)} \to \Gamma_*(Y, \mathcal{L})$ un homomorphisme d'anneaux
gradués tel que $\mathcal{L}$ soit engendré par les sections appartenant à
l'image de $\psi|_{S_d} : S_d \to \Gamma(Y, \mathcal{L})$.
Il existe alors un morphisme $\varphi : Y \to X$ tel que
$\varphi(Y) \subset U_d$ et un isomorphisme
$\alpha : \varphi^*\mathcal{O}_{U_d}(d) \to \mathcal{L}$ tels que
$\psi_\varphi^d$ coïncide avec $\psi$ par l'intermédiaire de $\alpha$ :
$$
\xymatrix{
\Gamma_*(Y, \mathcal{L}) &
\Gamma_*(Y, \varphi^*\mathcal{O}_{U_d}(d)) \ar[l]^-\alpha &
\Gamma_*(U_d, \mathcal{O}_{U_d}(d)) \ar[l]^-{\varphi^*} \\
S^{(d)} \ar[u]^\psi & &
S^{(d)} \ar[u]^{\psi^d} \ar[ul]^{\psi^d_\varphi} \ar[ll]_{\text{id}}
}
$$
est commutatif. De plus, le couple $(\varphi, \alpha)$ est unique.
\end{lemma}

\begin{proof}
Choisissons $f \in S_d$ et notons
$s = \psi(f) \in \Gamma(Y, \mathcal{L})$. Sur l'ouvert $Y_s$ où $s$ ne
s'annule pas, la multiplication par $s$ induit un isomorphisme
$\mathcal{O}_{Y_s} \to \mathcal{L}|_{Y_s}$ ; voir Modules,
Lemme \ref{modules-lemma-s-open}. Nous noterons l'inverse de cet isomorphisme
$x \mapsto x/s$, et procéderons de même pour les puissances de
$\mathcal{L}$. On définit ainsi un homomorphisme d'anneaux
$\psi_{(f)} : S_{(f)} \to \Gamma(Y_s, \mathcal{O}_Y)$ en envoyant la fraction
$a/f^n$ sur $\psi(a)/s^n$. Par Schémas,
Lemme \ref{schemes-lemma-morphism-into-affine}, cet homomorphisme correspond à
un morphisme $\varphi_f : Y_s \to \Spec(S_{(f)}) = D_{+}(f)$.
Introduisons aussi l'isomorphisme
$\alpha_f : \varphi_f^*\mathcal{O}_{D_{+}(f)}(d) \to \mathcal{L}|_{Y_s}$ qui
envoie l'image inverse de la section trivialisante $f$ sur $D_{+}(f)$ vers la
section trivialisante $s$ sur $Y_s$. Avec ce choix, le diagramme du lemme est
commutatif lorsque $Y$, $\varphi$ et $\alpha$ sont remplacés respectivement
par $Y_s$, $\varphi_f$ et $\alpha_f$ ; nous omettons la vérification.

\medskip\noindent
Soit $f' \in S_d$ un second élément et notons
$s' = \psi(f') \in \Gamma(Y, \mathcal{L})$. Alors
$Y_s \cap Y_{s'} = Y_{ss'}$, et de même
$D_{+}(f) \cap D_{+}(f') = D_{+}(ff')$. Nous avons vu au
Lemme \ref{lemma-ample-on-proj} que $D_{+}(f') \cap D_{+}(f)$ est l'ensemble
des points de $D_{+}(f)$ où la section de $\mathcal{O}_X(d)$ définie par $f'$
ne s'annule pas. Ainsi,
$\varphi_f^{-1}(D_{+}(f') \cap D_{+}(f)) = Y_s \cap Y_{s'}
= \varphi_{f'}^{-1}(D_{+}(f') \cap D_{+}(f))$.
Sur $D_{+}(f) \cap D_{+}(f')$, la fraction $f/f'$ est une section inversible
du faisceau structural, d'inverse $f'/f$. Remarquons que
$\psi_{(f')}(f/f') = \psi(f)/s' = s/s'$ et
$\psi_{(f)}(f'/f) = \psi(f')/s = s'/s$. Nous affirmons qu'il existe un unique
homomorphisme d'anneaux
$S_{(ff')} \to \Gamma(Y_{ss'}, \mathcal{O}_Y)$ rendant commutatif le diagramme
$$
\xymatrix{
\Gamma(Y_s, \mathcal{O}_Y) \ar[r] &
\Gamma(Y_{ss'}, \mathcal{O}_Y) &
\Gamma(Y_{s'}, \mathcal{O}_Y) \ar[l]\\
S_{(f)} \ar[r] \ar[u]^{\psi_{(f)}} &
S_{(ff')} \ar[u] &
S_{(f')} \ar[l] \ar[u]^{\psi_{(f')}}
}
$$
Il existe, car on peut utiliser la règle
$x/(ff')^n \mapsto \psi(x)/(ss')^n$, qui est bien définie d'après les formules
précédentes. L'unicité résulte de ce que $\text{Proj}(S)$ est séparé ; voir le
Lemme \ref{lemma-proj-separated} et sa démonstration. Ainsi, les morphismes
$\varphi_f$ et $\varphi_{f'}$ coïncident sur $Y_s \cap Y_{s'}$.
Les restrictions de $\alpha_f$ et $\alpha_{f'}$ y coïncident aussi, puisque
les fonctions régulières $s/s'$ et $\psi_{(f')}(f/f')$ coïncident.
Les morphismes $\varphi_f$ se recollent donc en un morphisme global de $Y$
vers $U_d \subset X$, et les morphismes $\alpha_f$ se recollent en un
isomorphisme satisfaisant les conditions du lemme.

\medskip\noindent
Il reste à montrer l'unicité du couple $(\varphi, \alpha)$. Supposons que
$(\varphi', \alpha')$ soit un second couple satisfaisant les mêmes conditions.
Soit $f \in S_d$. La commutativité des diagrammes du lemme montre que les
images inverses de $D_{+}(f)$ par $\varphi$ et par $\varphi'$ sont toutes deux
égales à $Y_{\psi(f)}$. Comme les ouverts $D_{+}(f)$ forment une base de la
topologie de $X$ et que l'espace topologique $X$ est sobre (voir Schémas,
Lemme \ref{schemes-lemma-scheme-sober}), les applications continues sous-jacentes
à $\varphi$ et à $\varphi'$ sont égales. Utilisons $s = \psi(f)$ pour
trivialiser le faisceau inversible $\mathcal{L}$ sur $Y_{\psi(f)}$.
La commutativité des diagrammes donne
$\alpha^{\otimes n}(\psi^d_{\varphi}(x)) =
\psi(x) = (\alpha')^{\otimes n}(\psi^d_{\varphi'}(x))$
pour tout $x \in S_{nd}$. Par construction de $\psi^d_{\varphi}$ et de
$\psi^d_{\varphi'}$, on a
$\psi^d_{\varphi}(x) = \varphi^\sharp(x/f^n) \psi^d_{\varphi}(f^n)$
sur $Y_{\psi(f)}$, et de même pour $\psi^d_{\varphi'}$.
La commutativité des diagrammes du lemme entraîne alors
$\varphi^\sharp(x/f^n) = (\varphi')^\sharp(x/f^n)$.
Ainsi, $\varphi$ et $\varphi'$ induisent le même morphisme de
$Y_{\psi(f)}$ vers le schéma affine $D_{+}(f) = \Spec(S_{(f)})$ ; ce sont donc
les mêmes morphismes. Enfin, il reste à montrer que, $\varphi$ étant donné, la
commutativité du diagramme du lemme détermine un unique $\alpha$. Nous omettons
la vérification.
\end{proof}

\noindent
Reprenons la discussion qui précède le lemme. Soit $S$ un anneau gradué et
soit $Y$ un schéma. Nous considérerons les {\it triplets}
$(d, \mathcal{L}, \psi)$ où
\begin{enumerate}
\item $d \geq 1$ est un entier ;
\item $\mathcal{L}$ est un $\mathcal{O}_Y$-module inversible ;
\item $\psi : S^{(d)} \to \Gamma_*(Y, \mathcal{L})$ est un homomorphisme
d'anneaux gradués tel que $\mathcal{L}$ soit engendré par les sections
globales $\psi(f)$, avec $f \in S_d$.
\end{enumerate}
Étant donnés un morphisme $h : Y' \to Y$ et un triplet
$(d, \mathcal{L}, \psi)$ sur $Y$, on peut en prendre l'image inverse : c'est
le triplet $(d, h^*\mathcal{L}, h^* \circ \psi)$.
Deux triplets $(d, \mathcal{L}, \psi)$ et
$(d, \mathcal{L}', \psi')$ associés au même entier $d$ sont dits
{\it strictement équivalents} s'il existe un isomorphisme
$\beta : \mathcal{L} \to \mathcal{L}'$ tel que
$\beta \circ \psi = \psi'$ comme homomorphismes d'anneaux gradués
$S^{(d)} \to \Gamma_*(Y, \mathcal{L}')$.

\medskip\noindent
Pour tout entier $d \geq 1$, nous définissons
\begin{eqnarray*}
F_d : \Sch^{opp} & \longrightarrow & \textit{Ens}, \\
Y & \longmapsto &
\{\text{classes d'équivalence stricte de triplets }
(d, \mathcal{L}, \psi)
\text{ comme ci-dessus}\}
\end{eqnarray*}
les images inverses étant définies comme ci-dessus.

\begin{lemma}
\label{lemma-proj-functor-strict}
Soit $S$ un anneau gradué et soit $X = \text{Proj}(S)$. Le sous-schéma ouvert
$U_d \subset X$ défini en (\ref{equation-Ud}) représente le foncteur $F_d$,
et le triplet $(d, \mathcal{O}_{U_d}(d), \psi^d)$ défini ci-dessus est la
famille universelle (voir Schémas,
Section \ref{schemes-section-representable}).
\end{lemma}

\begin{proof}
C'est une reformulation du Lemme \ref{lemma-converse-construction}.
\end{proof}

\begin{lemma}
\label{lemma-apply}
Soit $S$ un anneau gradué engendré, comme $S_0$-algèbre, par les éléments de
$S_1$. Alors le schéma $X = \text{Proj}(S)$ représente le foncteur qui associe
à un schéma $Y$ l'ensemble des couples $(\mathcal{L}, \psi)$ où
\begin{enumerate}
\item $\mathcal{L}$ est un $\mathcal{O}_Y$-module inversible ;
\item $\psi : S \to \Gamma_*(Y, \mathcal{L})$ est un homomorphisme d'anneaux
gradués tel que $\mathcal{L}$ soit engendré par les sections globales
$\psi(f)$, avec $f \in S_1$,
\end{enumerate}
à équivalence stricte près, au sens défini ci-dessus.
\end{lemma}

\begin{proof}
Sous les hypothèses du lemme, on a $X = U_1$ ; le résultat est donc une
reformulation du Lemme \ref{lemma-proj-functor-strict}.
\end{proof}

\noindent
Terminons cette section par la description d'un foncteur correspondant à
$\text{Proj}(S)$ pour un anneau gradué quelconque $S$. Nous conseillons au
lecteur de ne pas lire la suite de cette section.

\medskip\noindent
Fixons un anneau gradué quelconque $S$ et soit $T$ un schéma. Deux triplets
$(d, \mathcal{L}, \psi)$ et $(d', \mathcal{L}', \psi')$ sur $T$, où les
entiers $d$ et $d'$ peuvent être distincts, sont dits {\it équivalents} s'il
existe un isomorphisme de faisceaux inversibles sur $T$
$\beta : \mathcal{L}^{\otimes d'} \to (\mathcal{L}')^{\otimes d}$ tel que
$\beta \circ \psi|_{S^{(dd')}}$ et $\psi'|_{S^{(dd')}}$ coïncident comme
homomorphismes d'anneaux gradués
$S^{(dd')} \to \Gamma_*(T, (\mathcal{L}')^{\otimes d})$.

\begin{lemma}
\label{lemma-equivalent}
Soit $S$ un anneau gradué, posons $X = \text{Proj}(S)$ et soit $T$ un schéma.
Soient $(d, \mathcal{L}, \psi)$ et $(d', \mathcal{L}', \psi')$ deux triplets
sur $T$. Les assertions suivantes sont équivalentes :
\begin{enumerate}
\item Posons $n = \text{lcm}(d, d')$ et écrivons $n = ad = a'd'$.
Il existe un isomorphisme
$\beta : \mathcal{L}^{\otimes a} \to (\mathcal{L}')^{\otimes a'}$ tel que
$\beta \circ \psi|_{S^{(n)}}$ et $\psi'|_{S^{(n)}}$ coïncident comme
homomorphismes d'anneaux gradués
$S^{(n)} \to \Gamma_*(T, (\mathcal{L}')^{\otimes a'})$.
\item Les triplets $(d, \mathcal{L}, \psi)$ et
$(d', \mathcal{L}', \psi')$ sont équivalents.
\item Pour un certain entier strictement positif $n = ad = a'd'$, il existe un
isomorphisme
$\beta : \mathcal{L}^{\otimes a} \to (\mathcal{L}')^{\otimes a'}$ tel que
$\beta \circ \psi|_{S^{(n)}}$ et $\psi'|_{S^{(n)}}$ coïncident comme
homomorphismes d'anneaux gradués
$S^{(n)} \to \Gamma_*(T, (\mathcal{L}')^{\otimes a'})$.
\item Les morphismes $\varphi : T \to X$ et $\varphi' : T \to X$ associés
respectivement à $(d, \mathcal{L}, \psi)$ et
$(d', \mathcal{L}', \psi')$ sont égaux.
\end{enumerate}
\end{lemma}

\begin{proof}
Il est clair que (1) entraîne (2), et (2) entraîne (3) en se restreignant à
des degrés plus divisibles et en prenant des puissances des faisceaux
inversibles. En outre, (3) entraîne (4) par l'assertion d'unicité du
Lemme \ref{lemma-converse-construction}. Il reste donc à montrer que (4)
entraîne (1). Supposons (4), c'est-à-dire $\varphi = \varphi'$.
On peut alors écrire
$\mathcal{L} = \varphi^*\mathcal{O}_X(d)$ et
$\mathcal{L}' = \varphi^*\mathcal{O}_X(d')$. En outre, par ces identifications,
les homomorphismes d'anneaux gradués $\psi$ et $\psi'$ correspondent à la
restriction de l'homomorphisme canonique d'anneaux gradués
$$
S \longrightarrow \bigoplus\nolimits_{n \geq 0} \Gamma(X, \mathcal{O}_X(n))
$$
à $S^{(d)}$ et à $S^{(d')}$, composée avec l'image inverse par $\varphi$
(encore une fois par le Lemme \ref{lemma-converse-construction}). Il suffit
donc de prendre pour $\beta$ l'isomorphisme
$$
(\varphi^*\mathcal{O}_X(d))^{\otimes a} =
\varphi^*\mathcal{O}_X(n) =
(\varphi^*\mathcal{O}_X(d'))^{\otimes a'}.
$$
\end{proof}

\noindent
Soit $S$ un anneau gradué et soit $X = \text{Proj}(S)$. Sur le sous-schéma
ouvert $U_d \subset X = \text{Proj}(S)$ défini en (\ref{equation-Ud}), nous
disposons du triplet $(d, \mathcal{O}_{U_d}(d), \psi^d)$. Si $d | d'$, les
triplets $(d, \mathcal{O}_{U_d}(d), \psi^d)$ et
$(d', \mathcal{O}_{U_{d'}}(d'), \psi^{d'})$, restreints à l'ouvert $U_d$
(qui est contenu dans $U_{d'}$), sont manifestement équivalents. Ce fait et le
Lemme \ref{lemma-converse-construction} montrent que les morphismes
$Y \to X$ correspondent approximativement aux classes d'équivalence de
triplets sur $Y$. Cette affirmation n'est pas tout à fait exacte : si $Y$
n'est pas quasi-compact, il se peut qu'aucun triplet unique ne convienne.
Il faut donc prendre quelques précautions pour définir le foncteur
correspondant.

\medskip\noindent
Voici une façon de procéder. Supposons $d' = ad$. Considérons la transformation
de foncteurs $F_d \to F_{d'}$ qui associe au triplet
$(d, \mathcal{L}, \psi)$ sur $T$ le triplet
$(d', \mathcal{L}^{\otimes a}, \psi|_{S^{(d')}})$.
L'une des implications du Lemme \ref{lemma-equivalent} montre que cette
transformation $F_d \to F_{d'}$ est injective. Pour un schéma quasi-compact
$T$, posons
$$
F(T) = \bigcup\nolimits_{d \in \mathbf{N}} F_d(T).
$$
Les applications de transition sont celles décrites ci-dessus. On obtient
ainsi un foncteur contravariant de la catégorie des schémas quasi-compacts vers
la catégorie des ensembles. Pour un schéma quelconque $T$, posons
$$
F(T)
=
\lim_{V \subset T\text{ ouvert quasi-compact}} F(V).
$$
Autrement dit, un élément $\xi$ de $F(T)$ correspond au choix compatible
d'éléments $\xi_V \in F(V)$, où $V$ parcourt les ouverts quasi-compacts de
$T$. Nous omettons la définition de l'application d'image inverse
$F(T) \to F(T')$ associée à un morphisme de schémas $T' \to T$.
Nous avons ainsi défini le foncteur
\begin{eqnarray*}
F : \Sch^{opp} & \longrightarrow & \textit{Ens}
\end{eqnarray*}

\begin{lemma}
\label{lemma-proj-functor}
Soit $S$ un anneau gradué et soit $X = \text{Proj}(S)$. Le foncteur $F$ défini
ci-dessus est représentable par le schéma $X$.
\end{lemma}

\begin{proof}
Nous avons vu que le foncteur $F_d$ correspond au sous-schéma ouvert
$U_d \subset X$. En outre, la transformation de foncteurs
$F_d \to F_{d'}$ (si $d | d'$) définie ci-dessus correspond au morphisme
d'inclusion $U_d \to U_{d'}$ ; voir la discussion précédente. Pour montrer
que $F$ est représenté par $X$, il suffit donc de vérifier que tout morphisme
$T \to X$, où $T$ est quasi-compact, se factorise par un certain $U_d$, et que,
pour tout schéma $T$, on a
$$
\Mor(T, X)
=
\lim_{V \subset T\text{ ouvert quasi-compact}} \Mor(V, X).
$$
Nous omettons ces vérifications.
\end{proof}


\section{Espace projectif}
\label{section-projective-space}

\noindent
L'espace projectif est l'un des objets fondamentaux de la géométrie
algébrique. Dans cette section, nous nous contentons de le construire comme le
$\text{Proj}$ d'un anneau de polynômes. Nous découvrirons plus loin nombre de
ses belles propriétés.

\begin{lemma}
\label{lemma-projective-space}
Soit $S = \mathbf{Z}[T_0, \ldots, T_n]$ avec $\deg(T_i) = 1$. Le schéma
$$
\mathbf{P}^n_{\mathbf{Z}} = \text{Proj}(S)
$$
représente le foncteur qui associe à un schéma $Y$ les couples
$(\mathcal{L}, (s_0, \ldots, s_n))$ tels que
\begin{enumerate}
\item $\mathcal{L}$ soit un $\mathcal{O}_Y$-module inversible ;
\item $s_0, \ldots, s_n$ soient des sections globales de $\mathcal{L}$ qui
engendrent $\mathcal{L}$,
\end{enumerate}
à l'équivalence suivante près :
$(\mathcal{L}, (s_0, \ldots, s_n)) \sim
(\mathcal{N}, (t_0, \ldots, t_n))$ si et seulement s'il existe un
isomorphisme $\beta : \mathcal{L} \to \mathcal{N}$ tel que
$\beta(s_i) = t_i$ pour $i = 0, \ldots, n$.
\end{lemma}

\begin{proof}
C'est un cas particulier du Lemme \ref{lemma-apply}. En effet, pour tout
anneau gradué $A$, on a
\begin{eqnarray*}
\Mor_{\text{anneaux gradués}}(\mathbf{Z}[T_0, \ldots, T_n], A)
& = &
A_1 \times \ldots \times A_1 \\
\psi & \mapsto & (\psi(T_0), \ldots, \psi(T_n))
\end{eqnarray*}
et la partie de degré $1$ de $\Gamma_*(Y, \mathcal{L})$ n'est autre que
$\Gamma(Y, \mathcal{L})$.
\end{proof}

\begin{definition}
\label{definition-projective-space}
Le schéma
$\mathbf{P}^n_{\mathbf{Z}} = \text{Proj}(\mathbf{Z}[T_0, \ldots, T_n])$ est
appelé l'{\it espace projectif de dimension $n$ sur $\mathbf{Z}$}.
Son changement de base $\mathbf{P}^n_S$ à un schéma $S$ est appelé l'
{\it espace projectif de dimension $n$ sur $S$}. Si $R$ est un anneau, le
changement de base à $\Spec(R)$ est noté $\mathbf{P}^n_R$ et appelé l'
{\it espace projectif de dimension $n$ sur $R$}.
\end{definition}

\noindent
Soient $Y$ un schéma sur $S$ et
$(\mathcal{L}, (s_0, \ldots, s_n))$ un couple comme dans le
Lemme \ref{lemma-projective-space}. Le morphisme induit vers
$\mathbf{P}^n_S$ est noté
$$
\varphi_{(\mathcal{L}, (s_0, \ldots, s_n))} :
Y \longrightarrow \mathbf{P}^n_S.
$$
Cette notation a un sens : le couple définit un morphisme vers
$\mathbf{P}^n_{\mathbf{Z}}$, tandis que le morphisme structural vers $S$ est
déjà donné ; ensemble, ils définissent un morphisme vers
$\mathbf{P}^n_S = \mathbf{P}^n_{\mathbf{Z}} \times S$.
Par le Lemme \ref{lemma-relative-proj-base-change}, on peut identifier
$\mathbf{P}^n_S$ au Proj relatif de la $\mathcal{O}_S$-algèbre graduée
$\mathcal{O}_S[T_0, \ldots, T_n]$. On peut alors caractériser
$\varphi_{(\mathcal{L}, (s_0, \ldots, s_n))}$ comme l'unique $S$-morphisme tel
que
$$
\mathcal{L} =
\varphi_{(\mathcal{L}, (s_0, \ldots, s_n))}^*\mathcal{O}_{\mathbf{P}^n_S}(1)
\quad
\text{et}
\quad
s_i = \varphi_{(\mathcal{L}, (s_0, \ldots, s_n))}^*T_i,
$$
où l'on considère $T_i$ comme une section globale de
$\mathcal{O}_{\mathbf{P}^n_S}(1)$ par l'intermédiaire de l'application $\psi$
du Lemme \ref{lemma-glue-relative-proj-twists} ; nous omettons les détails.

\begin{lemma}
\label{lemma-standard-covering-projective-space}
L'espace projectif de dimension $n$ sur $\mathbf{Z}$ est recouvert par les
$n + 1$ ouverts standards
$$
\mathbf{P}^n_{\mathbf{Z}} =
\bigcup\nolimits_{i = 0, \ldots, n} D_{+}(T_i),
$$
chacun des $D_{+}(T_i)$ étant isomorphe à l'espace affine de dimension $n$
$\mathbf{A}^n_{\mathbf{Z}}$ sur $\mathbf{Z}$.
\end{lemma}

\begin{proof}
Cela résulte de
$\mathbf{Z}[T_0, \ldots, T_n]_{+} = (T_0, \ldots, T_n)$ et de l'isomorphisme
évident
$$
\Spec
\left(
\mathbf{Z}
\left[\frac{T_0}{T_i}, \ldots, \frac{T_n}{T_i}
\right]
\right)
\cong
\mathbf{A}^n_{\mathbf{Z}}.
$$
\end{proof}

\begin{lemma}
\label{lemma-projective-space-separated}
Soit $S$ un schéma. Le morphisme structural $\mathbf{P}^n_S \to S$
possède les propriétés suivantes :
\begin{enumerate}
\item il est séparé ;
\item il est quasi-compact ;
\item il satisfait les parties existence et unicité du critère valuatif ;
\item il est universellement fermé.
\end{enumerate}
\end{lemma}

\begin{proof}
Toutes ces propriétés sont stables par changement de base (c'est clair pour
les deux dernières ; pour les deux premières, voir Schémas, Lemmes
\ref{schemes-lemma-separated-permanence} et
\ref{schemes-lemma-quasi-compact-preserved-base-change}). Il suffit donc de les
démontrer pour le morphisme
$\mathbf{P}^n_{\mathbf{Z}} \to \Spec(\mathbf{Z})$.
La séparation est le Lemme \ref{lemma-proj-separated}. La quasi-compacité
résulte du Lemme \ref{lemma-standard-covering-projective-space}.
Les parties existence et unicité du critère valuatif résultent du
Lemme \ref{lemma-proj-valuative-criterion}. Enfin, le caractère universellement
fermé découle des assertions précédentes et de Schémas,
Proposition \ref{schemes-proposition-characterize-universally-closed}.
\end{proof}

\begin{remark}
\label{remark-missing-finite-type}
Que manque-t-il à la liste de propriétés ci-dessus ? Certainement la propriété
d'être de type fini. Si nous ne l'énonçons pas ici, c'est que nous n'avons pas
encore défini la notion de type fini à ce stade. Une autre propriété absente
est la « lissité » ; il y en a sans doute beaucoup d'autres auxquelles le
lecteur pourra penser.
\end{remark}

\begin{lemma}[Plongement de Segre]
\label{lemma-segre-embedding}
Soit $S$ un schéma. Il existe une immersion fermée
$$
\mathbf{P}^n_S \times_S \mathbf{P}^m_S
\longrightarrow
\mathbf{P}^{nm + n + m}_S
$$
appelée le {\it plongement de Segre}.
\end{lemma}

\begin{proof}
Il suffit de le démontrer lorsque $S = \Spec(\mathbf{Z})$. Nous omettrons donc
l'indice $S$ et travaillerons dans la situation absolue. Écrivons
$\mathbf{P}^n = \text{Proj}(\mathbf{Z}[X_0, \ldots, X_n])$,
$\mathbf{P}^m = \text{Proj}(\mathbf{Z}[Y_0, \ldots, Y_m])$ et
$\mathbf{P}^{nm + n + m} =
\text{Proj}(\mathbf{Z}[Z_0, \ldots, Z_{nm + n + m}])$.
Pour définir un morphisme vers $\mathbf{P}^{nm + n + m}$, il faut se donner
sur le membre de gauche un faisceau inversible $\mathcal{L}$ et
$(n + 1)(m + 1)$ sections $s_i$ qui l'engendrent ; voir le
Lemme \ref{lemma-projective-space}. Prenons le faisceau inversible
$$
\mathcal{L} =
\text{pr}_1^*\mathcal{O}_{\mathbf{P}^n}(1)
\otimes
\text{pr}_2^*\mathcal{O}_{\mathbf{P}^m}(1)
$$
et les sections
$$
s_0 = X_0Y_0, \ s_1 = X_1Y_0, \ldots, \ s_n = X_nY_0,
\ s_{n + 1} = X_0Y_1, \ldots, \ s_{nm + n + m} = X_nY_m.
$$
Elles engendrent $\mathcal{L}$, puisque les sections $X_i$ engendrent
$\mathcal{O}_{\mathbf{P}^n}(1)$ et les sections $Y_j$ engendrent
$\mathcal{O}_{\mathbf{P}^m}(1)$. Le morphisme induit $\varphi$ vérifie
$$
\varphi^{-1}(D_{+}(Z_{i + (n + 1)j})) = D_{+}(X_i) \times D_{+}(Y_j).
$$
C'est donc un morphisme affine. Dans le cas $(i, j) = (0, 0)$,
l'homomorphisme d'anneaux correspondant est
$$
\mathbf{Z}[Z_1/Z_0, \ldots, Z_{nm + n + m}/Z_0]
\longrightarrow
\mathbf{Z}[X_1/X_0, \ldots, X_n/X_0, Y_1/Y_0, \ldots, Y_m/Y_0]
$$
qui envoie $Z_i/Z_0$ sur $X_i/X_0$ pour $i \leq n$ et
$Z_{(n + 1)j}/Z_0$ sur $Y_j/Y_0$. Il est donc surjectif. Le même raisonnement
s'applique aux autres ouverts affines. Ainsi, $\varphi$ est une immersion
fermée (voir Schémas, Lemme \ref{schemes-lemma-closed-local-target} et
Exemple \ref{schemes-example-closed-immersion-affines}).
\end{proof}

\noindent
Les deux lemmes suivants sont des cas particuliers de résultats plus généraux
établis plus loin, mais il est sans doute utile de les démontrer directement
ici. Le premier est un cas particulier de Diviseurs,
Lemme \ref{divisors-lemma-closed-subscheme-proj}.

\begin{lemma}
\label{lemma-closed-in-projective-space}
Soit $R$ un anneau et soit $Z \subset \mathbf{P}^n_R$ un sous-schéma fermé.
Posons
$$
I_d = \Ker\left(
R[T_0, \ldots, T_n]_d
\longrightarrow
\Gamma(Z, \mathcal{O}_{\mathbf{P}^n_R}(d)|_Z)\right).
$$
Alors $I = \bigoplus I_d \subset R[T_0, \ldots, T_n]$ est un idéal gradué et
$Z = \text{Proj}(R[T_0, \ldots, T_n]/I)$.
\end{lemma}

\begin{proof}
Il est clair que $I$ est un idéal gradué. Posons
$Z' = \text{Proj}(R[T_0, \ldots, T_n]/I)$. D'après le
Lemme \ref{lemma-surjective-graded-rings-generated-degree-1-map-proj}, $Z'$ est
un sous-schéma fermé de $\mathbf{P}^n_R$. Pour vérifier l'égalité $Z = Z'$,
il suffit de se placer sur un ouvert affine standard $D_{+}(T_i)$.
Après renumérotation des coordonnées homogènes, on peut supposer $i = 0$.
Écrivons que $Z \cap D_{+}(T_0)$, respectivement
$Z' \cap D_{+}(T_0)$, est défini par l'idéal $J$, respectivement $J'$, de
$R[T_1/T_0, \ldots, T_n/T_0]$. Alors $J'$ est engendré par les éléments
$F/T_0^{\deg(F)}$, où $F \in I$ est homogène.
Supposons que $F \in I$ soit de degré $d$. Comme $F$ s'annule en tant que
section de $\mathcal{O}_{\mathbf{P}^n_R}(d)$ restreinte à $Z$, l'élément
$F/T_0^d$ appartient à $J$. Ainsi, $J' \subset J$.

\medskip\noindent
Réciproquement, soit $f \in J$. Si $f$ est de degré total $d$ en
$T_1/T_0, \ldots, T_n/T_0$, on peut écrire
$f = F/T_0^d$ pour un certain $F \in R[T_0, \ldots, T_n]_d$.
Choisissons $i \in \{1, \ldots, n\}$. Le sous-schéma
$Z \cap D_{+}(T_i)$ est défini par un certain idéal
$J_i \subset R[T_0/T_i, \ldots, T_n/T_i]$. En outre,
$$
J \cdot
R\left[
\frac{T_1}{T_0}, \ldots, \frac{T_n}{T_0},
\frac{T_0}{T_i}, \ldots, \frac{T_n}{T_i}
\right]
=
J_i \cdot
R\left[
\frac{T_1}{T_0}, \ldots, \frac{T_n}{T_0},
\frac{T_0}{T_i}, \ldots, \frac{T_n}{T_i}
\right].
$$
Le membre de gauche est le localisé de $J$ par rapport à $T_i/T_0$, et celui
de droite le localisé de $J_i$ par rapport à $T_0/T_i$. Il s'ensuit que
$T_0^{d_i}F/T_i^{d + d_i}$ appartient à $J_i$ pour un entier $d_i$
suffisamment grand. Par conséquent, $T_0^{\max(d_i)}F$ appartient à $I$,
car sa restriction à chacun des ouverts affines standards $D_{+}(T_i)$
s'annule sur le sous-schéma fermé $Z \cap D_{+}(T_i)$.
Ainsi, $f \in J'$, d'où $J \subset J'$, comme voulu.
\end{proof}

\noindent
Le lemme suivant est un cas particulier des résultats plus généraux de
Propriétés, Lemmes \ref{properties-lemma-ample-quasi-coherent} ou
\ref{properties-lemma-proj-quasi-coherent}.

\begin{lemma}
\label{lemma-quasi-coherent-projective-space}
Soit $R$ un anneau et soit $\mathcal{F}$ un faisceau quasi-cohérent sur
$\mathbf{P}^n_R$. Pour $d \geq 0$, posons
$$
M_d
=
\Gamma(\mathbf{P}^n_R,
\mathcal{F} \otimes_{\mathcal{O}_{\mathbf{P}^n_R}}
\mathcal{O}_{\mathbf{P}^n_R}(d))
=
\Gamma(\mathbf{P}^n_R, \mathcal{F}(d)).
$$
Alors $M = \bigoplus_{d \geq 0} M_d$ est un
$R[T_0, \ldots, T_n]$-module gradué et il existe un isomorphisme canonique
$\mathcal{F} = \widetilde{M}$.
\end{lemma}

\begin{proof}
Les applications de multiplication
$$
R[T_0, \ldots, T_n]_e \times M_d \longrightarrow M_{d + e}
$$
proviennent des isomorphismes naturels
$$
\mathcal{O}_{\mathbf{P}^n_R}(e)
\otimes_{\mathcal{O}_{\mathbf{P}^n_R}}
\mathcal{F}(d)
\longrightarrow
\mathcal{F}(e + d) ;
$$
voir l'Équation (\ref{equation-global-sections-module}). Construisons le
morphisme $c : \widetilde{M} \to \mathcal{F}$. Sur chacun des ouverts affines
standards $U_i = D_{+}(T_i)$, on a
$\Gamma(U_i, \widetilde{M}) = (M[1/T_i])_0$, où l'indice ${}_0$ désigne la
partie de degré $0$. Tout élément de ce module s'écrit $m/T_i^d$ avec
$m \in M_d$. Comme $T_i$ engendre $\mathcal{O}(1)$ au-dessus de $U_i$, on
peut toujours écrire $m|_{U_i} = m_i \otimes T_i^d$, avec une unique section
$m_i \in \Gamma(U_i, \mathcal{F})$. Il est donc naturel de poser
$c(m/T_i^d) = m_i$. Un raisonnement élémentaire, omis ici, montre que cela
définit un morphisme $c : \widetilde{M} \to \mathcal{F}$ dès que l'on vérifie
$$
(T_i/T_j)^d m_i|_{U_i \cap U_j} = m_j|_{U_i \cap U_j}
$$
dans $M[1/T_iT_j]$. Or cela est clair, car sur l'intersection les générateurs
$T_i$ et $T_j$ de $\mathcal{O}(1)$ diffèrent par la fonction inversible
$T_i/T_j$.

\medskip\noindent
Injectivité de $c$. On peut la vérifier sur les ouverts affines $U_i$.
Soit $i \in \{0, \ldots, n\}$ et soit
$s = m/T_i^d \in \Gamma(U_i, \widetilde{M})$ tel que
$c(m/T_i^d) = 0$. D'après la description de $c$, on a $m_i = 0$, donc
$m|_{U_i} = 0$. Il existe ainsi un entier $e$ tel que $T_i^em = 0$ dans $M$.
Par conséquent,
$s = m/T_i^d = T_i^e m/T_i^{e + d} = 0$, comme voulu.

\medskip\noindent
Surjectivité de $c$. On peut la vérifier sur les ouverts affines $U_i$ et,
après renumérotation, il suffit de la vérifier sur $U_0$.
Soit $s \in \mathcal{F}(U_0)$. Écrivons
$\mathcal{F}|_{U_i} = \widetilde{N_i}$, où $N_i$ est un
$R[T_0/T_i, \ldots, T_n/T_i]$-module ; c'est possible parce que
$\mathcal{F}$ est quasi-cohérent. La section $s$ correspond alors à un élément
$x \in N_0$. On a
$$
(N_i)_{T_j/T_i} \cong (N_j)_{T_i/T_j}
$$
où les indices désignent la localisation principale, comme modules sur
l'anneau
$$
R\left[
\frac{T_0}{T_i}, \ldots, \frac{T_n}{T_i},
\frac{T_0}{T_j}, \ldots, \frac{T_n}{T_j}
\right].
$$
Il existe donc un entier $d$ assez grand et des éléments
$s_i \in N_i$, $i = 1, \ldots, n$, tels que
$$
s = (T_i/T_0)^d s_i
$$
sur $U_0 \cap U_i$. Considérons ensuite la différence
$$
t_{ij} = s_i - (T_j/T_i)^d s_j
$$
sur $U_i \cap U_j$, $0 < i < j$. Par le choix des $s_i$, on a
$t_{ij}|_{U_0 \cap U_i \cap U_j} = 0$. Il existe donc un entier $e$ assez
grand tel que $(T_0/T_i)^et_{ij} = 0$. Posons
$s_i' = (T_0/T_i)^es_i$ et $s_0' = s$. On a alors
$$
s_a' = (T_b/T_a)^{e + d} s_b'
$$
sur $U_a \cap U_b$ pour tous $a, b$. C'est exactement la condition pour que
les éléments $s'_a$ se recollent en une section globale
$m \in \Gamma(\mathbf{P}^n_R, \mathcal{F}(e + d))$.
Par construction, $c(m/T_0^{e + d}) = s$. Le morphisme $c$ est donc surjectif,
ce qui achève la démonstration.
\end{proof}

\begin{lemma}
\label{lemma-globally-generated-omega-twist-1}
Soit $X$ un schéma. Soit $\mathcal{L}$ un faisceau inversible, et soient
$s_0, \ldots, s_n$ des sections globales de $\mathcal{L}$ qui l'engendrent.
Soit $\mathcal{F}$ le noyau du morphisme induit
$\mathcal{O}_X^{\oplus n + 1} \to \mathcal{L}$. Alors
$\mathcal{F} \otimes \mathcal{L}$ est engendré par ses sections globales.
\end{lemma}

\begin{proof}
Le résultat est même vrai lorsque $X$ est un espace localement annelé
quelconque. Le faisceau $\mathcal{F}$ est un $\mathcal{O}_X$-module localement
libre de rang fini $n$. Les éléments
$$
s_{ij} = (0, \ldots, 0, s_j, 0, \ldots, 0, -s_i, 0, \ldots, 0)
\in \Gamma(X, \mathcal{L}^{\oplus n + 1})
$$
où $s_j$ occupe la $i$-ième place et $-s_i$ la $j$-ième s'envoient sur zéro
dans $\mathcal{L}^{\otimes 2}$. Ainsi,
$s_{ij} \in \Gamma(X, \mathcal{F} \otimes_{\mathcal{O}_X} \mathcal{L})$.
Un calcul local montre que ces sections engendrent
$\mathcal{F} \otimes \mathcal{L}$.

\medskip\noindent
Autre démonstration. Considérons le morphisme
$\varphi : X \to \mathbf{P}^n_\mathbf{Z}$ associé au couple
$(\mathcal{L}, (s_0, \ldots, s_n))$. Comme l'image inverse de
$\mathcal{O}(1)$ est $\mathcal{L}$ et celle de $T_i$ est $s_i$, il suffit de
démontrer le lemme pour $\mathbf{P}^n_\mathbf{Z}$. Dans ce cas, le faisceau
$\mathcal{F}$ correspond au module gradué $M$ sur
$S = \mathbf{Z}[T_0, \ldots, T_n]$ qui s'insère dans la suite exacte courte
$$
0 \to M \to S^{\oplus n + 1} \to S(1) \to 0
$$
où le second morphisme est donné par $T_0, \ldots, T_n$.
L'assertion précédente se traduit alors par le fait que les éléments
$$
T_{ij} = (0, \ldots, 0, T_j, 0, \ldots, 0, -T_i, 0, \ldots, 0)
\in M(1)_0
$$
engendrent le module gradué $M(1)$ sur $S$. Nous omettons les détails.
\end{proof}


\section{Faisceaux inversibles et morphismes vers Proj}
\label{section-invertible-proj}

\noindent
Soit $T$ un schéma et soit $\mathcal{L}$ un faisceau inversible sur $T$.
Pour une section $s \in \Gamma(T, \mathcal{L})$, on note $T_s$ l'ouvert
des points où $s$ ne s'annule pas. Voir Modules,
Lemme \ref{modules-lemma-s-open}. On peut considérer le lemme suivant comme
une légère généralisation du Lemme \ref{lemma-apply}. C'est également une
généralisation du Lemme \ref{lemma-morphism-proj}.

\begin{lemma}
\label{lemma-invertible-map-into-proj}
Soit $A$ un anneau gradué et posons $X = \text{Proj}(A)$.
Soient $T$ un schéma, $\mathcal{L}$ un $\mathcal{O}_T$-module inversible et
$\psi : A \to \Gamma_*(T, \mathcal{L})$ un homomorphisme d'anneaux gradués.
Posons
$$
U(\psi) = \bigcup\nolimits_{f \in A_{+}\text{ homogène}} T_{\psi(f)}
$$
Le morphisme $\psi$ induit un morphisme canonique de schémas
$$
r_{\mathcal{L}, \psi} :
U(\psi) \longrightarrow X
$$
muni d'un morphisme de $\mathcal{O}_T$-algèbres $\mathbf{Z}$-graduées
$$
\theta :
r_{\mathcal{L}, \psi}^*\left(
\bigoplus\nolimits_{d \in \mathbf{Z}} \mathcal{O}_X(d)
\right)
\longrightarrow
\bigoplus\nolimits_{d \in \mathbf{Z}} \mathcal{L}^{\otimes d}|_{U(\psi)}.
$$
Le triplet $(U(\psi), r_{\mathcal{L}, \psi}, \theta)$ est caractérisé par
les propriétés suivantes :
\begin{enumerate}
\item Pour tout $f \in A_{+}$ homogène, on a
$r_{\mathcal{L}, \psi}^{-1}(D_{+}(f)) = T_{\psi(f)}$.
\item Pour tout $d \geq 0$, le diagramme
$$
\xymatrix{
A_d \ar[d]_{(\ref{equation-global-sections})} \ar[r]_{\psi} &
\Gamma(T, \mathcal{L}^{\otimes d}) \ar[d]^{\text{restriction}} \\
\Gamma(X, \mathcal{O}_X(d)) \ar[r]^{\theta} &
\Gamma(U(\psi), \mathcal{L}^{\otimes d})
}
$$
est commutatif.
\end{enumerate}
De plus, pour tout $d \geq 1$ et tout sous-schéma ouvert $V \subset T$ tels
que les sections de $\psi(A_d)$ engendrent $\mathcal{L}^{\otimes d}|_V$,
le morphisme $r_{\mathcal{L}, \psi}|_V$ coïncide avec le morphisme
$\varphi : V \to \text{Proj}(A)$ et le morphisme $\theta|_V$ coïncide avec
$\alpha : \varphi^*\mathcal{O}_X(d) \to \mathcal{L}^{\otimes d}|_V$,
où $(\varphi, \alpha)$ est le couple du Lemme
\ref{lemma-converse-construction} associé à
$\psi|_{A^{(d)}} : A^{(d)} \to \Gamma_*(V, \mathcal{L}^{\otimes d})$.
\end{lemma}

\begin{proof}
Supposons que deux triplets $(U, r : U \to X, \theta)$ et
$(U', r' : U' \to X, \theta')$ satisfassent (1) et (2).
La propriété (1) entraîne $U = U' = U(\psi)$ ainsi que l'égalité
$r = r'$ des applications continues sous-jacentes : les ouverts $D_{+}(f)$
forment en effet une base de la topologie de $X$, et $X$ est un espace
topologique sobre (voir Algèbre, Section \ref{algebra-section-proj} et
Schémas, Lemme \ref{schemes-lemma-scheme-sober}).
Soit $f \in A_{+}$ homogène. On a
$\Gamma(D_{+}(f), \bigoplus_{n \in \mathbf{Z}} \mathcal{O}_X(n)) = A_f$
comme algèbres $\mathbf{Z}$-graduées. Considérons les deux homomorphismes
d'anneaux $\mathbf{Z}$-gradués
$$
\theta, \theta' :
A_f
\longrightarrow
\Gamma(T_{\psi(f)}, \bigoplus \mathcal{L}^{\otimes n}).
$$
La multiplication par $f$ (resp.\ $\psi(f)$) est un isomorphisme dans le
membre de gauche (resp.\ de droite). On sait aussi, par (2), que
$\theta(x/1) = \theta'(x/1) = \psi(x)|_{T_{\psi(f)}}$ pour tout $x \in A$.
On en déduit aisément $\theta = \theta'$. En considérant les composantes
de degré $0$, on obtient $r^\sharp = (r')^\sharp$, c'est-à-dire
$r = r'$ comme morphismes de schémas. Cela prouve l'unicité.

\medskip\noindent
Passons à l'existence. Grâce à l'unicité que nous venons de démontrer,
il suffit de construire le couple $(r, \theta)$ localement sur $T$.
On peut donc supposer $T = \Spec(R)$ affine, $\mathcal{L} = \mathcal{O}_T$,
et qu'il existe $f \in A_{+}$ homogène tel que $\psi(f)$ engendre
$\mathcal{O}_T = \mathcal{O}_T^{\otimes \deg(f)}$.
Autrement dit, $\psi(f) = u \in R^*$ est inversible. Dans ce cas, $\psi$
est un homomorphisme d'anneaux gradués
$$
A \longrightarrow R[x] = \Gamma_*(T, \mathcal{O}_T)
$$
qui envoie $f$ sur $ux^{\deg(f)}$. Il se prolonge de façon unique en un
homomorphisme d'anneaux $\mathbf{Z}$-gradués
$\theta : A_f \to R[x, x^{-1}]$ en envoyant $1/f$ sur
$u^{-1}x^{-\deg(f)}$. En degré zéro, on obtient l'homomorphisme d'anneaux
$A_{(f)} \to R$ qui définit le morphisme
$r : T = \Spec(R) \to \Spec(A_{(f)}) = D_{+}(f) \subset X$.
Nous avons donc construit $(r, \theta)$ dans ce cas particulier.

\medskip\noindent
Montrons la dernière assertion du lemme. D'après le
Lemme \ref{lemma-converse-construction}, le morphisme qui y est construit
est l'unique morphisme rendant commutatif le diagramme de son énoncé.
La commutativité du diagramme du présent lemme entraîne celle du diagramme
restreint à $V$ et à $A^{(d)}$. D'où le résultat.
\end{proof}

\begin{remark}
\label{remark-not-in-invertible-locus}
Sous les hypothèses du Lemme \ref{lemma-invertible-map-into-proj}, l'image
du morphisme $r_{\mathcal{L}, \psi}$ n'est pas nécessairement contenue dans
le lieu où le faisceau $\mathcal{O}_X(1)$ est inversible. Voici un exemple.
Soit $k$ un corps et soit $S = k[A, B, C]$ gradué par
$\deg(A) = 1$, $\deg(B) = 2$, $\deg(C) = 3$.
Posons $X = \text{Proj}(S)$ et soit
$T = \mathbf{P}^2_k = \text{Proj}(k[X_0, X_1, X_2])$.
Rappelons que $\mathcal{L} = \mathcal{O}_T(1)$ est inversible et que
$\mathcal{O}_T(n) = \mathcal{L}^{\otimes n}$.
Considérons la composée $\psi$ des homomorphismes
$$
S \to k[X_0, X_1, X_2] \to \Gamma_*(T, \mathcal{L}).
$$
Le premier envoie $A$ sur $X_0$, $B$ sur $X_1^2$ et $C$ sur $X_2^3$ ;
le second est (\ref{equation-global-sections}). D'après le lemme, cette
composée correspond à un morphisme
$r_{\mathcal{L}, \psi} : T \to X = \text{Proj}(S)$, dont on voit aisément
qu'il est surjectif. Or, dans la Remarque \ref{remark-not-isomorphism},
nous avons montré que $\mathcal{O}_X(1)$ n'est pas inversible en tout point
de $X$.
\end{remark}


\section{Proj relatif par recollement}
\label{section-relative-proj-via-glueing}

\begin{situation}
\label{situation-relative-proj}
Ici, $S$ est un schéma et $\mathcal{A}$ une $\mathcal{O}_S$-algèbre graduée
quasi-cohérente.
\end{situation}

\noindent
Dans cette section, nous indiquons comment construire un morphisme de schémas
$$
\underline{\text{Proj}}_S(\mathcal{A}) \longrightarrow S
$$
en recollant les spectres homogènes $\text{Proj}(\Gamma(U, \mathcal{A}))$,
où $U$ parcourt les ouverts affines de $S$. Montrons d'abord que les spectres
homogènes des algèbres de sections de $\mathcal{A}$ sur les ouverts affines
forment une famille de schémas satisfaisant les conditions du
Lemme \ref{lemma-relative-glueing}.

\begin{lemma}
\label{lemma-proj-inclusion}
Plaçons-nous dans la Situation \ref{situation-relative-proj}.
Soient $U \subset U' \subset S$ des ouverts affines et posons
$A = \mathcal{A}(U)$, $A' = \mathcal{A}(U')$.
L'homomorphisme d'anneaux gradués $A' \to A$ induit un morphisme
$r : \text{Proj}(A) \to \text{Proj}(A')$, et le diagramme
$$
\xymatrix{
\text{Proj}(A) \ar[r] \ar[d] &
\text{Proj}(A') \ar[d] \\
U \ar[r] &
U'
}
$$
est cartésien. De plus, il existe des isomorphismes canoniques
$\theta : r^*\mathcal{O}_{\text{Proj}(A')}(n) \to
\mathcal{O}_{\text{Proj}(A)}(n)$ compatibles avec les morphismes de
multiplication.
\end{lemma}

\begin{proof}
Posons $R = \mathcal{O}_S(U)$ et $R' = \mathcal{O}_S(U')$.
L'homomorphisme $R \otimes_{R'} A' \to A$ est un isomorphisme puisque
$\mathcal{A}$ est quasi-cohérente (voir, par exemple, Schémas,
Lemme \ref{schemes-lemma-widetilde-pullback}). Le résultat découle donc du
Lemme \ref{lemma-base-change-map-proj}.
\end{proof}

\noindent
En particulier, le morphisme $\text{Proj}(A) \to \text{Proj}(A')$ du lemme
est une immersion ouverte.

\begin{lemma}
\label{lemma-transitive-proj}
Plaçons-nous dans la Situation \ref{situation-relative-proj}.
Soient $U \subset U' \subset U'' \subset S$ des ouverts affines et posons
$A = \mathcal{A}(U)$, $A' = \mathcal{A}(U')$, $A'' = \mathcal{A}(U'')$.
La composée des morphismes $r : \text{Proj}(A) \to \text{Proj}(A')$ et
$r' : \text{Proj}(A') \to \text{Proj}(A'')$ du
Lemme \ref{lemma-proj-inclusion} est le morphisme
$r'' : \text{Proj}(A) \to \text{Proj}(A'')$ du
Lemme \ref{lemma-proj-inclusion}. Une assertion analogue vaut pour les
isomorphismes $\theta$.
\end{lemma}

\begin{proof}
Cela résulte du Lemme \ref{lemma-morphism-proj-transitive}, puisque
$A'' \to A$ est la composée de $A'' \to A'$ et de $A' \to A$.
\end{proof}

\begin{lemma}
\label{lemma-glue-relative-proj}
Dans la Situation \ref{situation-relative-proj}, il existe un morphisme
de schémas
$$
\pi : \underline{\text{Proj}}_S(\mathcal{A}) \longrightarrow S
$$
ayant les propriétés suivantes :
\begin{enumerate}
\item pour tout ouvert affine $U \subset S$, il existe un isomorphisme
$i_U : \pi^{-1}(U) \to \text{Proj}(A)$, où $A = \mathcal{A}(U)$ ;
\item pour des ouverts affines $U \subset U' \subset S$, la composée
$$
\xymatrix{
\text{Proj}(A) \ar[r]^{i_U^{-1}} &
\pi^{-1}(U) \ar[rr]^{\text{inclusion}} & &
\pi^{-1}(U') \ar[r]^{i_{U'}} &
\text{Proj}(A')
}
$$
où $A = \mathcal{A}(U)$, $A' = \mathcal{A}(U')$, est l'immersion ouverte du
Lemme \ref{lemma-proj-inclusion}.
\end{enumerate}
\end{lemma}

\begin{proof}
Cela résulte immédiatement des Lemmes \ref{lemma-relative-glueing},
\ref{lemma-proj-inclusion} et \ref{lemma-transitive-proj}.
\end{proof}

\begin{lemma}
\label{lemma-glue-relative-proj-twists}
Dans la Situation \ref{situation-relative-proj}, le morphisme
$\pi : \underline{\text{Proj}}_S(\mathcal{A}) \to S$ du
Lemme \ref{lemma-glue-relative-proj} est muni de la structure supplémentaire
suivante. Il existe un faisceau quasi-cohérent de
$\mathcal{O}_{\underline{\text{Proj}}_S(\mathcal{A})}$-algèbres
$\mathbf{Z}$-graduées
$\bigoplus\nolimits_{n \in \mathbf{Z}}
\mathcal{O}_{\underline{\text{Proj}}_S(\mathcal{A})}(n)$
et un morphisme de $\mathcal{O}_S$-algèbres graduées
$$
\psi :
\mathcal{A}
\longrightarrow
\bigoplus\nolimits_{n \geq 0}
\pi_*\left(\mathcal{O}_{\underline{\text{Proj}}_S(\mathcal{A})}(n)\right)
$$
uniquement déterminés par la propriété suivante : pour tout ouvert affine
$U \subset S$, en posant $A = \mathcal{A}(U)$, il existe un isomorphisme
$$
\theta_U :
i_U^*\left(
\bigoplus\nolimits_{n \in \mathbf{Z}} \mathcal{O}_{\text{Proj}(A)}(n)
\right)
\longrightarrow
\left(
\bigoplus\nolimits_{n \in \mathbf{Z}}
\mathcal{O}_{\underline{\text{Proj}}_S(\mathcal{A})}(n)
\right)|_{\pi^{-1}(U)}
$$
de $\mathcal{O}_{\pi^{-1}(U)}$-algèbres $\mathbf{Z}$-graduées tel que le
diagramme
$$
\xymatrix{
A_n
\ar[rr]_\psi
\ar[dr]_-{(\ref{equation-global-sections})}
& &
\Gamma(\pi^{-1}(U),
\mathcal{O}_{\underline{\text{Proj}}_S(\mathcal{A})}(n)) \\
&
\Gamma(\text{Proj}(A),
\mathcal{O}_{\text{Proj}(A)}(n))
\ar[ru]_-{\theta_U}
&
}
$$
soit commutatif.
\end{lemma}

\begin{proof}
Nous allons appliquer le Lemme \ref{lemma-relative-glueing-sheaves} pour
recoller, sur le schéma $\underline{\text{Proj}}_S(\mathcal{A})$, les
faisceaux d'algèbres $\mathbf{Z}$-graduées
$\bigoplus_{n \in \mathbf{Z}} \mathcal{O}_{\text{Proj}(A)}(n)$,
où $A = \mathcal{A}(U)$ et $U \subset S$ est un ouvert affine.
Nous avons construit les données nécessaires dans le
Lemme \ref{lemma-proj-inclusion} et vérifié la condition (d) du
Lemme \ref{lemma-relative-glueing-sheaves} dans le
Lemme \ref{lemma-transitive-proj}. On obtient ainsi le faisceau de
$\mathcal{O}_{\underline{\text{Proj}}_S(\mathcal{A})}$-algèbres
$\mathbf{Z}$-graduées
$\bigoplus_{n \in \mathbf{Z}}
\mathcal{O}_{\underline{\text{Proj}}_S(\mathcal{A})}(n)$,
muni des isomorphismes $\theta_U$ pour tout ouvert affine $U \subset S$
et tout $n \in \mathbf{Z}$. Pour chaque ouvert affine $U \subset S$,
en posant $A = \mathcal{A}(U)$, on dispose d'un homomorphisme
$A \to \Gamma(\text{Proj}(A),
\bigoplus_{n \geq 0} \mathcal{O}_{\text{Proj}(A)}(n))$.
Le morphisme $\psi$ existe donc par fonctorialité du recollement relatif ;
voir la Remarque \ref{remark-relative-glueing-functorial}.
Le diagramme du lemme est commutatif par construction. Cette propriété
caractérise le faisceau de
$\mathcal{O}_{\underline{\text{Proj}}_S(\mathcal{A})}$-algèbres
$\mathbf{Z}$-graduées
$\bigoplus \mathcal{O}_{\underline{\text{Proj}}_S(\mathcal{A})}(n)$ :
la démonstration du Lemme \ref{lemma-morphism-proj} montre en effet que la
commutativité de ces diagrammes détermine de façon unique les morphismes
$\theta_U$. Nous omettons certains détails.
\end{proof}


\section{Le Proj relatif comme foncteur}
\label{section-relative-proj}

\noindent
Plaçons-nous dans la Situation \ref{situation-relative-proj}.
Ainsi, $S$ est un schéma et
$\mathcal{A} = \bigoplus_{d \geq 0} \mathcal{A}_d$ est une
$\mathcal{O}_S$-algèbre graduée quasi-cohérente. Nous allons donner une
version relative de la construction de $\text{Proj}$, en définissant un
foncteur que représentera le spectre homogène relatif. Nous construirons
ainsi un morphisme de schémas
$$
\underline{\text{Proj}}_S(\mathcal{A}) \longrightarrow S
$$
qui, au-dessus des ouverts affines de $S$, sera le spectre homogène d'un
anneau gradué. Nous suivrons la démarche adoptée pour le spectre relatif
dans la Section \ref{section-spec}. La méthode plus simple consistant à
recoller les schémas $\text{Proj}(A)$, où $A = \Gamma(U, \mathcal{A})$ et
$U \subset S$ est un ouvert affine, est exposée dans la
Section \ref{section-relative-proj-via-glueing}. Nous montrerons que les
deux constructions donnent des schémas isomorphes.

\medskip\noindent
Fixons pour l'instant un entier $d \geq 1$. Nous notons
$\mathcal{A}^{(d)} = \bigoplus_{n \geq 0} \mathcal{A}_{nd}$, par analogie
avec la notation d'Algèbre, Section \ref{algebra-section-graded}.
Soit $T$ un schéma. Considérons les {\it quadruplets
$(d, f : T \to S, \mathcal{L}, \psi)$ sur $T$} où
\begin{enumerate}
\item $d$ est l'entier fixé ci-dessus ;
\item $f : T \to S$ est un morphisme de schémas ;
\item $\mathcal{L}$ est un $\mathcal{O}_T$-module inversible ;
\item
$\psi : f^*\mathcal{A}^{(d)} \to \bigoplus_{n \geq 0}\mathcal{L}^{\otimes n}$
est un homomorphisme de $\mathcal{O}_T$-algèbres graduées tel que
$f^*\mathcal{A}_d \to \mathcal{L}$ soit surjectif.
\end{enumerate}
Étant donnés un morphisme $h : T' \to T$ et un quadruplet
$(d, f, \mathcal{L}, \psi)$ sur $T$, on peut en prendre l'image inverse,
qui est le quadruplet $(d, f \circ h, h^*\mathcal{L}, h^*\psi)$ sur $T'$.
Deux quadruplets $(d, f, \mathcal{L}, \psi)$ et
$(d, f', \mathcal{L}', \psi')$ sur $T$, associés au même entier $d$, sont
dits {\it strictement équivalents} si $f = f'$ et s'il existe un isomorphisme
$\beta : \mathcal{L} \to \mathcal{L}'$ tel que $\beta \circ \psi = \psi'$
comme morphismes de $\mathcal{O}_T$-algèbres graduées
$f^*\mathcal{A}^{(d)} \to \bigoplus_{n \geq 0} (\mathcal{L}')^{\otimes n}$.

\medskip\noindent
Pour tout entier $d \geq 1$, définissons
\begin{eqnarray*}
F_d : \Sch^{opp} & \longrightarrow & \textit{Ens}, \\
T & \longmapsto &
\{\text{classes d'équivalence stricte de }
(d, f : T \to S, \mathcal{L}, \psi)
\text{ comme ci-dessus}\}
\end{eqnarray*}
avec les images inverses définies ci-dessus.

\begin{lemma}
\label{lemma-proj-base-change}
Dans la Situation \ref{situation-relative-proj}, soit $d \geq 1$ et soit
$F_d$ le foncteur associé ci-dessus à $(S, \mathcal{A})$.
Soit $g : S' \to S$ un morphisme de schémas. Posons
$\mathcal{A}' = g^*\mathcal{A}$ et soit $F_d'$ le foncteur associé à
$(S', \mathcal{A}')$. Il existe un isomorphisme canonique de foncteurs
$$
F'_d \cong h_{S'} \times_{h_S} F_d
$$
\end{lemma}

\begin{proof}
Se donner un quadruplet
$(d, f' : T \to S', \mathcal{L}',
\psi' : (f')^*(\mathcal{A}')^{(d)} \to
\bigoplus_{n \geq 0} (\mathcal{L}')^{\otimes n})$
revient à se donner un quadruplet
$(d, f, \mathcal{L},
\psi : f^*\mathcal{A}^{(d)} \to
\bigoplus_{n \geq 0} \mathcal{L}^{\otimes n})$
et une factorisation $f = g \circ f'$. La correspondance est donnée par
$f = g \circ f'$, $\mathcal{L} = \mathcal{L}'$ et $\psi = \psi'$, via les
identifications
$(f')^*(\mathcal{A}')^{(d)} = (f')^*g^*(\mathcal{A}^{(d)}) =
f^*\mathcal{A}^{(d)}$. D'où le lemme.
\end{proof}

\begin{lemma}
\label{lemma-relative-proj-affine}
Dans la Situation \ref{situation-relative-proj}, soit $F_d$ le foncteur
associé ci-dessus à $(d, S, \mathcal{A})$. Si $S$ est affine, alors $F_d$
est représentable par le sous-schéma ouvert $U_d$ de (\ref{equation-Ud})
du schéma $\text{Proj}(\Gamma(S, \mathcal{A}))$.
\end{lemma}

\begin{proof}
Écrivons $S = \Spec(R)$ et $A = \Gamma(S, \mathcal{A})$.
Alors $A$ est une $R$-algèbre graduée et $\mathcal{A} = \widetilde A$.
Il nous faut identifier $F_d$ au foncteur $F_d^{triples}$ des triplets
défini dans la Section \ref{section-morphisms-proj}.

\medskip\noindent
Soit $(d, f : T \to S, \mathcal{L}, \psi)$ un quadruplet. On peut voir
$\psi$ comme un morphisme de $\mathcal{O}_S$-modules
$\mathcal{A}^{(d)} \to \bigoplus_{n \geq 0} f_*\mathcal{L}^{\otimes n}$.
Puisque $\mathcal{A}^{(d)}$ est quasi-cohérente, cela revient à se donner
un homomorphisme $R$-linéaire d'anneaux gradués
$A^{(d)} \to \Gamma(S, \bigoplus_{n \geq 0} f_*\mathcal{L}^{\otimes n})$.
Or
$\Gamma(S, \bigoplus_{n \geq 0} f_*\mathcal{L}^{\otimes n}) =
\Gamma_*(T, \mathcal{L})$.
On peut donc associer au quadruplet le triplet $(d, \mathcal{L}, \psi)$.

\medskip\noindent
Réciproquement, soit $(d, \mathcal{L}, \psi)$ un triplet.
La composée $R \to A_0 \to \Gamma(T, \mathcal{O}_T)$ détermine un morphisme
$f : T \to S = \Spec(R)$ ; voir Schémas,
Lemme \ref{schemes-lemma-morphism-into-affine}. Avec ce choix de $f$,
l'homomorphisme
$A^{(d)} \to \Gamma(S, \bigoplus_{n \geq 0} f_*\mathcal{L}^{\otimes n})$
est $R$-linéaire et correspond donc à un morphisme $\psi$ donnant un
quadruplet $(d, f : T \to S, \mathcal{L}, \psi)$.
Nous omettons la vérification que ces constructions définissent un
isomorphisme de foncteurs $F_d = F_d^{triples}$.
\end{proof}

\begin{lemma}
\label{lemma-relative-proj-d}
Dans la Situation \ref{situation-relative-proj}, le foncteur $F_d$ est
représentable par un schéma.
\end{lemma}

\begin{proof}
Nous allons appliquer Schémas, Lemme \ref{schemes-lemma-glue-functors}.

\medskip\noindent
Vérifions d'abord que $F_d$ est un faisceau pour la topologie de Zariski.
Soient $T$ un schéma, $T = \bigcup_{i \in I} U_i$ un recouvrement ouvert et
$(d, f_i, \mathcal{L}_i, \psi_i) \in F_d(U_i)$ tels que
$(d, f_i, \mathcal{L}_i, \psi_i)|_{U_i \cap U_j}$ et
$(d, f_j, \mathcal{L}_j, \psi_j)|_{U_i \cap U_j}$ soient strictement
équivalents. Les morphismes $f_i : U_i \to S$ se recollent alors en un
morphisme de schémas $f : T \to S$ tel que $f|_{U_i} = f_i$ ; voir
Schémas, Section \ref{schemes-section-glueing-schemes}.
On a donc $f_i^*\mathcal{A}^{(d)} = f^*\mathcal{A}^{(d)}|_{U_i}$.
Il existe aussi des isomorphismes
$\beta_{ij} : \mathcal{L}_i|_{U_i \cap U_j} \to \mathcal{L}_j|_{U_i \cap U_j}$
tels que $\beta_{ij} \circ \psi_i = \psi_j$ sur $U_i \cap U_j$.
Cette condition détermine les $\beta_{ij}$ de façon unique puisque les
morphismes $f_i^*\mathcal{A}_d \to \mathcal{L}_i$ sont surjectifs.
En particulier, $\beta_{jk} \circ \beta_{ij} = \beta_{ik}$ sur
$U_i \cap U_j \cap U_k$. Par Faisceaux,
Section \ref{sheaves-section-glueing-sheaves}, les faisceaux inversibles
$\mathcal{L}_i$ se recollent donc en un $\mathcal{O}_T$-module inversible
$\mathcal{L}$, et les morphismes $\psi_i$ se recollent en un morphisme de
$\mathcal{O}_T$-algèbres
$\psi : f^*\mathcal{A}^{(d)} \to \bigoplus_{n \geq 0} \mathcal{L}^{\otimes n}$.
Cela prouve que $F_d$ satisfait la condition de faisceau pour la topologie
de Zariski.

\medskip\noindent
Soit $S = \bigcup_{i \in I} U_i$ un recouvrement par des ouverts affines.
Soit $F_{d, i} \subset F_d$ le sous-foncteur constitué des quadruplets
$(d, f : T \to S, \mathcal{L}, \psi)$ tels que $f(T) \subset U_i$.

\medskip\noindent
Montrons que chaque $F_{d, i}$ est représentable. D'après le
Lemme \ref{lemma-proj-base-change}, ce foncteur s'identifie au foncteur
associé à $U_i$ muni de la $\mathcal{O}_{U_i}$-algèbre graduée
quasi-cohérente $\mathcal{A}|_{U_i}$. La représentabilité résulte donc du
Lemme \ref{lemma-relative-proj-affine}.

\medskip\noindent
Montrons ensuite que $F_{d, i} \subset F_d$ est représentable par des
immersions ouvertes. Soit
$(d, f : T \to S, \mathcal{L}, \psi) \in F_d(T)$ et posons
$V_i = f^{-1}(U_i)$. Par définition de $F_{d, i}$, pour un morphisme
$a : T' \to T$, on a
$a^*(d, f, \mathcal{L}, \psi) \in F_{d, i}(T')$ si et seulement si
$a(T') \subset V_i$. C'est ce qu'il fallait montrer.

\medskip\noindent
Enfin, montrons que la famille $(F_{d, i})_{i \in I}$ recouvre $F_d$.
Soit $(d, f : T \to S, \mathcal{L}, \psi) \in F_d(T)$ et posons
$V_i = f^{-1}(U_i)$. Puisque $S = \bigcup_{i \in I} U_i$ est un
recouvrement ouvert de $S$, on obtient un recouvrement ouvert
$T = \bigcup_{i \in I} V_i$ de $T$.
De plus, $(d, f, \mathcal{L}, \psi)|_{V_i} \in F_{d, i}(V_i)$.
Cela achève la démonstration du lemme.
\end{proof}

\noindent
Nous pouvons maintenant reprendre, dans le cadre relatif, la fin de la
Section \ref{section-morphisms-proj} et définir un foncteur représentable
par $\underline{\text{Proj}}_S(\mathcal{A})$.
Introduisons pour cela la notion d'équivalence entre deux quadruplets
$(d, f : T \to S, \mathcal{L}, \psi)$ et
$(d', f' : T \to S, \mathcal{L}', \psi')$, les entiers $d$ et $d'$
pouvant être distincts. Ils sont dits {\it équivalents} si $f = f'$ et
s'il existe un isomorphisme
$\beta : \mathcal{L}^{\otimes d'} \to (\mathcal{L}')^{\otimes d}$ tel que
$\beta \circ \psi|_{f^*\mathcal{A}^{(dd')}} = \psi'|_{f^*\mathcal{A}^{(dd')}}$.
Le lemme suivant montre que l'on obtient bien une relation d'équivalence
(ce qui n'est pas tout à fait immédiat).

\begin{lemma}
\label{lemma-equivalent-relative}
Dans la Situation \ref{situation-relative-proj}, soient $T$ un schéma et
$(d, f, \mathcal{L}, \psi)$, $(d', f', \mathcal{L}', \psi')$ deux
quadruplets sur $T$. Les assertions suivantes sont équivalentes :
\begin{enumerate}
\item Posons $m = \text{lcm}(d, d')$ et écrivons $m = ad = a'd'$.
On a $f = f'$ et il existe un isomorphisme
$\beta : \mathcal{L}^{\otimes a} \to (\mathcal{L}')^{\otimes a'}$
tel que $\beta \circ \psi|_{f^*\mathcal{A}^{(m)}}$ et
$\psi'|_{f^*\mathcal{A}^{(m)}}$ coïncident comme homomorphismes d'anneaux
gradués
$f^*\mathcal{A}^{(m)} \to \bigoplus_{n \geq 0} (\mathcal{L}')^{\otimes a'n}$.
\item Les quadruplets $(d, f, \mathcal{L}, \psi)$ et
$(d', f', \mathcal{L}', \psi')$ sont équivalents.
\item On a $f = f'$ et, pour un certain entier strictement positif
$m = ad = a'd'$, il existe un isomorphisme
$\beta : \mathcal{L}^{\otimes a} \to (\mathcal{L}')^{\otimes a'}$
tel que $\beta \circ \psi|_{f^*\mathcal{A}^{(m)}}$ et
$\psi'|_{f^*\mathcal{A}^{(m)}}$ coïncident comme homomorphismes d'anneaux
gradués
$f^*\mathcal{A}^{(m)} \to \bigoplus_{n \geq 0} (\mathcal{L}')^{\otimes a'n}$.
\end{enumerate}
\end{lemma}

\begin{proof}
Il est clair que (1) entraîne (2) et que (2) entraîne (3), en se restreignant
à des degrés plus divisibles et en prenant des puissances des faisceaux
inversibles. Supposons (3) pour un entier $m = ad = a'd'$.
Posons $m_0 = \text{lcm}(d, d')$ et écrivons $m_0 = a_0d = a'_0d'$.
Nous disposons d'un isomorphisme
$\beta : \mathcal{L}^{\otimes a} \to (\mathcal{L}')^{\otimes a'}$
ayant la propriété (3), et cherchons un isomorphisme
$\beta_0 : \mathcal{L}^{\otimes a_0} \to (\mathcal{L}')^{\otimes a'_0}$
ayant la même propriété. Les morphismes
$\psi : f^*\mathcal{A}_d \to \mathcal{L}$ et
$\psi' : (f')^*\mathcal{A}_{d'} \to \mathcal{L}'$ étant surjectifs par
hypothèse, il en est de même des morphismes
$\psi : f^*\mathcal{A}_{m_0} \to \mathcal{L}^{\otimes a_0}$ et
$\psi' : (f')^*\mathcal{A}_{m_0} \to (\mathcal{L}')^{\otimes a'_0}$.
Ainsi, s'il existe, $\beta_0$ est uniquement déterminé par la condition
$\beta_0 \circ \psi = \psi'$. On peut donc travailler localement sur $T$,
et supposer que $f = f' : T \to S$ est à valeurs dans un ouvert affine,
autrement dit que $S$ est affine. Le résultat découle alors de l'assertion
correspondante pour les triplets (Lemme \ref{lemma-equivalent}) et de la
correspondance entre triplets et quadruplets lorsque la base est affine
(voir la démonstration du Lemme \ref{lemma-relative-proj-affine}).
\end{proof}

\noindent
Supposons $d' = ad$. Considérons la transformation de foncteurs
$F_d \to F_{d'}$ qui associe au quadruplet $(d, f, \mathcal{L}, \psi)$ sur
$T$ le quadruplet
$(d', f, \mathcal{L}^{\otimes a}, \psi|_{f^*\mathcal{A}^{(d')}})$.
Une des implications du Lemme \ref{lemma-equivalent-relative} montre que
cette transformation est injective. Pour un schéma quasi-compact $T$,
définissons
$$
F(T) = \bigcup\nolimits_{d \in \mathbf{N}} F_d(T)
$$
avec les applications de transition ci-dessus. On obtient ainsi un foncteur
contravariant de la catégorie des schémas quasi-compacts vers celle des
ensembles. Pour un schéma quelconque $T$, définissons
$$
F(T)
=
\lim_{V \subset T\text{ ouvert quasi-compact}} F(V).
$$
Autrement dit, un élément $\xi$ de $F(T)$ correspond à une famille
compatible d'éléments $\xi_V \in F(V)$, où $V$ parcourt les ouverts
quasi-compacts de $T$. Nous omettons la définition de l'application d'image
inverse $F(T) \to F(T')$ associée à un morphisme de schémas $T' \to T$.
Nous avons ainsi défini le foncteur
\begin{equation}
\label{equation-proj}
F : \Sch^{opp} \longrightarrow \textit{Ens}
\end{equation}

\begin{lemma}
\label{lemma-relative-proj}
Dans la Situation \ref{situation-relative-proj}, le foncteur $F$ ci-dessus
est représentable par un schéma.
\end{lemma}

\begin{proof}
Soit $U_d \to S$ le schéma représentant le foncteur $F_d$ défini ci-dessus,
et soit $\mathcal{L}_d$,
$\psi^d : \pi_d^*\mathcal{A}^{(d)} \to
\bigoplus_{n \geq 0} \mathcal{L}_d^{\otimes n}$ l'objet universel.
Si $d | d'$, on peut considérer le quadruplet
$(d', \pi_d, \mathcal{L}_d^{\otimes d'/d}, \psi^d|_{\mathcal{A}^{(d')}})$,
qui détermine un morphisme canonique $U_d \to U_{d'}$ au-dessus de $S$.
Par construction, ce morphisme correspond à la transformation de foncteurs
$F_d \to F_{d'}$ définie ci-dessus.

\medskip\noindent
Pour tout ouvert affine $\Spec(R) = V \subset S$, en posant
$A = \Gamma(V, \mathcal{A})$, on dispose d'une identification canonique
du changement de base $U_{d, V}$ avec le sous-schéma ouvert correspondant
de $\text{Proj}(A)$ ; voir le Lemme \ref{lemma-relative-proj-affine}.
De plus, les morphismes $U_{d, V} \to U_{d', V}$ construits ci-dessus
correspondent aux inclusions d'ouverts dans $\text{Proj}(A)$.
On en déduit que $U_d \to U_{d'}$ est une immersion ouverte.

\medskip\noindent
On peut donc construire $X$ en recollant les schémas $U_d$ suivant les
immersions ouvertes $U_d \to U_{d'}$. Il est commode de choisir une suite
$d_1 | d_2 | d_3 | \ldots$ telle que tout entier strictement positif divise
l'un des $d_i$, puis de poser $X = \bigcup U_{d_i}$ au moyen des immersions
ouvertes ci-dessus. Il est alors aisé de montrer que $X$ représente $F$.
\end{proof}

\begin{lemma}
\label{lemma-glueing-gives-functor-proj}
Dans la Situation \ref{situation-relative-proj}, le schéma
$\pi : \underline{\text{Proj}}_S(\mathcal{A}) \to S$ construit dans le
Lemme \ref{lemma-glue-relative-proj} est canoniquement isomorphe, comme
schéma sur $S$, au schéma représentant le foncteur $F$.
\end{lemma}

\begin{proof}
Soit $X$ le schéma représentant $F$. C'est un schéma sur $S$, puisque $F$
est muni d'une transformation naturelle $F \to h_S$.
Posons $Y = \underline{\text{Proj}}_S(\mathcal{A})$. Nous devons montrer
que $X \cong Y$ comme $S$-schémas. Donnons deux démonstrations.

\medskip\noindent
La première utilise la construction de $X$ comme réunion des schémas $U_d$
représentant $F_d$ dans la démonstration du Lemme \ref{lemma-relative-proj}.
Au-dessus de chaque ouvert affine de $S$, on peut identifier $X$ au spectre
homogène des sections de $\mathcal{A}$ sur cet ouvert, puisque c'est le cas
pour les ouverts $U_d$. Ces identifications sont de plus compatibles avec
les restrictions à des ouverts affines plus petits. Or $Y$ a été construit
en recollant ces spectres homogènes. On peut donc recoller ces isomorphismes
en un isomorphisme entre $X$ et $\underline{\text{Proj}}_S(\mathcal{A})$.
Nous omettons les détails.

\medskip\noindent
Voici la seconde démonstration. Le
Lemme \ref{lemma-glue-relative-proj-twists} fournit un morphisme d'algèbres
graduées
$$
\psi : \pi^*\mathcal{A}
\longrightarrow
\bigoplus\nolimits_{n \geq 0} \mathcal{O}_Y(n)
$$
sur $Y$ qui, sur les sections au-dessus des ouverts affines de $S$, coïncide
avec (\ref{equation-global-sections}). Pour tout $y \in Y$, il existe donc
un voisinage ouvert $V \subset Y$ de $y$ et un entier $d \geq 1$ tels que,
pour $d | n$, le faisceau $\mathcal{O}_Y(n)|_V$ soit inversible et que les
morphismes de multiplication
$\mathcal{O}_Y(n)|_V \otimes_{\mathcal{O}_V} \mathcal{O}_Y(m)|_V
\to \mathcal{O}_Y(n + m)|_V$ soient des isomorphismes.
La restriction de $\psi$ à $\pi^*\mathcal{A}^{(d)}|_V$ donne ainsi un
élément de $F_d(V)$. Les ouverts $V$ recouvrant $Y$, on voit que « $\psi$ »
définit un élément de $F(Y)$, donc un morphisme canonique $Y \to X$ sur $S$.
La construction étant entièrement canonique, on peut vérifier localement
sur $S$ qu'il s'agit d'un isomorphisme. On se ramène donc au cas où $S$ est
affine, dans lequel le résultat est clair.
\end{proof}

\begin{definition}
\label{definition-relative-proj}
Soient $S$ un schéma et $\mathcal{A}$ un faisceau quasi-cohérent de
$\mathcal{O}_S$-algèbres graduées. Le {\it spectre homogène relatif de
$\mathcal{A}$ sur $S$}, ou {\it spectre homogène de $\mathcal{A}$ sur $S$},
ou encore {\it Proj relatif de $\mathcal{A}$ sur $S$}, est le schéma
construit dans le Lemme \ref{lemma-glue-relative-proj} et représentant le
foncteur $F$ de (\ref{equation-proj}) ; voir le
Lemme \ref{lemma-glueing-gives-functor-proj}.
On le note $\pi : \underline{\text{Proj}}_S(\mathcal{A}) \to S$.
\end{definition}

\noindent
Le Proj relatif est muni d'un faisceau quasi-cohérent d'algèbres
$\mathbf{Z}$-graduées
$\bigoplus_{n \in \mathbf{Z}}
\mathcal{O}_{\underline{\text{Proj}}_S(\mathcal{A})}(n)$
(les faisceaux tordus du faisceau structural) et d'un homomorphisme
« universel » d'algèbres graduées
$$
\psi_{univ} :
\mathcal{A}
\longrightarrow
\pi_*\left(
\bigoplus\nolimits_{n \geq 0}
\mathcal{O}_{\underline{\text{Proj}}_S(\mathcal{A})}(n)
\right)
$$
voir le Lemme \ref{lemma-glue-relative-proj-twists}. On peut également le
considérer comme un homomorphisme
$$
\psi_{univ} :
\pi^*\mathcal{A}
\longrightarrow
\bigoplus\nolimits_{n \geq 0}
\mathcal{O}_{\underline{\text{Proj}}_S(\mathcal{A})}(n)
$$
Le lemme suivant exprime l'universalité de cet objet.

\begin{lemma}
\label{lemma-tie-up-psi}
Dans la Situation \ref{situation-relative-proj}, soit
$(f : T \to S, d, \mathcal{L}, \psi)$ un quadruplet et soit
$r_{d, \mathcal{L}, \psi} : T \to \underline{\text{Proj}}_S(\mathcal{A})$
le $S$-morphisme associé. Il existe un isomorphisme de
$\mathcal{O}_T$-algèbres $\mathbf{Z}$-graduées
$$
\theta :
r_{d, \mathcal{L}, \psi}^*\left(
\bigoplus\nolimits_{n \in \mathbf{Z}}
\mathcal{O}_{\underline{\text{Proj}}_S(\mathcal{A})}(nd)
\right)
\longrightarrow
\bigoplus\nolimits_{n \in \mathbf{Z}} \mathcal{L}^{\otimes n}
$$
rendant commutatif le diagramme suivant
$$
\xymatrix{
\mathcal{A}^{(d)} \ar[rr]_-{\psi}
 \ar[rd]_-{\psi_{univ}} & &
f_*\left(
\bigoplus\nolimits_{n \in \mathbf{Z}}
\mathcal{L}^{\otimes n}
\right) \\
 &
\pi_*\left(
\bigoplus\nolimits_{n \geq 0}
\mathcal{O}_{\underline{\text{Proj}}_S(\mathcal{A})}(nd)
\right) \ar[ru]_\theta
}
$$
La commutativité de ce diagramme détermine $\theta$ de façon unique.
\end{lemma}

\begin{proof}
Le quadruplet $(f : T \to S, d, \mathcal{L}, \psi)$ définit un élément de
$F_d(T)$. Soit $U_d \subset \underline{\text{Proj}}_S(\mathcal{A})$ le lieu
où le faisceau $\mathcal{O}_{\underline{\text{Proj}}_S(\mathcal{A})}(d)$ est
inversible et engendré par l'image de
$\psi_{univ} : \pi^*\mathcal{A}_d \to
\mathcal{O}_{\underline{\text{Proj}}_S(\mathcal{A})}(d)$.
Rappelons que $U_d$ représente $F_d$ ; voir la démonstration du
Lemme \ref{lemma-relative-proj}. Il suffit donc de montrer que le quadruplet
$(U_d \to S, d, \mathcal{O}_{U_d}(d), \psi_{univ}|_{\mathcal{A}^{(d)}})$
est la famille universelle, c'est-à-dire l'objet représentant dans
$F_d(U_d)$. On peut le vérifier après restriction à un ouvert affine de
$S$, car (a) la formation des foncteurs $F_d$ commute au changement de base
(Lemme \ref{lemma-proj-base-change}) et (b) le couple
$(\bigoplus_{n \in \mathbf{Z}}
\mathcal{O}_{\underline{\text{Proj}}_S(\mathcal{A})}(n),
\psi_{univ})$
est construit par recollement au-dessus des ouverts affines de $S$
(Lemme \ref{lemma-glue-relative-proj-twists}).
On peut donc supposer $S$ affine. Alors le foncteur $F_d$ des quadruplets
coïncide avec le foncteur $F_d$ des triplets (voir la démonstration du
Lemme \ref{lemma-relative-proj-affine}), et le
Lemme \ref{lemma-proj-functor-strict} montre que
$(d, \mathcal{O}_{U_d}(d), \psi^d)$ est le triplet universel sur $U_d$.
En remontant les identifications de la démonstration du
Lemme \ref{lemma-relative-proj-affine}, on voit que
$(U_d \to S, d, \mathcal{O}_{U_d}(d), \psi_{univ}|_{\mathcal{A}^{(d)}})$
est bien le quadruplet universel cherché.
\end{proof}

\begin{lemma}
\label{lemma-relative-proj-separated}
Soient $S$ un schéma et $\mathcal{A}$ un faisceau quasi-cohérent de
$\mathcal{O}_S$-algèbres graduées. Le morphisme
$\pi : \underline{\text{Proj}}_S(\mathcal{A}) \to S$ est séparé.
\end{lemma}

\begin{proof}
Pour vérifier qu'un morphisme est séparé, on peut travailler localement sur
la base ; voir Schémas, Section \ref{schemes-section-separation-axioms}.
Par construction, $\underline{\text{Proj}}_S(\mathcal{A})$ est isomorphe,
au-dessus de chaque ouvert affine $U \subset S$, à $\text{Proj}(A)$, où
$A = \mathcal{A}(U)$. Par le Lemme \ref{lemma-proj-separated},
$\text{Proj}(A)$ est séparé. Le morphisme $\text{Proj}(A) \to U$ est donc
séparé (voir Schémas, Lemme \ref{schemes-lemma-compose-after-separated}),
comme voulu.
\end{proof}

\begin{lemma}
\label{lemma-relative-proj-base-change}
Soient $S$ un schéma et $\mathcal{A}$ un faisceau quasi-cohérent de
$\mathcal{O}_S$-algèbres graduées. Pour tout morphisme de schémas
$g : S' \to S$, il existe un isomorphisme canonique
$$
r :
\underline{\text{Proj}}_{S'}(g^*\mathcal{A})
\longrightarrow
S' \times_S \underline{\text{Proj}}_S(\mathcal{A})
$$
ainsi qu'un isomorphisme correspondant
$$
\theta :
r^*\text{pr}_2^*\left(\bigoplus\nolimits_{d \in \mathbf{Z}}
\mathcal{O}_{\underline{\text{Proj}}_S(\mathcal{A})}(d)\right)
\longrightarrow
\bigoplus\nolimits_{d \in \mathbf{Z}}
\mathcal{O}_{\underline{\text{Proj}}_{S'}(g^*\mathcal{A})}(d)
$$
de $\mathcal{O}_{\underline{\text{Proj}}_{S'}(g^*\mathcal{A})}$-algèbres
$\mathbf{Z}$-graduées.
\end{lemma}

\begin{proof}
Cela résulte du Lemme \ref{lemma-proj-base-change} et de la construction
de $\underline{\text{Proj}}_S(\mathcal{A})$, dans le
Lemme \ref{lemma-relative-proj}, comme réunion des schémas $U_d$ représentant
les foncteurs $F_d$. Du point de vue de la construction par recollement,
cet isomorphisme est donné par les isomorphismes du
Lemme \ref{lemma-base-change-map-proj}, qui fournit aussi $\theta$.
Nous omettons certains détails.
\end{proof}

\begin{lemma}
\label{lemma-apply-relative}
Soient $S$ un schéma et $\mathcal{A}$ un faisceau quasi-cohérent de
$\mathcal{O}_S$-modules gradués, engendré comme $\mathcal{A}_0$-algèbre par
$\mathcal{A}_1$. Alors $X = \underline{\text{Proj}}_S(\mathcal{A})$
représente le foncteur $F_1$ qui associe à un schéma $f : T \to S$ sur $S$
l'ensemble des couples $(\mathcal{L}, \psi)$ où
\begin{enumerate}
\item $\mathcal{L}$ est un $\mathcal{O}_T$-module inversible ;
\item $\psi : f^*\mathcal{A} \to \bigoplus_{n \geq 0} \mathcal{L}^{\otimes n}$
est un homomorphisme de $\mathcal{O}_T$-algèbres graduées tel que
$f^*\mathcal{A}_1 \to \mathcal{L}$ soit surjectif,
\end{enumerate}
à équivalence stricte près au sens ci-dessus. De plus, dans ce cas, tous les
faisceaux quasi-cohérents
$\mathcal{O}_{\underline{\text{Proj}}(\mathcal{A})}(n)$ sont des
$\mathcal{O}_{\underline{\text{Proj}}(\mathcal{A})}$-modules inversibles,
et les morphismes de multiplication induisent des isomorphismes
$
\mathcal{O}_{\underline{\text{Proj}}(\mathcal{A})}(n)
\otimes_{\mathcal{O}_{\underline{\text{Proj}}(\mathcal{A})}}
\mathcal{O}_{\underline{\text{Proj}}(\mathcal{A})}(m) =
\mathcal{O}_{\underline{\text{Proj}}(\mathcal{A})}(n + m)$.
\end{lemma}

\begin{proof}
Sous les hypothèses du lemme, les faisceaux
$\mathcal{O}_{\underline{\text{Proj}}(\mathcal{A})}(n)$ sont inversibles
et les morphismes de multiplication sont des isomorphismes, par le
Lemme \ref{lemma-relative-proj} et le Lemme \ref{lemma-apply} appliqués
au-dessus des ouverts affines de $S$. Ainsi, $X$ représente bien $F_1$ ;
voir la démonstration du Lemme \ref{lemma-relative-proj}.
\end{proof}











\section{Faisceaux quasi-cohérents sur le Proj relatif}
\label{section-quasi-coherent-relative-proj}

\noindent
Nous indiquons brièvement comment traiter les modules gradués dans le cadre
relatif.

\medskip\noindent
Plaçons-nous dans la Situation \ref{situation-relative-proj} : $S$ est un
schéma et $\mathcal{A}$ une $\mathcal{O}_S$-algèbre graduée quasi-cohérente.
Soit $\mathcal{M} = \bigoplus_{n \in \mathbf{Z}} \mathcal{M}_n$ un
$\mathcal{A}$-module gradué, quasi-cohérent comme $\mathcal{O}_S$-module.
Nous allons décrire le faisceau quasi-cohérent de modules associé sur
$\underline{\text{Proj}}_S(\mathcal{A})$. Nous décrivons d'abord son image
inverse sur les schémas $T$ munis d'un morphisme vers le Proj relatif.

\medskip\noindent
Soient $T$ un schéma et $(d, f : T \to S, \mathcal{L}, \psi)$ un quadruplet
sur $T$, au sens de la Section \ref{section-relative-proj}. Définissons le
faisceau quasi-cohérent de $\mathcal{O}_T$-modules
$\widetilde{\mathcal{M}}_T$ par
\begin{equation}
\label{equation-widetilde-M}
\widetilde{\mathcal{M}}_T =
\left(
f^*\mathcal{M}^{(d)}
\otimes_{f^*\mathcal{A}^{(d)}}
\left(\bigoplus\nolimits_{n \in \mathbf{Z}} \mathcal{L}^{\otimes n}\right)
\right)_0
\end{equation}
Ainsi, $\widetilde{\mathcal{M}}_T$ est la composante de degré $0$ du produit
tensoriel des $f^*\mathcal{A}^{(d)}$-modules gradués
$f^*\mathcal{M}^{(d)}$ et
$\bigoplus\nolimits_{n \in \mathbf{Z}} \mathcal{L}^{\otimes n}$.
Ce faisceau dépend du quadruplet, bien que la notation ne l'indique pas.
Cette construction a la propriété agréable suivante : pour tout morphisme
$g : T' \to T$, on a
$\widetilde{\mathcal{M}}_{T'} = g^*\widetilde{\mathcal{M}}_T$,
où $\widetilde{\mathcal{M}}_{T'}$ désigne le faisceau quasi-cohérent associé
au quadruplet image inverse $(d, f \circ g, g^*\mathcal{L}, g^*\psi)$.

\medskip\noindent
Tous les faisceaux de (\ref{equation-widetilde-M}) étant quasi-cohérents,
on peut expliciter la construction sur un ouvert affine
$\Spec(C) = V \subset T$ dont l'image est contenue dans un ouvert affine
$\Spec(R) = U \subset S$. Supposons que $\mathcal{A}|_U$ corresponde à la
$R$-algèbre graduée $A$, que $\mathcal{M}|_U$ corresponde au $A$-module
gradué $M$ et que $\mathcal{L}|_V$ corresponde au $C$-module inversible $L$.
Le morphisme $\psi$ induit un homomorphisme de $R$-algèbres graduées
$\gamma : A^{(d)} \to \bigoplus_{n \geq 0} L^{\otimes n}$
(les puissances tensorielles de $L$ sont prises sur $C$).
Alors $(\widetilde{\mathcal{M}}_T)|_V$ est le faisceau quasi-cohérent
associé au $C$-module
$$
N_{R, C, A, M, \gamma} =
\left(
M^{(d)} \otimes_{A^{(d)}, \gamma}
\left(\bigoplus\nolimits_{n \in \mathbf{Z}} L^{\otimes n}\right)
\right)_0
$$
Par hypothèse, on peut même recouvrir $T$ par des ouverts affines $V$ tels
qu'il existe $a \in A_d$ dont l'image $\gamma(a) \in L$ soit une base du
$C$-module $L$. Dans ce cas, tout élément de $N_{R, C, A, M, \gamma}$ est
une somme de tenseurs purs $\sum m_i \otimes \gamma(a)^{-n_i}$, avec
$m_i \in M_{n_id}$. En multipliant chaque $m_i$ par une puissance positive
convenable de $a$, puis en regroupant les termes, on voit que tout élément
de $N_{R, C, A, M, \gamma}$ s'écrit $m \otimes \gamma(a)^{-n}$ avec
$m \in M_{nd}$ et $n \gg 0$. Autrement dit, dans ce cas,
$$
N_{R, C, A, M, \gamma} = M_{(a)} \otimes_{A_{(a)}} C
$$
où $A_{(a)} \to C$ est donné par
$x/a^n \mapsto \gamma(x)/\gamma(a)^n$.
Il s'agit donc de l'image inverse sur $\Spec(C)$ de la restriction de
$\widetilde{M}$ à $D_{+}(a) \subset \text{Proj}(A)$, par le morphisme
$\Spec(C) \to D_{+}(a)$ défini par $\gamma$.

\begin{lemma}
\label{lemma-relative-proj-modules}
Dans la Situation \ref{situation-relative-proj}, à tout faisceau
quasi-cohérent de $\mathcal{A}$-modules gradués $\mathcal{M}$ sur $S$ est
canoniquement associé un faisceau
$\widetilde{\mathcal{M}}$ de
$\mathcal{O}_{\underline{\text{Proj}}_S(\mathcal{A})}$-modules ayant les
propriétés suivantes :
\begin{enumerate}
\item Pour tout schéma $T$ et tout quadruplet
$(T \to S, d, \mathcal{L}, \psi)$ sur $T$ correspondant à un morphisme
$h : T \to \underline{\text{Proj}}_S(\mathcal{A})$, il existe un
isomorphisme canonique
$\widetilde{\mathcal{M}}_T = h^*\widetilde{\mathcal{M}}$,
où $\widetilde{\mathcal{M}}_T$ est défini par (\ref{equation-widetilde-M}).
\item Les isomorphismes de (1) sont compatibles avec les images inverses.
\item Il existe un morphisme canonique
$$
\pi^*\mathcal{M}_0 \longrightarrow \widetilde{\mathcal{M}}.
$$
\item La construction $\mathcal{M} \mapsto \widetilde{\mathcal{M}}$ est
fonctorielle en $\mathcal{M}$.
\item La construction $\mathcal{M} \mapsto \widetilde{\mathcal{M}}$ est
exacte.
\item Il existe des morphismes canoniques
$$
\widetilde{\mathcal{M}}
\otimes_{\mathcal{O}_{\underline{\text{Proj}}_S(\mathcal{A})}}
\widetilde{\mathcal{N}}
\longrightarrow
\widetilde{\mathcal{M} \otimes_\mathcal{A} \mathcal{N}}
$$
comme dans le Lemme \ref{lemma-widetilde-tensor}.
\item Il existe des morphismes canoniques
$$
\pi^*\mathcal{M}
\longrightarrow
\bigoplus\nolimits_{n \in \mathbf{Z}}
\widetilde{\mathcal{M}(n)}
$$
généralisant (\ref{equation-global-sections-more-generally}).
\item La formation de $\widetilde{\mathcal{M}}$ commute au changement de base.
\end{enumerate}
\end{lemma}

\begin{proof}
La démonstration est omise. Il conviendrait de scinder ce lemme en plusieurs
énoncés et de les démontrer séparément.
\end{proof}


\section{Fonctorialité du Proj relatif}
\label{section-functoriality-relative-proj}

\noindent
Cette section est l'analogue de la Section \ref{section-functoriality-proj}
pour le Proj relatif. Soit $S$ un schéma. Un morphisme de
$\mathcal{O}_S$-algèbres graduées $\psi : \mathcal{A} \to \mathcal{B}$
n'induit pas toujours un morphisme entre les Proj relatifs associés.
L'énoncé correct est le suivant.

\begin{lemma}
\label{lemma-morphism-relative-proj}
Soient $S$ un schéma et $\mathcal{A}$, $\mathcal{B}$ deux
$\mathcal{O}_S$-algèbres graduées quasi-cohérentes. Posons
$p : X = \underline{\text{Proj}}_S(\mathcal{A}) \to S$ et
$q : Y = \underline{\text{Proj}}_S(\mathcal{B}) \to S$.
Soit $\psi : \mathcal{A} \to \mathcal{B}$ un homomorphisme de
$\mathcal{O}_S$-algèbres graduées. Il existe un ouvert canonique
$U(\psi) \subset Y$, un morphisme canonique de schémas sur $S$
$$
r_\psi :
U(\psi)
\longrightarrow
X
$$
et un morphisme de $\mathcal{O}_{U(\psi)}$-algèbres $\mathbf{Z}$-graduées
$$
\theta = \theta_\psi :
r_\psi^*\left(
\bigoplus\nolimits_{d \in \mathbf{Z}} \mathcal{O}_X(d)
\right)
\longrightarrow
\bigoplus\nolimits_{d \in \mathbf{Z}} \mathcal{O}_{U(\psi)}(d).
$$
Le triplet $(U(\psi), r_\psi, \theta)$ est caractérisé par la propriété
suivante : pour tout ouvert affine $W \subset S$, le triplet
$$
(U(\psi) \cap q^{-1}W,\quad
r_\psi|_{U(\psi) \cap q^{-1}W} : U(\psi) \cap q^{-1}W \to p^{-1}W,\quad
\theta|_{U(\psi) \cap q^{-1}W})
$$
est égal au triplet associé à
$\psi : \mathcal{A}(W) \to \mathcal{B}(W)$ dans le
Lemme \ref{lemma-morphism-proj}, via les identifications
$p^{-1}W = \text{Proj}(\mathcal{A}(W))$ et
$q^{-1}W = \text{Proj}(\mathcal{B}(W))$ de la
Section \ref{section-relative-proj-via-glueing}.
\end{lemma}

\begin{proof}
Le lemme s'obtient en recollant les triplets locaux.
\end{proof}

\begin{lemma}
\label{lemma-morphism-relative-proj-transitive}
Soient $S$ un schéma et $\mathcal{A}$, $\mathcal{B}$, $\mathcal{C}$ des
$\mathcal{O}_S$-algèbres graduées quasi-cohérentes. Posons
$X = \underline{\text{Proj}}_S(\mathcal{A})$,
$Y = \underline{\text{Proj}}_S(\mathcal{B})$ et
$Z = \underline{\text{Proj}}_S(\mathcal{C})$.
Soient $\varphi : \mathcal{A} \to \mathcal{B}$ et
$\psi : \mathcal{B} \to \mathcal{C}$ des morphismes de
$\mathcal{O}_S$-algèbres graduées. On a
$$
U(\psi \circ \varphi) = r_\psi^{-1}(U(\varphi))
\quad
\text{et}
\quad
r_{\psi \circ \varphi}
=
r_\varphi \circ r_\psi|_{U(\psi \circ \varphi)}.
$$
De plus,
$$
\theta_\psi \circ r_\psi^*\theta_\varphi
=
\theta_{\psi \circ \varphi}
$$
avec les notations évidentes.
\end{lemma}

\begin{proof}
La démonstration est omise.
\end{proof}

\begin{lemma}
\label{lemma-surjective-graded-rings-map-relative-proj}
Avec les hypothèses et notations du Lemme \ref{lemma-morphism-relative-proj},
supposons que $\mathcal{A}_d \to \mathcal{B}_d$ soit surjectif pour
$d \gg 0$. Alors
\begin{enumerate}
\item $U(\psi) = Y$ ;
\item $r_\psi : Y \to X$ est une immersion fermée ;
\item les morphismes $\theta : r_\psi^*\mathcal{O}_X(n) \to \mathcal{O}_Y(n)$
sont surjectifs, mais ne sont pas des isomorphismes en général, même si
$\mathcal{A} \to \mathcal{B}$ est surjectif.
\end{enumerate}
\end{lemma}

\begin{proof}
Cela résulte des Lemmes \ref{lemma-morphism-relative-proj} et
\ref{lemma-surjective-graded-rings-map-proj}.
\end{proof}

\begin{lemma}
\label{lemma-eventual-iso-graded-rings-map-relative-proj}
Avec les hypothèses et notations du Lemme \ref{lemma-morphism-relative-proj},
supposons que $\mathcal{A}_d \to \mathcal{B}_d$ soit un isomorphisme pour
tout $d \gg 0$. Alors
\begin{enumerate}
\item $U(\psi) = Y$ ;
\item $r_\psi : Y \to X$ est un isomorphisme ;
\item les morphismes $\theta : r_\psi^*\mathcal{O}_X(n) \to \mathcal{O}_Y(n)$
sont des isomorphismes.
\end{enumerate}
\end{lemma}

\begin{proof}
Cela résulte des Lemmes \ref{lemma-morphism-relative-proj} et
\ref{lemma-eventual-iso-graded-rings-map-proj}.
\end{proof}

\begin{lemma}
\label{lemma-surjective-generated-degree-1-map-relative-proj}
Avec les hypothèses et notations du Lemme \ref{lemma-morphism-relative-proj},
supposons que $\mathcal{A}_d \to \mathcal{B}_d$ soit surjectif pour
$d \gg 0$ et que $\mathcal{A}$ soit engendrée par $\mathcal{A}_1$ sur
$\mathcal{A}_0$. Alors
\begin{enumerate}
\item $U(\psi) = Y$ ;
\item $r_\psi : Y \to X$ est une immersion fermée ;
\item les morphismes $\theta : r_\psi^*\mathcal{O}_X(n) \to \mathcal{O}_Y(n)$
sont des isomorphismes.
\end{enumerate}
\end{lemma}

\begin{proof}
Cela résulte des Lemmes \ref{lemma-morphism-relative-proj} et
\ref{lemma-surjective-graded-rings-generated-degree-1-map-proj}.
\end{proof}


\section{Faisceaux inversibles et morphismes vers le Proj relatif}
\label{section-invertible-relative-proj}

\noindent
Le lemme suivant pourra nous être utile. Voici les données considérées :
\begin{enumerate}
\item $S$ est un schéma.
\item $\mathcal{A}$ est une $\mathcal{O}_S$-algèbre graduée quasi-cohérente.
\item On note $\pi : \underline{\text{Proj}}_S(\mathcal{A}) \to S$ le spectre
homogène relatif sur $S$.
\item $f : X \to S$ est un morphisme de schémas.
\item $\mathcal{L}$ est un $\mathcal{O}_X$-module inversible.
\item $\psi : f^*\mathcal{A} \to
\bigoplus_{d \geq 0} \mathcal{L}^{\otimes d}$ est un homomorphisme de
$\mathcal{O}_X$-algèbres graduées.
\end{enumerate}
Ces données étant fixées, posons
$$
U(\psi) = \bigcup\nolimits_{(U, V, a)} U_{\psi(a)}
$$
où $(U, V, a)$ vérifie les conditions suivantes :
\begin{enumerate}
\item $V \subset S$ est un ouvert affine ;
\item $U = f^{-1}(V)$ ;
\item $a \in \mathcal{A}(V)_{+}$ est homogène.
\end{enumerate}
On a alors $\psi(a) \in \Gamma(U, \mathcal{L}^{\otimes \deg(a)})$, et
$U_{\psi(a)}$ est l'ouvert correspondant (voir Modules,
Lemme \ref{modules-lemma-s-open}).

\begin{lemma}
\label{lemma-invertible-map-into-relative-proj}
Avec les hypothèses et notations ci-dessus, le morphisme $\psi$ induit un
morphisme canonique de schémas sur $S$
$$
r_{\mathcal{L}, \psi} :
U(\psi) \longrightarrow \underline{\text{Proj}}_S(\mathcal{A})
$$
muni d'un morphisme de $\mathcal{O}_{U(\psi)}$-algèbres graduées
$$
\theta :
r_{\mathcal{L}, \psi}^*\left(
\bigoplus\nolimits_{d \geq 0}
\mathcal{O}_{\underline{\text{Proj}}_S(\mathcal{A})}(d)
\right)
\longrightarrow
\bigoplus\nolimits_{d \geq 0} \mathcal{L}^{\otimes d}|_{U(\psi)}
$$
caractérisés par les propriétés suivantes :
\begin{enumerate}
\item Pour tout ouvert $V \subset S$ et tout $d \geq 0$, le diagramme
$$
\xymatrix{
\mathcal{A}_d(V) \ar[d]_{\psi} \ar[r]_{\psi} &
\Gamma(f^{-1}(V), \mathcal{L}^{\otimes d}) \ar[d]^{\text{restriction}} \\
\Gamma(\pi^{-1}(V),
\mathcal{O}_{\underline{\text{Proj}}_S(\mathcal{A})}(d)) \ar[r]^{\theta} &
\Gamma(f^{-1}(V) \cap U(\psi), \mathcal{L}^{\otimes d})
}
$$
est commutatif.
\item Pour tout $d \geq 1$ et tout sous-schéma ouvert $W \subset X$ tels que
$\psi|_W : f^*\mathcal{A}_d|_W \to \mathcal{L}^{\otimes d}|_W$ soit
surjectif, la restriction du morphisme $r_{\mathcal{L}, \psi}$ coïncide avec
le morphisme $W \to \underline{\text{Proj}}_S(\mathcal{A})$ fourni par la
construction du spectre homogène relatif ; voir la
Définition \ref{definition-relative-proj}.
\item Pour tout ouvert affine $V \subset S$, la restriction
$$
(U(\psi) \cap f^{-1}(V), r_{\mathcal{L}, \psi}|_{U(\psi) \cap f^{-1}(V)},
\theta|_{U(\psi) \cap f^{-1}(V)})
$$
coïncide, via $i_V$ (voir le Lemme \ref{lemma-glue-relative-proj}), avec le
triplet $(U(\psi'), r_{\mathcal{L}, \psi'}, \theta')$ du
Lemme \ref{lemma-invertible-map-into-proj} associé à l'homomorphisme
$\psi' : A = \mathcal{A}(V) \to \Gamma_*(f^{-1}(V), \mathcal{L}|_{f^{-1}(V)})$
induit par $\psi$.
\end{enumerate}
\end{lemma}

\begin{proof}
On utilise la caractérisation (3) pour construire $r_{\mathcal{L}, \psi}$
et $\theta$ localement sur $S$, puis l'unicité du
Lemme \ref{lemma-invertible-map-into-proj} pour montrer que ces constructions
se recollent. Nous omettons les détails.
\end{proof}


\section{Torsion par des faisceaux inversibles et Proj relatif}
\label{section-twisting-and-proj}

\noindent
Soient $S$ un schéma,
$\mathcal{A} = \bigoplus_{d \geq 0} \mathcal{A}_d$ une
$\mathcal{O}_S$-algèbre graduée quasi-cohérente et $\mathcal{L}$ un faisceau
inversible sur $S$. On obtient une autre $\mathcal{O}_S$-algèbre graduée
quasi-cohérente en posant
$$
\mathcal{B}
=
\bigoplus\nolimits_{d \geq 0}
\mathcal{A}_d \otimes_{\mathcal{O}_S} \mathcal{L}^{\otimes d}
$$
Les spectres homogènes relatifs de $\mathcal{A}$ et de $\mathcal{B}$ sont
isomorphes.

\begin{lemma}
\label{lemma-twisting-and-proj}
Avec les notations $S$, $\mathcal{A}$, $\mathcal{L}$ et $\mathcal{B}$
ci-dessus, il existe un isomorphisme canonique
$$
\xymatrix{
P = \underline{\text{Proj}}_S(\mathcal{A})
\ar[rr]_g \ar[rd]_\pi & &
\underline{\text{Proj}}_S(\mathcal{B}) = P'
\ar[ld]^{\pi'} \\
& S &
}
$$
ayant les propriétés suivantes :
\begin{enumerate}
\item Il existe des isomorphismes
$\theta_n : g^*\mathcal{O}_{P'}(n)
\to
\mathcal{O}_P(n) \otimes \pi^*\mathcal{L}^{\otimes n}$
qui forment un isomorphisme d'algèbres $\mathbf{Z}$-graduées
$$
\theta :
g^*\left(
\bigoplus\nolimits_{n \in \mathbf{Z}} \mathcal{O}_{P'}(n)
\right)
\longrightarrow
\bigoplus\nolimits_{n \in \mathbf{Z}} \mathcal{O}_P(n)
\otimes \pi^*\mathcal{L}^{\otimes n}
$$
\item Pour tout ouvert $V \subset S$, les diagrammes
$$
\xymatrix{
\mathcal{A}_n(V) \otimes \mathcal{L}^{\otimes n}(V)
\ar[r]_{\text{multiplication}} \ar[d]^{\psi \otimes \pi^*}
&
\mathcal{B}_n(V) \ar[dd]^\psi \\
\Gamma(\pi^{-1}V, \mathcal{O}_P(n)) \otimes
\Gamma(\pi^{-1}V, \pi^*\mathcal{L}^{\otimes n})
\ar[d]^{\text{multiplication}} \\
\Gamma(\pi^{-1}V, \mathcal{O}_P(n) \otimes \pi^*\mathcal{L}^{\otimes n})
&
\Gamma(\pi'^{-1}V, \mathcal{O}_{P'}(n)) \ar[l]_-{\theta_n}
}
$$
sont commutatifs.
\item Ajouter d'autres propriétés ici au besoin.
\end{enumerate}
\end{lemma}

\begin{proof}
Lorsque $\mathcal{L} \cong \mathcal{O}_S$, il s'agit de l'identité.
Dans le cas général, on choisit un recouvrement ouvert de $S$ sur les
éléments duquel $\mathcal{L}$ se trivialise, puis on recolle les morphismes
correspondants. Nous omettons les détails.
\end{proof}


\section{Fibrés projectifs}
\label{section-projective-bundle}

\noindent
Soient $S$ un schéma et $\mathcal{E}$ un faisceau quasi-cohérent de
$\mathcal{O}_S$-modules. Par Modules,
Lemme \ref{modules-lemma-whole-tensor-algebra-permanence}, l'algèbre
symétrique $\text{Sym}(\mathcal{E})$ de $\mathcal{E}$ sur $\mathcal{O}_S$
est un faisceau quasi-cohérent de $\mathcal{O}_S$-algèbres. Elle est
engendrée en degré $1$ sur $\mathcal{O}_S$. On peut donc lui appliquer la
construction de la section précédente, notamment les Lemmes
\ref{lemma-relative-proj} et \ref{lemma-apply-relative}.

\begin{definition}
\label{definition-projective-bundle}
Soient $S$ un schéma et $\mathcal{E}$ un $\mathcal{O}_S$-module
quasi-cohérent\footnote{Le lecteur pourrait s'attendre à la condition que
$\mathcal{E}$ soit localement libre de rang fini. Nous ne l'imposons pas,
afin de rester cohérents avec \cite[II, Definition 4.1.1]{EGA}.}.
On note
$$
\pi :
\mathbf{P}(\mathcal{E}) = \underline{\text{Proj}}_S(\text{Sym}(\mathcal{E}))
\longrightarrow
S
$$
et on l'appelle le {\it fibré projectif associé à $\mathcal{E}$}.
Le symbole $\mathcal{O}_{\mathbf{P}(\mathcal{E})}(n)$ désigne le
$\mathcal{O}_{\mathbf{P}(\mathcal{E})}$-module inversible du
Lemme \ref{lemma-apply-relative}, appelé le {\it faisceau structural tordu}
d'indice $n$.
\end{definition}

\noindent
D'après le Lemme \ref{lemma-glue-relative-proj-twists}, il existe des
homomorphismes canoniques de $\mathcal{O}_S$-modules
$$
\text{Sym}^n(\mathcal{E})
\longrightarrow
\pi_*\mathcal{O}_{\mathbf{P}(\mathcal{E})}(n)
\quad\text{ou, de façon équivalente,}\quad
\pi^*\text{Sym}^n(\mathcal{E})
\longrightarrow
\mathcal{O}_{\mathbf{P}(\mathcal{E})}(n)
$$
pour tout $n \geq 0$. En particulier, pour $n = 1$, on a
$$
\mathcal{E}
\longrightarrow
\pi_*\mathcal{O}_{\mathbf{P}(\mathcal{E})}(1)
\quad\text{ou, de façon équivalente,}\quad
\pi^*\mathcal{E}
\longrightarrow
\mathcal{O}_{\mathbf{P}(\mathcal{E})}(1)
$$
et $\pi^*\mathcal{E} \to \mathcal{O}_{\mathbf{P}(\mathcal{E})}(1)$ est
surjectif par le Lemme \ref{lemma-apply-relative}. C'est une bonne façon
de retenir la convention adoptée pour construire $\mathbf{P}(\mathcal{E})$.

\medskip\noindent
Attention : certains ouvrages ne définissent $\mathbf{P}(\mathcal{E})$
que lorsque $\mathcal{E}$ est localement libre de rang fini sur $S$.
De plus, ils notent parfois $\mathbf{P}(\mathcal{E})$ ce que nous notons
$\mathbf{P}(\mathcal{E}^\vee)$, où $\mathcal{E}^\vee$ est le dual de
$\mathcal{E}$ (là encore, seulement dans le cas localement libre de rang
fini).

\medskip\noindent
Soient $S$, $\mathcal{E}$ et $\mathbf{P}(\mathcal{E}) \to S$ comme dans la
Définition \ref{definition-projective-bundle}. Soit $f : T \to S$ un
schéma sur $S$ et soit $\psi : f^*\mathcal{E} \to \mathcal{L}$ une
surjection, où $\mathcal{L}$ est un $\mathcal{O}_T$-module inversible.
Le morphisme induit de $\mathcal{O}_T$-algèbres graduées
$$
f^*\text{Sym}(\mathcal{E}) = \text{Sym}(f^*\mathcal{E}) \to
\text{Sym}(\mathcal{L}) = \bigoplus\nolimits_{n \geq 0} \mathcal{L}^{\otimes n}
$$
correspond à un morphisme
$$
\varphi_{\mathcal{L}, \psi} : T \longrightarrow \mathbf{P}(\mathcal{E})
$$
sur $S$, d'après la construction du Proj relatif comme schéma représentant
le foncteur $F$ dans la Section \ref{section-relative-proj}.
Réciproquement, pour un morphisme $\varphi : T \to \mathbf{P}(\mathcal{E})$
sur $S$, on pose
$\mathcal{L} = \varphi^*\mathcal{O}_{\mathbf{P}(\mathcal{E})}(1)$ et on
prend pour $\psi : f^*\mathcal{E} \to \mathcal{L}$ l'image inverse par
$\varphi$ de la surjection canonique
$\pi^*\mathcal{E} \to \mathcal{O}_{\mathbf{P}(\mathcal{E})}(1)$.
Par le Lemme \ref{lemma-apply-relative}, ces constructions sont des
bijections réciproques entre les classes d'isomorphisme de couples
$(\mathcal{L}, \psi)$ et les morphismes
$\varphi : T \to \mathbf{P}(\mathcal{E})$ sur $S$.
Ainsi, $\mathbf{P}(\mathcal{E})$ représente le foncteur qui associe à
$f : T \to S$ l'ensemble des quotients de $f^*\mathcal{E}$ qui sont des
$\mathcal{O}_T$-modules localement libres de rang $1$.

\begin{example}[Espace projectif d'un espace vectoriel]
\label{example-projective-space}
Soient $k$ un corps et $V$ un $k$-espace vectoriel. L'{\it espace projectif}
correspondant est le $k$-schéma
$$
\mathbf{P}(V) = \text{Proj}(\text{Sym}(V))
$$
où $\text{Sym}(V)$ est l'algèbre symétrique de $V$ sur $k$.
On a bien sûr $\mathbf{P}(V) \cong \mathbf{P}^n_k$ si $\dim(V) = n + 1$,
puisque l'algèbre symétrique de $V$ est alors isomorphe à un anneau de
polynômes en $n + 1$ variables. Si l'on considère $V$ comme un module
quasi-cohérent sur $\Spec(k)$, alors $\mathbf{P}(V)$ est le fibré projectif
correspondant sur $\Spec(k)$. D'après ce qui précède, un point $p$ de
$\mathbf{P}(V)$ à valeurs dans $k$ correspond à une surjection de
$k$-espaces vectoriels $V \to L_p$, avec $\dim(L_p) = 1$.
Plus généralement, soient $X$ un schéma sur $k$, $\mathcal{L}$ un
$\mathcal{O}_X$-module inversible et
$\psi : V \to \Gamma(X, \mathcal{L})$ une application $k$-linéaire telle
que $\mathcal{L}$ soit engendré, comme $\mathcal{O}_X$-module, par les
sections appartenant à l'image de $\psi$. On obtient un morphisme canonique
$$
\varphi_{\mathcal{L}, \psi} : X \longrightarrow \mathbf{P}(V)
$$
de schémas sur $k$, muni d'un isomorphisme
$\theta : \varphi_{\mathcal{L}, \psi}^*\mathcal{O}_{\mathbf{P}(V)}(1)
\to \mathcal{L}$ tel que $\psi$ coïncide avec la composée
$$
V \to
\Gamma(\mathbf{P}(V), \mathcal{O}_{\mathbf{P}(V)}(1))
\to
\Gamma(X, \varphi_{\mathcal{L}, \psi}^*\mathcal{O}_{\mathbf{P}(V)}(1))
\to
\Gamma(X, \mathcal{L})
$$
Voir le Lemme \ref{lemma-invertible-map-into-proj}.
Si $V \subset \Gamma(X, \mathcal{L})$ est un sous-espace, nous noterons
simplement $\varphi_{\mathcal{L}, V}$ le morphisme ainsi construit.
Si $\dim(V) = n + 1$ et si l'on choisit une base $v_0, \ldots, v_n$ de $V$,
le diagramme
$$
\xymatrix{
X \ar@{=}[d] \ar[rr]_{\varphi_{\mathcal{L}, \psi}} & &
\mathbf{P}(V) \ar[d]^{\cong} \\
X \ar[rr]^{\varphi_{(\mathcal{L}, (s_0, \ldots, s_n))}} & &
\mathbf{P}^n_k
}
$$
est commutatif, où $s_i = \psi(v_i) \in \Gamma(X, \mathcal{L})$,
où $\varphi_{(\mathcal{L}, (s_0, \ldots, s_n))}$ est défini comme dans la
Section \ref{section-projective-space}, et où la flèche verticale de
droite correspond à l'isomorphisme
$k[T_0, \ldots, T_n] \to \text{Sym}(V)$ envoyant $T_i$ sur $v_i$.
\end{example}

\begin{example}
\label{example-projective-bundle}
Le morphisme $\text{Sym}^n(\mathcal{E}) \to
\pi_*(\mathcal{O}_{\mathbf{P}(\mathcal{E})}(n))$ est un isomorphisme lorsque
$\mathcal{E}$ est localement libre, mais pas nécessairement dans le cas
général. Nous allons donner un exemple où il n'est pas injectif pour
$n = 1$. Posons $S = \Spec(A)$, avec
$$
A = k[u, v, s_1, s_2, t_1, t_2]/I
$$
où $k$ est un corps et
$$
I = (-us_1 + vt_1 + ut_2, vs_1 + us_2 - vt_2, vs_2, ut_1).
$$
Notons $\overline{u}$ la classe de $u$ dans $A$, et de même pour les autres
variables. Soit $M = (Ax \oplus Ay)/A(\overline{u}x + \overline{v}y)$,
de sorte que
$$
\text{Sym}(M) = A[x, y]/(\overline{u}x + \overline{v}y)
= k[x, y, u, v, s_1, s_2, t_1, t_2]/J
$$
où
$$
J = (-us_1 + vt_1 + ut_2, vs_1 + us_2 - vt_2, vs_2, ut_1, ux + vy).
$$
Le fibré projectif associé au faisceau quasi-cohérent
$\mathcal{E} = \widetilde{M}$ sur $S = \Spec(A)$ est alors le schéma
$$
P =
\text{Proj}(\text{Sym}(M)).
$$
Il admet le recouvrement ouvert affine $P = D_{+}(x) \cup D_{+}(y)$.
Considérons l'élément $m \in M$ qui est l'image de $us_1x + vt_2y$.
On a
$$
x(us_1x + vt_2y) = (s_1x + s_2y)(ux + vy) \bmod I
$$
et
$$
y(us_1x + vt_2y) = (t_1x + t_2y)(ux + vy) \bmod I.
$$
La première égalité montre que $m$ s'envoie sur la section nulle de
$\mathcal{O}_P(1)$ sur $D_{+}(x)$ ; la seconde donne la même conclusion
sur $D_{+}(y)$. Ainsi, l'image de $m$ dans
$\Gamma(P, \mathcal{O}_P(1))$ est nulle.
Montrons pourtant que $m \not = 0$ : on aura ainsi une section globale
non nulle de $\mathcal{E}$ dont l'image dans
$\Gamma(P, \mathcal{O}_P(1))$ est nulle. Supposons $m = 0$ pour obtenir
une contradiction. Il existe alors
$f \in k[u, v, s_1, s_2, t_1, t_2]$ tel que
$$
us_1x + vt_2y = f(ux + vy) \bmod I
$$
Puisque $I$ est engendré par des polynômes homogènes de degré $2$, on peut
décomposer $f$ en composantes homogènes et ne conserver que celle de degré
1. Autrement dit, on peut supposer
$$
f = au + bv + \alpha_1s_1 + \alpha_2s_2 + \beta_1t_1 + \beta_2t_2
$$
avec $a, b, \alpha_1, \alpha_2, \beta_1, \beta_2 \in k$.
On obtient les conditions
$$
\begin{matrix}
us_1 - u(au + bv + \alpha_1s_1 + \alpha_2s_2 + \beta_1t_1 + \beta_2t_2)
\in I \\
vt_2 - v(au + bv + \alpha_1s_1 + \alpha_2s_2 + \beta_1t_1 + \beta_2t_2)
\in I
\end{matrix}
$$
Les générateurs de $I$ ne contiennent aucun terme en $u^2, uv, v^2$ ;
on a donc $a = b = 0$. Il reste les relations
$$
\begin{matrix}
us_1 - u(\alpha_1s_1 + \alpha_2s_2 + \beta_1t_1 + \beta_2t_2)
\in I \\
vt_2 - v(\alpha_1s_1 + \alpha_2s_2 + \beta_1t_1 + \beta_2t_2)
\in I
\end{matrix}
$$
Le premier générateur de $I$ permet de remplacer $us_1$ par $vt_1 + ut_2$,
le deuxième de remplacer $vs_1$ par $-us_2 + vt_2$, le troisième d'éliminer
$vs_2$, et le quatrième d'éliminer $ut_1$. On obtient
$$
\begin{matrix}
(1 - \alpha_1)vt_1 + (1 - \alpha_1)ut_2 - \alpha_2us_2 - \beta_2ut_2 = 0 \\
(1 - \alpha_1)vt_2 + \alpha_1us_2 - \beta_1vt_1 - \beta_2vt_2 = 0
\end{matrix}
$$
Ces relations imposent à $\alpha_1$ d'être à la fois égal à $0$ et à $1$,
ce qui donne la contradiction cherchée.
\end{example}


\begin{lemma}
\label{lemma-projective-bundle-separated}
Soit $S$ un schéma. Le morphisme structural
$\mathbf{P}(\mathcal{E}) \to S$ d'un fibré projectif sur $S$ est séparé.
\end{lemma}

\begin{proof}
Cela résulte immédiatement du Lemme \ref{lemma-relative-proj-separated}.
\end{proof}

\begin{lemma}
\label{lemma-projective-space-bundle}
Soient $S$ un schéma et $n \geq 0$. Alors $\mathbf{P}^n_S$ est un fibré
projectif sur $S$.
\end{lemma}

\begin{proof}
On a
$$
\mathbf{P}^n_{\mathbf{Z}} =
\text{Proj}(\mathbf{Z}[T_0, \ldots, T_n]) =
\underline{\text{Proj}}_{\Spec(\mathbf{Z})}
\left(\widetilde{\mathbf{Z}[T_0, \ldots, T_n]}\right)
$$
où l'anneau $\mathbf{Z}[T_0, \ldots, T_n]$ est gradué par
$\deg(T_i) = 1$, les éléments de $\mathbf{Z}$ étant de degré $0$.
Rappelons que $\mathbf{P}^n_S$ est défini comme
$\mathbf{P}^n_{\mathbf{Z}} \times_{\Spec(\mathbf{Z})} S$.
De plus, la formation du spectre homogène relatif commute au changement de
base ; voir le Lemme \ref{lemma-relative-proj-base-change}.
Pour tout schéma $g : S \to \Spec(\mathbf{Z})$, on a
$g^*\mathcal{O}_{\Spec(\mathbf{Z})}[T_0, \ldots, T_n]
= \mathcal{O}_S[T_0, \ldots, T_n]$.
En combinant ces observations, on obtient
$$
\mathbf{P}^n_S = \underline{\text{Proj}}_S(\mathcal{O}_S[T_0, \ldots, T_n]).
$$
Enfin,
$\mathcal{O}_S[T_0, \ldots, T_n] = \text{Sym}(\mathcal{O}_S^{\oplus n + 1})$.
Ainsi, $\mathbf{P}^n_S$ est un fibré projectif sur $S$.
\end{proof}


\section{Grassmanniennes}
\label{section-grassmannian}

\noindent
Dans cette section, nous introduisons les foncteurs grassmanniens usuels et
montrons qu'ils sont représentables par des schémas. Fixons des entiers
$k$, $n$ tels que $0 < k < n$. Nous allons construire un foncteur
\begin{equation}
\label{equation-gkn}
G(k, n) : \Sch^{opp} \longrightarrow \textit{Ens}
\end{equation}
qui, en termes approximatifs, paramétrera les sous-espaces de dimension $k$
d'un espace de dimension $n$. Pour des raisons techniques, il est cependant
plus commode de paramétrer les quotients de dimension $n - k$ ; c'est la
convention que nous adopterons.

\medskip\noindent
Plus précisément, $G(k, n)$ associe à un schéma $S$ l'ensemble
$G(k, n)(S)$ des classes d'isomorphisme de surjections
$$
q : \mathcal{O}_S^{\oplus n} \longrightarrow \mathcal{Q}
$$
où $\mathcal{Q}$ est un $\mathcal{O}_S$-module localement libre de rang
fini $n - k$. Il s'agit bien d'un ensemble : cela résulte, par exemple,
de Modules,
Lemme \ref{modules-lemma-set-isomorphism-classes-finite-type-modules},
ou du fait que la classe d'isomorphisme de $q$ est déterminée par son noyau
(et que les sous-faisceaux d'un faisceau donné forment un ensemble).
À un morphisme de schémas $f : T \to S$, on associe l'application
$G(k, n)(f) : G(k, n)(S) \to G(k, n)(T)$ envoyant la classe de
$q : \mathcal{O}_S^{\oplus n} \longrightarrow \mathcal{Q}$ sur celle de
$f^*q : \mathcal{O}_T^{\oplus n} \longrightarrow f^*\mathcal{Q}$.
Cette définition a un sens, car (1) $f^*\mathcal{O}_S = \mathcal{O}_T$,
(2) $f^*$ est additif, (3) $f^*$ préserve les modules localement libres
(Modules, Lemme \ref{modules-lemma-pullback-locally-free}) et (4) $f^*$
est exact à droite (Modules,
Lemme \ref{modules-lemma-exactness-pushforward-pullback}).

\begin{lemma}
\label{lemma-gkn-representable}
Soit $0 < k < n$. Le foncteur $G(k, n)$ de (\ref{equation-gkn}) est
représentable par un schéma.
\end{lemma}

\begin{proof}
Posons $F = G(k, n)$ et appliquons le critère de Schémas,
Lemme \ref{schemes-lemma-glue-functors}.
Le foncteur $F$ est un faisceau pour la topologie de Zariski parce que
l'on peut recoller les faisceaux ; voir Faisceaux,
Section \ref{sheaves-section-glueing-sheaves} (certains détails sont omis).

\medskip\noindent
La famille de sous-foncteurs $F_i$.
Soit $I$ l'ensemble des parties de $\{1, \ldots, n\}$ de cardinal $n - k$.
Pour un schéma $S$ et $j \in \{1, \ldots, n\}$, notons $e_j$ la section
globale
$$
e_j = (0, \ldots, 0, 1, 0, \ldots, 0)\quad(1\text{ à la }j\text{-ième place})
$$
de $\mathcal{O}_S^{\oplus n}$. Ces sections forment une base de
$\mathcal{O}_S^{\oplus n}$. De même, pour $j \in \{1, \ldots, n - k\}$,
notons $f_j$ la section globale de $\mathcal{O}_S^{\oplus n - k}$ dont
toutes les composantes sont nulles, sauf la $j$-ième qui vaut $1$.
Pour $i \in I$, soit
$$
s_i : \mathcal{O}_S^{\oplus n - k} \longrightarrow \mathcal{O}_S^{\oplus n}
$$
la somme directe des injections canoniques
$\mathcal{O}_S \to \mathcal{O}_S^{\oplus n}$ correspondant aux éléments
de $i$. Plus précisément, si $i = \{i_1, \ldots, i_{n - k}\}$ avec
$i_1 < i_2 < \ldots < i_{n - k}$, le morphisme $s_i$ envoie $f_j$ sur
$e_{i_j}$ pour $j \in \{1, \ldots, n - k\}$. On pose
$$
F_i(S) = \{q : \mathcal{O}_S^{\oplus n} \to \mathcal{Q} \in F(S) \mid
q \circ s_i \text{ est surjectif}\}
\subset F(S)
$$
Pour un morphisme de schémas $f : T \to S$, l'image inverse $f^*s_i$ est
le morphisme correspondant sur $T$. Puisque $f^*$ est exact à droite
(Modules, Lemme \ref{modules-lemma-exactness-pushforward-pullback}),
on en déduit que $F_i$ est un sous-foncteur de $F$.

\medskip\noindent
Représentabilité de $F_i$. Après renumérotation, on peut supposer
$i = \{1, \ldots, n - k\}$ ; ainsi, $s_i$ est l'inclusion des
$n - k$ premiers facteurs directs. Si $q \circ s_i$ est surjectif, c'est
un isomorphisme, puisqu'il s'agit d'un morphisme surjectif entre modules
localement libres de même rang fini (Modules,
Lemme \ref{modules-lemma-map-finite-locally-free}).
Pour $q : \mathcal{O}_S^{\oplus n} \to \mathcal{Q}$ dans $F_i(S)$, on
peut donc identifier $\mathcal{Q}$ à $\mathcal{O}_S^{\oplus n - k}$ à
l'aide de $q \circ s_i$. On obtient ainsi
$$
q : \mathcal{O}_S^{\oplus n} \longrightarrow \mathcal{O}_S^{\oplus n - k}
$$
qui envoie $e_j$ sur $f_j$ pour $j = 1, \ldots, n - k$, avec les notations
ci-dessus. Pour déterminer entièrement $q$, il reste à choisir les images
$q(e_{n - k + 1}), \ldots, q(e_n)$ dans
$\Gamma(S, \mathcal{O}_S^{\oplus n - k})$.
Ainsi, $F_i$ est isomorphe au foncteur
$$
S \longmapsto
\prod\nolimits_{j = n - k + 1, \ldots, n}
\Gamma(S,  \mathcal{O}_S^{\oplus n - k})
$$
qui est lui-même isomorphe au produit de $k(n - k)$ copies du foncteur
$S \mapsto \Gamma(S, \mathcal{O}_S)$. Ce dernier est représentable par
$\mathbf{A}^1_\mathbf{Z}$, d'après Schémas,
Exemple \ref{schemes-example-global-sections}. Par conséquent, $F_i$ est
représentable par $\mathbf{A}^{k(n - k)}_\mathbf{Z}$, puisque le produit
fibré sur $\Spec(\mathbf{Z})$ est le produit dans la catégorie des schémas.

\medskip\noindent
L'inclusion $F_i \subset F$ est représentable par des immersions ouvertes.
Soient $S$ un schéma et $q : \mathcal{O}_S^{\oplus n} \to \mathcal{Q}$
un élément de $F(S)$. Par Modules,
Lemme \ref{modules-lemma-finite-type-surjective-on-stalk}, l'ensemble
$U_i = \{s \in S \mid (q \circ s_i)_s\text{ surjectif}\}$ est ouvert
dans $S$. Puisque $\mathcal{O}_{S, s}$ est un anneau local et
$\mathcal{Q}_s$ un $\mathcal{O}_{S, s}$-module de type fini, le lemme de
Nakayama (Algèbre, Lemme \ref{algebra-lemma-NAK}) donne
$$
s \in U_i \Leftrightarrow
\left(
\text{l'application }
\kappa(s)^{\oplus n - k} \to \mathcal{Q}_s/\mathfrak m_s\mathcal{Q}_s
\text{ induite par }
(q \circ s_i)_s
\text{ est surjective}
\right)
$$
Soient $f : T \to S$ un morphisme de schémas et $t \in T$ un point
d'image $s \in S$. On a
$(f^*\mathcal{Q})_t =
\mathcal{Q}_s \otimes_{\mathcal{O}_{S, s}} \mathcal{O}_{T, t}$
(Faisceaux, Lemme \ref{sheaves-lemma-stalk-pullback-modules}), et de même
pour les morphismes correspondants. Ainsi, l'application
$$
\kappa(t)^{\oplus n - k} \to (f^*\mathcal{Q})_t/\mathfrak m_t(f^*\mathcal{Q})_t
$$
induite par $(f^*q \circ f^*s_i)_t$ s'obtient par changement de base de
l'application
$\kappa(s)^{\oplus n - k} \to \mathcal{Q}_s/\mathfrak m_s\mathcal{Q}_s$
ci-dessus suivant l'extension de corps $\kappa(t)/\kappa(s)$.
Il en résulte que $s \in U_i$ si et seulement si $t$ appartient à l'ouvert
correspondant pour $f^*q$. En particulier, $T \to S$ se factorise par
$U_i$ si et seulement si $f^*q \in F_i(T)$, comme voulu.

\medskip\noindent
La famille des $F_i$, $i \in I$, recouvre $F$.
Soit $q : \mathcal{O}_S^{\oplus n} \to \mathcal{Q}$ un élément de $F(S)$.
Il faut montrer que, pour tout point $s$ de $S$, il existe $i \in I$ tel
que $q \circ s_i$ soit surjectif dans un voisinage de $s$.
Il suffit donc de montrer que l'une des composées
$$
\kappa(s)^{\oplus n - k} \xrightarrow{s_i}
\kappa(s)^{\oplus n} \rightarrow
\mathcal{Q}_s/\mathfrak m_s\mathcal{Q}_s
$$
est surjective (voir le paragraphe précédent). Comme
$\mathcal{Q}_s/\mathfrak m_s\mathcal{Q}_s$ est un espace vectoriel de
dimension $n - k$, cela résulte de l'algèbre linéaire.
\end{proof}

\begin{definition}
\label{definition-grassmannian}
Soit $0 < k < n$. Le schéma $\mathbf{G}(k, n)$ représentant le foncteur
$G(k, n)$ est appelé la {\it grassmannienne sur $\mathbf{Z}$}.
Son changement de base $\mathbf{G}(k, n)_S$ à un schéma $S$ est appelé la
{\it grassmannienne sur $S$}. Si $R$ est un anneau, le changement de base
à $\Spec(R)$ est noté $\mathbf{G}(k, n)_R$ et appelé la
{\it grassmannienne sur $R$}.
\end{definition}

\noindent
Cette définition a un sens, puisque la représentabilité de ces foncteurs a
été démontrée dans le Lemme \ref{lemma-gkn-representable}.

\begin{lemma}
\label{lemma-projective-space-grassmannian}
Soit $n \geq 1$. Il existe un isomorphisme canonique
$\mathbf{G}(n, n + 1) = \mathbf{P}^n_\mathbf{Z}$.
\end{lemma}

\begin{proof}
D'après le Lemme \ref{lemma-projective-space}, le schéma
$\mathbf{P}^n_\mathbf{Z}$ représente le foncteur qui associe à un schéma
$S$ l'ensemble des classes d'isomorphisme de couples
$(\mathcal{L}, (s_0, \ldots, s_n))$ formés d'un module inversible
$\mathcal{L}$ et de $n + 1$ sections globales qui l'engendrent.
Un tel couple définit un quotient
$$
\mathcal{O}_S^{\oplus n + 1} \longrightarrow \mathcal{L},\quad
(h_0, \ldots, h_n) \longmapsto \sum h_i s_i.
$$
Réciproquement, un élément
$q : \mathcal{O}_S^{\oplus n + 1} \to \mathcal{Q}$ de $G(n, n + 1)(S)$
donne le couple $(\mathcal{Q}, (q(e_1), \ldots, q(e_{n + 1})))$.
Ici, les $e_i$, $i = 1, \ldots, n + 1$, sont les sections de la base
canonique du module libre $\mathcal{O}_S^{\oplus n + 1}$.
Nous omettons la vérification que ces constructions définissent des
transformations de foncteurs réciproques.
\end{proof}


\begin{multicols}{2}[\section{Autres chapitres}]
\noindent
Préliminaires
\begin{enumerate}
\item \hyperref[introduction-section-phantom]{Introduction}
\item \hyperref[conventions-section-phantom]{Conventions}
\item \hyperref[sets-section-phantom]{Théorie des ensembles}
\item \hyperref[categories-section-phantom]{Catégories}
\item \hyperref[topology-section-phantom]{Topologie}
\item \hyperref[sheaves-section-phantom]{Faisceaux sur les espaces}
\item \hyperref[sites-section-phantom]{Sites et faisceaux}
\item \hyperref[stacks-section-phantom]{Champs}
\item \hyperref[fields-section-phantom]{Corps}
\item \hyperref[algebra-section-phantom]{Algèbre commutative}
\item \hyperref[brauer-section-phantom]{Groupes de Brauer}
\item \hyperref[homology-section-phantom]{Algèbre homologique}
\item \hyperref[derived-section-phantom]{Catégories dérivées}
\item \hyperref[simplicial-section-phantom]{Méthodes simpliciales}
\item \hyperref[more-algebra-section-phantom]{Compléments d'algèbre}
\item \hyperref[smoothing-section-phantom]{Lissage des morphismes d'anneaux}
\item \hyperref[modules-section-phantom]{Faisceaux de modules}
\item \hyperref[sites-modules-section-phantom]{Modules sur les sites}
\item \hyperref[injectives-section-phantom]{Objets injectifs}
\item \hyperref[cohomology-section-phantom]{Cohomologie des faisceaux}
\item \hyperref[sites-cohomology-section-phantom]{Cohomologie sur les sites}
\item \hyperref[dga-section-phantom]{Algèbre différentielle graduée}
\item \hyperref[dpa-section-phantom]{Algèbre à puissances divisées}
\item \hyperref[sdga-section-phantom]{Faisceaux différentiels gradués}
\item \hyperref[hypercovering-section-phantom]{Hyperrecouvrements}
\end{enumerate}
Schémas
\begin{enumerate}
\setcounter{enumi}{25}
\item \hyperref[schemes-section-phantom]{Schémas}
\item \hyperref[constructions-section-phantom]{Constructions de schémas}
\item \hyperref[properties-section-phantom]{Propriétés des schémas}
\item \hyperref[morphisms-section-phantom]{Morphismes de schémas}
\item \hyperref[coherent-section-phantom]{Cohomologie des schémas}
\item \hyperref[divisors-section-phantom]{Diviseurs}
\item \hyperref[limits-section-phantom]{Limites de schémas}
\item \hyperref[varieties-section-phantom]{Variétés}
\item \hyperref[topologies-section-phantom]{Topologies sur les schémas}
\item \hyperref[descent-section-phantom]{Descente}
\item \hyperref[perfect-section-phantom]{Catégories dérivées de schémas}
\item \hyperref[more-morphisms-section-phantom]{Compléments sur les morphismes}
\item \hyperref[flat-section-phantom]{Compléments sur la platitude}
\item \hyperref[groupoids-section-phantom]{Schémas en groupoïdes}
\item \hyperref[more-groupoids-section-phantom]{Compléments sur les schémas en groupoïdes}
\item \hyperref[etale-section-phantom]{Morphismes étales de schémas}
\end{enumerate}
Thèmes de théorie des schémas
\begin{enumerate}
\setcounter{enumi}{41}
\item \hyperref[chow-section-phantom]{Homologie de Chow}
\item \hyperref[intersection-section-phantom]{Théorie de l'intersection}
\item \hyperref[pic-section-phantom]{Schémas de Picard des courbes}
\item \hyperref[weil-section-phantom]{Théories cohomologiques de Weil}
\item \hyperref[adequate-section-phantom]{Modules adéquats}
\item \hyperref[dualizing-section-phantom]{Complexes dualisants}
\item \hyperref[duality-section-phantom]{Dualité pour les schémas}
\item \hyperref[discriminant-section-phantom]{Discriminants et différentes}
\item \hyperref[derham-section-phantom]{Cohomologie de de Rham}
\item \hyperref[local-cohomology-section-phantom]{Cohomologie locale}
\item \hyperref[algebraization-section-phantom]{Géométrie algébrique et formelle}
\item \hyperref[curves-section-phantom]{Courbes algébriques}
\item \hyperref[resolve-section-phantom]{Résolution des surfaces}
\item \hyperref[models-section-phantom]{Réduction semi-stable}
\item \hyperref[functors-section-phantom]{Foncteurs et morphismes}
\item \hyperref[equiv-section-phantom]{Catégories dérivées de variétés}
\item \hyperref[pione-section-phantom]{Groupes fondamentaux de schémas}
\item \hyperref[etale-cohomology-section-phantom]{Cohomologie étale}
\item \hyperref[crystalline-section-phantom]{Cohomologie cristalline}
\item \hyperref[proetale-section-phantom]{Cohomologie pro-étale}
\item \hyperref[relative-cycles-section-phantom]{Cycles relatifs}
\item \hyperref[more-etale-section-phantom]{Compléments de cohomologie étale}
\item \hyperref[trace-section-phantom]{La formule des traces}
\end{enumerate}
Espaces algébriques
\begin{enumerate}
\setcounter{enumi}{64}
\item \hyperref[spaces-section-phantom]{Espaces algébriques}
\item \hyperref[spaces-properties-section-phantom]{Propriétés des espaces algébriques}
\item \hyperref[spaces-morphisms-section-phantom]{Morphismes d'espaces algébriques}
\item \hyperref[decent-spaces-section-phantom]{Espaces algébriques décents}
\item \hyperref[spaces-cohomology-section-phantom]{Cohomologie des espaces algébriques}
\item \hyperref[spaces-limits-section-phantom]{Limites d'espaces algébriques}
\item \hyperref[spaces-divisors-section-phantom]{Diviseurs sur les espaces algébriques}
\item \hyperref[spaces-over-fields-section-phantom]{Espaces algébriques sur des corps}
\item \hyperref[spaces-topologies-section-phantom]{Topologies sur les espaces algébriques}
\item \hyperref[spaces-descent-section-phantom]{Descente et espaces algébriques}
\item \hyperref[spaces-perfect-section-phantom]{Catégories dérivées d'espaces}
\item \hyperref[spaces-more-morphisms-section-phantom]{Compléments sur les morphismes d'espaces}
\item \hyperref[spaces-flat-section-phantom]{Platitude sur les espaces algébriques}
\item \hyperref[spaces-groupoids-section-phantom]{Groupoïdes dans les espaces algébriques}
\item \hyperref[spaces-more-groupoids-section-phantom]{Compléments sur les groupoïdes dans les espaces}
\item \hyperref[bootstrap-section-phantom]{Amorçage}
\item \hyperref[spaces-pushouts-section-phantom]{Sommes amalgamées d'espaces algébriques}
\end{enumerate}
Thèmes de géométrie
\begin{enumerate}
\setcounter{enumi}{81}
\item \hyperref[spaces-chow-section-phantom]{Groupes de Chow des espaces}
\item \hyperref[groupoids-quotients-section-phantom]{Quotients de groupoïdes}
\item \hyperref[spaces-more-cohomology-section-phantom]{Compléments sur la cohomologie des espaces}
\item \hyperref[spaces-simplicial-section-phantom]{Espaces simpliciaux}
\item \hyperref[spaces-duality-section-phantom]{Dualité pour les espaces}
\item \hyperref[formal-spaces-section-phantom]{Espaces algébriques formels}
\item \hyperref[restricted-section-phantom]{Algébrisation des espaces formels}
\item \hyperref[spaces-resolve-section-phantom]{Résolution des surfaces revisitée}
\end{enumerate}
Théorie des déformations
\begin{enumerate}
\setcounter{enumi}{89}
\item \hyperref[formal-defos-section-phantom]{Théorie formelle des déformations}
\item \hyperref[defos-section-phantom]{Théorie des déformations}
\item \hyperref[cotangent-section-phantom]{Le complexe cotangent}
\item \hyperref[examples-defos-section-phantom]{Problèmes de déformation}
\end{enumerate}
Champs algébriques
\begin{enumerate}
\setcounter{enumi}{93}
\item \hyperref[algebraic-section-phantom]{Champs algébriques}
\item \hyperref[examples-stacks-section-phantom]{Exemples de champs}
\item \hyperref[stacks-sheaves-section-phantom]{Faisceaux sur les champs algébriques}
\item \hyperref[criteria-section-phantom]{Critères de représentabilité}
\item \hyperref[artin-section-phantom]{Axiomes d'Artin}
\item \hyperref[quot-section-phantom]{Espaces de Quot et de Hilbert}
\item \hyperref[stacks-properties-section-phantom]{Propriétés des champs algébriques}
\item \hyperref[stacks-morphisms-section-phantom]{Morphismes de champs algébriques}
\item \hyperref[stacks-limits-section-phantom]{Limites de champs algébriques}
\item \hyperref[stacks-cohomology-section-phantom]{Cohomologie des champs algébriques}
\item \hyperref[stacks-perfect-section-phantom]{Catégories dérivées de champs}
\item \hyperref[stacks-introduction-section-phantom]{Introduction aux champs algébriques}
\item \hyperref[stacks-more-morphisms-section-phantom]{Compléments sur les morphismes de champs}
\item \hyperref[stacks-geometry-section-phantom]{La géométrie des champs}
\end{enumerate}
Thèmes de théorie des espaces de modules
\begin{enumerate}
\setcounter{enumi}{107}
\item \hyperref[moduli-section-phantom]{Champs de modules}
\item \hyperref[moduli-curves-section-phantom]{Espaces de modules de courbes}
\end{enumerate}
Divers
\begin{enumerate}
\setcounter{enumi}{109}
\item \hyperref[examples-section-phantom]{Exemples}
\item \hyperref[exercises-section-phantom]{Exercices}
\item \hyperref[guide-section-phantom]{Guide bibliographique}
\item \hyperref[desirables-section-phantom]{Desiderata}
\item \hyperref[coding-section-phantom]{Style de programmation}
\item \hyperref[obsolete-section-phantom]{Obsolète}
\item \hyperref[fdl-section-phantom]{Licence de documentation libre GNU}
\item \hyperref[index-section-phantom]{Index généré automatiquement}
\end{enumerate}
\end{multicols}


\bibliography{my}
\bibliographystyle{amsalpha}

\end{document}
